% The appendices of the notes: Appendix A is the method, written for the
% students; then one chapter per backed claim, B (CLI efficiency), C (HPC
% and the command line) and D (the Unix lineage).  This is the first deck
% in the course's reading order, so Appendix A describes the method in
% full; the later decks' Appendix A points back here (the backing-claims
% skill, appendix-guide.md, "One method chapter across a collection").
% The prose stands in onlyin blocks, English first, each with its
% Swedish partner, as in contents.nw.

\chapter{\ensv{Checking a claim}{Att kontrollera ett påstående}} % sync english=e8053b swedish=1bf2a3
\label{app:checking}

\begin{onlyin}{english} % sync english=0a8ac2 swedish=3d6d87
These notes, like all teaching material, make a number of claims.  Most
of them you can check at the keyboard: run \texttt{ls | wc -l} and see
whether it prints what the notes say.  Others you can look up in a
source you can open, such as a dictionary or the documentation of a
program, and the notes cite it where they make the claim.  But some
claims are about what is true in the world in general, and a single
source you happen to know cannot settle them.  Three such questions
stand behind \cref{sec:where-cli}:
\begin{enumerate}
  \item whether a command line is more efficient than a graphical
    interface for repeated or bulk tasks, and makes such work easier to
    redo;
  \item whether the large computers of research, computing clusters, are
    reached through a command line;
  \item which of today's systems descend from Unix.
\end{enumerate}
How do you know the answer to such a question?  The same problem meets
you every time you read a textbook, a web page or an answer from a
language model, and every time you write a report yourself.  This
appendix describes a way of checking a claim against the research
literature.  \Cref{app:cli-efficiency,app:hpc-cli,app:unix-lineage}
carry it out for each of the three questions, and each of them can be
read on its own.
\end{onlyin}
\begin{onlyin}{swedish}
Den här texten gör, som allt undervisningsmaterial, ett antal
påståenden.  De flesta kan du kontrollera vid tangentbordet: kör
\texttt{ls | wc -l} och se om det skriver ut det som texten säger.
Andra kan du slå upp i en källa som du kan öppna, som en ordbok eller
dokumentationen till ett program, och texten hänvisar till den där
påståendet görs.  Men en del påståenden gäller vad som är sant i
världen i allmänhet, och en enskild källa som du råkar känna till kan
inte avgöra dem.  Tre sådana frågor står bakom \cref{sec:where-cli}:
\begin{enumerate}
  \item om en kommandorad är effektivare än ett grafiskt gränssnitt för
    upprepade uppgifter och massuppgifter, och gör sådant arbete lättare
    att göra om;
  \item om forskningens stora datorer, beräkningskluster, nås genom en
    kommandorad;
  \item vilka av dagens system som härstammar från Unix.
\end{enumerate}
Hur vet du svaret på en sådan fråga?  Samma problem möter dig varje gång
du läser en lärobok, en webbsida eller ett svar från en språkmodell, och
varje gång du själv skriver en rapport.  Den här bilagan beskriver ett
sätt att kontrollera ett påstående mot forskningslitteraturen.
\Cref{app:cli-efficiency,app:hpc-cli,app:unix-lineage} gör det för var
och en av de tre frågorna, och var och en av dem kan läsas för sig.
\end{onlyin}

\section{\ensv{Turn the claim into a question}%
  {Gör påståendet till en fråga}} % sync english=9f21d5 swedish=aad785

\begin{onlyin}{english} % sync english=1ac72f swedish=d80b7b
A claim can be checked only once it is clear what would show it to be
\emph{wrong}.  So turn it into a question with an answer.  \enquote{Are
computing clusters reached through a command line?} can be answered yes
or no, while \enquote{the command line is important} cannot: no source
could show it wrong, since \enquote{important} is a judgement, not a
fact.  Whoever says it usually means something else, such as \enquote{is
a command line faster for this kind of task?}, and that question can get
the answer no, which is why it can be tested.  Often one claim hides
several.  \enquote{Linux, macOS and Android descend from Unix} is three
questions, and they need different sources.
\end{onlyin}
\begin{onlyin}{swedish}
Ett påstående kan kontrolleras först när det är klart vad som skulle
visa att det är \emph{fel}.  Gör det därför till en fråga med ett svar.
\enquote{Nås beräkningskluster genom en kommandorad?} kan besvaras med
ja eller nej, medan \enquote{kommandoraden är viktig} inte kan det:
ingen källa kunde visa att det är fel, eftersom \enquote{viktig} är ett
omdöme, inte ett faktum.  Den som säger det menar oftast något annat,
som \enquote{är en kommandorad snabbare för den här sortens uppgift?},
och den frågan kan få svaret nej, och därför går den att pröva.  Ofta
döljer ett påstående flera.  \enquote{Linux, macOS och Android
härstammar från Unix} är tre frågor, och de behöver olika källor.
\end{onlyin}

\section{\ensv{Choose the right kind of source}%
  {Välj rätt sorts källa}} % sync english=ed2b21 swedish=a4a560

\begin{onlyin}{english} % sync english=b0ea8e swedish=c84432
Different questions are answered by different sources.
\begin{description}
  \item[Dictionaries and reference works] for what a word \emph{means}.
    The analysis of \emph{interface} in \cref{sec:what-interface} comes
    from a dictionary, and its footnote names the edition and the date
    it was read: dictionaries change.
  \item[The primary source] for where something \emph{comes from}, or
    for what a system or a standard \emph{is}: the paper that first
    described it, the organization that certifies it, the documentation
    of its makers.
  \item[Reviews] for whether something is \emph{established}, or for
    what the research says \emph{taken together}.  A systematic review
    goes through hundreds of studies by a stated method, and so weighs
    more than a single study.
  \item[Single studies] for one specific result, with the reservation
    that a single study shows only that something was observed once.
\end{description}
Where to find them: the university library's search service, and
databases such as Scopus, Web of Science, IEEE Xplore, the ACM Digital
Library, DBLP (computer science) and OpenAlex (open).  Google Scholar
covers a lot but ranks in ways nobody can inspect.  The appendices use
the command-line tool \texttt{scholar}\footnote{The package
\texttt{scholarcli} on PyPI:
\url{https://pypi.org/project/scholarcli/}.}, which asks several
databases the same question at once and records everything that
happened in a \emph{session}.  What you get from the library's search
service is the same, only slower.
\end{onlyin}
\begin{onlyin}{swedish}
Olika frågor besvaras av olika källor.
\begin{description}
  \item[Ordböcker och uppslagsverk] för vad ett ord \emph{betyder}.
    Analysen av \emph{gränssnitt} i \cref{sec:what-interface} kommer
    från en ordbok, och dess fotnot anger utgåvan och dagen den lästes:
    ordböcker ändras.
  \item[Primärkällan] för var något \emph{kommer ifrån}, eller för vad
    ett system eller en standard \emph{är}: artikeln som först beskrev
    det, organisationen som certifierar det, dokumentationen från dem
    som gjort det.
  \item[Översikter] för om något är \emph{etablerat}, eller för vad
    forskningen säger \emph{sammantaget}.  En systematisk översikt går
    igenom hundratals studier enligt en angiven metod, och väger därför
    tyngre än en enskild studie.
  \item[Enskilda studier] för ett bestämt resultat, med förbehållet att
    en enskild studie bara visar att något har observerats en gång.
\end{description}
Var du hittar dem: universitetsbibliotekets söktjänst, och databaser som
Scopus, Web of Science, IEEE Xplore, ACM Digital Library, DBLP
(datavetenskap) och OpenAlex (öppen).  Google Scholar täcker mycket men
rangordnar på sätt som ingen kan granska.  Bilagorna använder
kommandoradsverktyget \texttt{scholar}\footnote{Paketet
\texttt{scholarcli} på PyPI:
\url{https://pypi.org/project/scholarcli/}.}, som ställer samma fråga
till flera databaser på en gång och sparar allt som hände i en
\emph{session}.  Det du får från bibliotekets söktjänst är detsamma,
bara långsammare.
\end{onlyin}

\section{\ensv{Search so that you can be wrong}%
  {Sök så att du kan ha fel}} % sync english=9f82b4 swedish=c2fcb3

\begin{onlyin}{english} % sync english=477fe8 swedish=4d62d6
If you search for the names you already know, you find what you
expected.  So search by \emph{topic}, by concept, not by author.  All
the searches in these appendices were made that way; the sources that
could only be reached as \emph{known items} (web pages, documentation,
books, papers without a DOI) are said to be, with the reason.

A few practical rules:
\begin{itemize}
  \item Keep the queries short.  Several short queries beat one long
    one.
  \item Ask \emph{the same} questions in every database, and search for
    every name of a concept, in both directions.  A counter-search that
    tried only one of the names, in one of the databases, does not show
    that there are no objections.
  \item Search for \emph{objections} too: add words such as
    \enquote{critique}, \enquote{limitations} or \enquote{no
    difference}.  A claim that has survived a counter-search is worth
    more than one that was never tested, and an objection you find
    becomes a \emph{reservation} in the answer, not a defeat.
  \item Expect the databases' ranking to miss things: older journals,
    conference papers without a DOI, web pages.  Then fetch the source
    as a \emph{known item}: you know it exists and look it up directly.
    That is in order, as long as you say so.
\end{itemize}
\end{onlyin}
\begin{onlyin}{swedish}
Söker du efter de namn du redan känner till hittar du det du väntade
dig.  Sök därför efter \emph{ämne}, efter begrepp, inte efter
författare.  Alla sökningar i de här bilagorna gjordes så; de källor som
bara kunde nås som \emph{kända titlar} (webbsidor, dokumentation,
böcker, artiklar utan DOI) sägs vara det, med skälet.

Några praktiska regler:
\begin{itemize}
  \item Håll frågorna korta.  Flera korta frågor slår en lång.
  \item Ställ \emph{samma} frågor i varje databas, och sök på varje namn
    som ett begrepp har, i båda riktningarna.  En motsökning som bara
    prövade ett av namnen, i en av databaserna, visar inte att det inte
    finns några invändningar.
  \item Sök efter \emph{invändningar} också: lägg till ord som
    \enquote{critique}, \enquote{limitations} eller \enquote{no
    difference}.  Ett påstående som har klarat en motsökning är värt
    mer än ett som aldrig prövats, och en invändning som du hittar blir
    ett \emph{förbehåll} i svaret, inte ett nederlag.
  \item Räkna med att databasernas rangordning missar saker: äldre
    tidskrifter, konferensbidrag utan DOI, webbsidor.  Hämta då källan
    som en \emph{känd titel}: du vet att den finns och slår upp den
    direkt.  Det är i sin ordning, så länge du säger det.
\end{itemize}
\end{onlyin}

\section{\ensv{Read in full text, and write down how}%
  {Läs i fulltext, och skriv ner hur}} % sync english=c1a202 swedish=f9f4d6

\begin{onlyin}{english} % sync english=24ea6d swedish=971798
A hit in a database is not evidence.  Nor is the abstract: it says what
the authors want to have said, not what they showed.  Read the whole
text and find \emph{the place} that supports the claim: a verbatim
quotation with its page number.  If you find no such place, the source
does not support the claim, however well its title fits.  A title is
not evidence: a study whose title seems to support a claim may, when you
read it, qualify the claim or have no result on it at all.  Where only
the abstract could be had, say so, and weigh the source accordingly.

Then write down how you went about it, so that someone else can follow
the trail.  In these notes that record stands as a comment above each
entry in the bibliography's source file, with fixed fields:
\begin{description}
  \item[\texttt{CLAIM}] the claim the source is to support;
  \item[\texttt{FOUND-VIA}] how it was found: the query and the
    database, or \enquote{known item};
  \item[\texttt{PICKED}] why this source rather than another hit;
  \item[\texttt{QUOTE}] the verbatim quotation, with its page;
  \item[\texttt{VERIFIED}] that the full text was read, and where;
  \item[\texttt{COUNTER}] the counter-search, and what it found;
  \item[\texttt{DATE}] when.
\end{description}
It looks bureaucratic and takes a couple of minutes per source.  It is
also the only thing that lets you answer, a year later, the question
\enquote{how do you know?}
\end{onlyin}
\begin{onlyin}{swedish}
En träff i en databas är inget belägg.  Det är inte sammanfattningen
heller: den säger vad författarna vill ha sagt, inte vad de visade.  Läs
hela texten och hitta \emph{stället} som stöder påståendet: ett ordagrant
citat med sitt sidnummer.  Hittar du inget sådant ställe stöder källan
inte påståendet, hur väl titeln än passar.  En titel är inget belägg: en
studie vars titel tycks stödja ett påstående kan, när du läser den,
kvalificera påståendet eller sakna resultat om det alls.  Där bara
sammanfattningen gick att få tag på, säg det, och väg källan därefter.

Skriv sedan ner hur du gick till väga, så att någon annan kan följa
spåret.  I den här texten står den uppteckningen som en kommentar ovanför
varje post i källförteckningens källfil, med fasta fält:
\begin{description}
  \item[\texttt{CLAIM}] påståendet som källan ska stödja;
  \item[\texttt{FOUND-VIA}] hur den hittades: frågan och databasen,
    eller \enquote{känd titel};
  \item[\texttt{PICKED}] varför just den här källan och inte en annan
    träff;
  \item[\texttt{QUOTE}] det ordagranna citatet, med sin sida;
  \item[\texttt{VERIFIED}] att fulltexten lästes, och var;
  \item[\texttt{COUNTER}] motsökningen, och vad den gav;
  \item[\texttt{DATE}] när.
\end{description}
Det ser byråkratiskt ut och tar ett par minuter per källa.  Det är också
det enda som låter dig svara, ett år senare, på frågan \enquote{hur vet
du det?}
\end{onlyin}

\section{\ensv{Draw the conclusion, with its reservations}%
  {Dra slutsatsen, med förbehåll}} % sync english=aa6687 swedish=2e6213

\begin{onlyin}{english} % sync english=8a7f23 swedish=9f88e7
The answer is written out together with what speaks against it.  The
reservations are not weaknesses of the answer.  They \emph{are} the
answer, and they decide how the claim may be used in the notes: a claim
is stated in the main text at the strength the appendix supports, never
stronger.  Each appendix that reports a search therefore ends with a
section \emph{Further work}: what that search left undone, and what
would have to be studied for a surer answer.
\end{onlyin}
\begin{onlyin}{swedish}
Svaret skrivs ut tillsammans med det som talar emot det.  Förbehållen är
inte svagheter i svaret.  De \emph{är} svaret, och de avgör hur
påståendet får användas i texten: ett påstående står i huvudtexten med
den styrka som bilagan bär, aldrig starkare.  Varje bilaga som redovisar
en sökning slutar därför med ett avsnitt \emph{Fortsatt arbete}: vad
den sökningen lämnade ogjort, och vad som skulle behöva studeras för ett
säkrare svar.
\end{onlyin}

% TODO(claims): a table of the three questions and their answers,
% tab:three-claims, once appendices B and C are final.

\section{\ensv{A note on language models}%
  {En anmärkning om språkmodeller}} % sync english=9c8834 swedish=eaaa32

\begin{onlyin}{english} % sync english=032880 swedish=acb44b
The searches in these appendices gave hundreds of hits each.  A language
model was used to \emph{sort} them (does the hit belong to the topic at
all, and which way does it point?) so that the screening could be
checked.  Notice the division of labour: the model sorted, but every
source that is cited was read and chosen by a person, and the sorting
too is to be confirmed by a person.  The model's classification
therefore stands in the tables with its confidence, so that it can be
questioned, and the start of each appendix says how far the human
review has come.  Do use language models to find and sort; never let
them be the last link.  They can be wrong, and it is you who must answer
for the content.
\end{onlyin}
\begin{onlyin}{swedish}
Sökningarna i de här bilagorna gav hundratals träffar var.  En
språkmodell användes för att \emph{sortera} dem (hör träffen alls till
ämnet, och åt vilket håll pekar den?) så att gallringen kunde
kontrolleras.  Lägg märke till arbetsfördelningen: modellen sorterade,
men varje källa som citeras lästes och valdes av en människa, och även
sorteringen ska bekräftas av en människa.  Modellens klassificering står
därför i tabellerna med sin konfidens, så att den kan ifrågasättas, och
början av varje bilaga säger hur långt den mänskliga granskningen har
kommit.  Använd gärna språkmodeller för att hitta och sortera; låt dem
aldrig vara den sista länken.  De kan ha fel, och det är du som måste
stå för innehållet.
\end{onlyin}

% --- Appendix B: claim B, from its search of 2026-11-18 (scholar session
% cli-efficiency).  The main text of sec:where-cli keeps within the bounds
% its conclusion sets.
% The full question is the table of contents' entry; the running head,
% which has room for one line, gets the short form.
\chapter[\ensv{Is a command line more efficient or reproducible than a
  graphical interface for repeated or bulk tasks?}%
  {Är en kommandorad effektivare eller lättare att upprepa än ett
  grafiskt gränssnitt för upprepade uppgifter och massuppgifter?}]% % sync english=db83f6 swedish=2061b9
  [\ensv{Command line or graphical interface?}%
  {Kommandorad eller grafiskt gränssnitt?}]% % sync english=b0df85 swedish=b14a13
  {\ensv{Is a command line more efficient or reproducible than a
  graphical interface for repeated or bulk tasks?}%
  {Är en kommandorad effektivare eller lättare att upprepa än ett
  grafiskt gränssnitt för upprepade uppgifter och massuppgifter?}} % sync english=db83f6 swedish=2061b9
\label{app:cli-efficiency}

\chapterprecis{\ensv{The author has not yet reviewed the results in this
  appendix in full.}%
  {Författaren har ännu inte granskat resultaten i den här bilagan i sin
  helhet.}} % sync english=21ea16 swedish=60ed7a

\begin{onlyin}{english} % sync english=589f3a swedish=819ab2
In \cref{sec:cli-both} we let you try a bulk task in your head: put the
contents of 368 daily reports, picked out of 3287 by their dates, into
one file, first in a file manager and then at the command line.  The
point of the section is that at a command line a program can take your
place, so that work done once can be written down as commands and run
again, on any number of files.  In \cref{sec:where-cli} we add that in
research you will want to process the files your experiments produce
with commands rather than by hand.  Behind both sections lies a claim:
for repeated or bulk tasks a
command-line interface, where commands can be combined and saved as a
script, is more efficient than a graphical interface and makes the work
easier to reproduce; but command lines are harder for novices to learn.

That the line of commands can be saved and run again is not in doubt:
it is a fact about the shell, and the lecture \emph{Scripts} shows it
by running it.  What we test here are the empirical parts.  The
question of this chapter is: \emph{is a command line more efficient or
reproducible than a graphical interface for repeated or bulk tasks?}
We split it into three sub-questions.
\begin{description}
  \item[Efficiency] Do people work faster or with fewer errors with a
    command interface than with a graphical, menu or
    direct-manipulation interface, and for which tasks and users?
  \item[Reproducibility] Does scripted, command-line work make
    research easier to reproduce than manual steps in a graphical
    program, and is there evidence for this beyond advice?
  \item[Novices] Are command lines harder for novices to learn than
    graphical interfaces?
\end{description}
The strength we test is that of the deck's sections: not that a
command line is always better, but that it pays off for repeated or
bulk work and costs more at the start.  We describe the search in
\cref{sec:cli-efficiency-method}, report what it found for each
sub-question in \cref{sec:cli-efficiency-results}, discuss its limits
in \cref{sec:cli-efficiency-discussion}, answer the question in
\cref{sec:cli-efficiency-conclusion} and list what remains to be done
in \cref{sec:cli-efficiency-further}.  The classified hits are listed
at the end, in \cref{sec:cli-efficiency-hits}.
\end{onlyin}
\begin{onlyin}{swedish}
I \cref{sec:cli-both} lät vi dig pröva en massuppgift i huvudet: att
lägga innehållet i 368 dagliga rapporter, utvalda bland 3287 efter sina
datum, i en fil, först i en filhanterare och sedan på kommandoraden.
Poängen med avsnittet är att ett program kan ta din plats på en
kommandorad, så att arbete som gjorts en gång kan skrivas ner som
kommandon och köras igen, på hur många filer som helst.  I
\cref{sec:where-cli} lägger vi till att du i forskningen kommer att
vilja bearbeta de filer som dina experiment ger med kommandon snarare än
för hand.  Bakom båda avsnitten ligger ett påstående: för upprepade
uppgifter och massuppgifter är ett kommandoradsgränssnitt, där kommandon
kan kombineras och sparas som ett skript, effektivare än ett grafiskt
gränssnitt och gör arbetet lättare att upprepa; men kommandorader är
svårare för nybörjare att lära sig.

Att raden med kommandon kan sparas och köras igen är inte i tvivel: det
är ett faktum om skalet, och föreläsningen \emph{Skript} visar det genom
att köra den.  Det vi prövar här är de empiriska delarna.  Kapitlets
fråga är: \emph{är en kommandorad effektivare eller lättare att upprepa
än ett grafiskt gränssnitt för upprepade uppgifter och massuppgifter?}
Vi delar upp den i tre delfrågor.
\begin{description}
  \item[Effektivitet] Arbetar människor snabbare eller med färre fel
    med ett kommandogränssnitt än med ett grafiskt gränssnitt, ett
    menygränssnitt eller ett gränssnitt för direkt manipulation, och för
    vilka uppgifter och användare?
  \item[Upprepbarhet] Gör skriptat arbete på kommandoraden forskning
    lättare att upprepa än manuella steg i ett grafiskt program, och
    finns det belägg för det utöver råd?
  \item[Nybörjare] Är kommandorader svårare för nybörjare att lära sig
    än grafiska gränssnitt?
\end{description}
Styrkan vi prövar är föreläsningsavsnittens: inte att en kommandorad
alltid är bättre, utan att den lönar sig för upprepat arbete och
massarbete och kostar mer i början.  Vi beskriver sökningen i
\cref{sec:cli-efficiency-method}, redovisar vad den fann för varje
delfråga i \cref{sec:cli-efficiency-results}, diskuterar dess gränser i
\cref{sec:cli-efficiency-discussion}, besvarar frågan i
\cref{sec:cli-efficiency-conclusion} och räknar upp det som återstår i
\cref{sec:cli-efficiency-further}.  De klassade träffarna listas i
slutet, i \cref{sec:cli-efficiency-hits}.
\end{onlyin}

\section{\ensv{Method}{Metod}} % sync english=52a0f9 swedish=f41b92
\label{sec:cli-efficiency-method}

\begin{onlyin}{english} % sync english=153e41 swedish=3c05f3
We searched on 18 November 2026 with the tool \texttt{scholar}, which
sends one query to several bibliographic databases and records every
query, its results and how each database answered.  Before the search
we checked the databases.  Web of Science, Scopus, IEEE Xplore and
OpenAlex answered.  DBLP answered every request with a page that stops
automated clients instead of with data, so DBLP could not be searched;
its column in \cref{tab:cli-efficiency-queries} marks every query as
failed.  Semantic Scholar first rejected our access key; with a new
key later the same day, 18 November 2026, its search answered, but it
takes no Boolean operators: query S1 returned nothing, which is not a
real negative, and we did not use it for the other queries.  We used
it, and Crossref, to look up missing abstracts.
\end{onlyin}
\begin{onlyin}{swedish}
Vi sökte den 18 november 2026 med verktyget \texttt{scholar}, som
skickar en fråga till flera bibliografiska databaser och sparar varje
fråga, dess resultat och hur varje databas svarade.  Före sökningen
kontrollerade vi databaserna.  Web of Science, Scopus, IEEE Xplore och
OpenAlex svarade.  DBLP svarade på varje förfrågan med en sida som
stoppar automatiska klienter i stället för med data, så DBLP kunde inte
genomsökas; dess kolumn i \cref{tab:cli-efficiency-queries} markerar
varje fråga som misslyckad.  Semantic Scholar avvisade först vår
åtkomstnyckel; med en ny nyckel senare samma dag, den 18 november 2026,
svarade dess sökning, men den tar inga booleska operatorer: frågan S1
gav ingenting, vilket inte är ett verkligt negativt svar, och vi använde
den inte för de andra frågorna.  Vi använde den, och Crossref, för att
slå upp saknade sammanfattningar.
\end{onlyin}

\paragraph{\ensv{Concepts and their wordings.}%
  {Begreppen och deras formuleringar.}} % sync english=ad9fc5 swedish=437460
\begin{onlyin}{english} % sync english=29e2b4 swedish=bf4f18
The claim rests on a small number of concepts, and we gave each of
them several wordings, since a study that says \enquote{command
language} is not found by a query for \enquote{command line}.  Words in
quotation marks were searched as phrases; the blocks below were joined
with \texttt{AND} and their words with \texttt{OR}.
% Ragged right: the items are runs of long quoted query phrases.
\begingroup\raggedright
\begin{description}
  \item[CLI] \enquote{command line}, \enquote{command-line},
    \enquote{command language(s)}, \enquote{command-based},
    \enquote{command-driven}, \enquote{command interface(s)},
    \enquote{typed commands}, shell, Unix.
  \item[CMD] command, commands: the bare words, used only beside the
    narrow block DM, since alone they also mean a military command or a
    voice command.
  \item[GUI] \enquote{graphical user interface}, \enquote{graphical
    interface}, GUI(s), \enquote{direct manipulation}, menu(s),
    \enquote{menu-driven}, \enquote{point-and-click}, icon(s), iconic.
  \item[DM] \enquote{direct manipulation}, iconic, \enquote{menu
    selection}, \enquote{menu-driven}, \enquote{menu-based},
    \enquote{menu interface(s)}, \enquote{graphical user
    interface(s)}.
  \item[EFF] \enquote{task performance}, \enquote{user performance},
    \enquote{completion time}, \enquote{task time}, speed, efficiency,
    errors, \enquote{error rate}, productivity.
  \item[EXP] expert(s), expertise, \enquote{experienced users},
    \enquote{frequent users}, \enquote{power users}, repetitive,
    \enquote{repeated tasks}, \enquote{batch processing}, bulk.
  \item[SCR] script(s), scripted, scripting.
  \item[REP] reproducible, reproducibility, replicable, replicability,
    \enquote{reproducible research}; in C7 also replication.
  \item[MAN] \enquote{point-and-click}, \enquote{graphical user
    interface}, \enquote{graphical interface}, GUI, \enquote{by hand},
    \enquote{manual steps}, \enquote{manual data}, \enquote{manual
    processing}.
  \item[ADV] \enquote{best practices}, \enquote{good practices},
    \enquote{good enough practices}, \enquote{simple rules},
    \enquote{ten simple rules}.
  \item[RES] \enquote{scientific computing}, \enquote{research
    computing}, \enquote{computational research}, \enquote{computational
    science}, \enquote{research software}; in S6 also \enquote{data
    analysis} and bioinformatics.
  \item[AUT] \enquote{command line}, \enquote{command-line}, shell,
    script(s), scripted, scripting, automate(d), automation, automating.
  \item[NOV] novice(s), beginner(s), \enquote{new users},
    \enquote{casual users}, \enquote{occasional users},
    \enquote{inexperienced users}, learnability, \enquote{ease of
    learning}, \enquote{ease of use}.
  \item[DIF] difficult, difficulty, difficulties, errors, mistakes,
    misconception(s), barrier(s), \enquote{hard to learn}.
  \item[STYLE] the field's own names for the concept:
    \enquote{interaction style(s)}, \enquote{dialogue style},
    \enquote{dialog style}, \enquote{interface style(s)}; in C2 and C3
    also WIMP beside GUI and DM.
  \item[NEG] \enquote{no significant difference}, \enquote{no
    difference}, superior, outperform(ed), preferred, preference,
    satisfaction, \enquote{easier to use}, \enquote{easier to learn}.
  \item[FAIL] \enquote{not reproducible}, irreproducible,
    \enquote{re-execution}, \enquote{re-execute}, \enquote{failed to
    run}, \enquote{fail to run}, \enquote{execution errors},
    \enquote{run successfully}, \enquote{executed successfully}; in C7
    also \enquote{re-executed}, \enquote{without error(s)}.
  \item[GREP] \enquote{graphical user interface}, \enquote{graphical
    interface}, GUI, \enquote{point-and-click}, \enquote{workflow
    system(s)}, \enquote{workflow management}.
  \item[PROV] provenance, \enquote{without programming}, \enquote{no
    programming}, \enquote{non-programmers}, \enquote{without coding}.
\end{description}
\par\endgroup
The three title queries T1--T3 use wider wordings that would drown an
abstract search: in T1, command(s), \enquote{command line},
\enquote{command-line}, \enquote{command-based}, \enquote{command
language}, \enquote{character-based} and \enquote{text-based} beside
menu(s), \enquote{direct manipulation}, graphical, graphic, GUI,
iconic, \enquote{point-and-click} and WIMP; in T2, the CLI block with
\enquote{Unix shell}, \enquote{shell commands}, \enquote{shell
scripting} and bash in place of the bare word shell, beside novice(s),
beginner(s), learning, learn, learners, teaching, teach, students,
errors, mistakes, difficulties and literacy; in T3, \enquote{command
name(s)}, \enquote{command syntax}, CUI and \enquote{character user
interface} beside the words of T2's second block, of T1's second block
or desktop.
\end{onlyin}
\begin{onlyin}{swedish}
Påståendet vilar på ett litet antal begrepp, och vi gav vart och ett av
dem flera formuleringar, eftersom en studie som säger \enquote{command
language} inte hittas av en fråga efter \enquote{command line}.  Ord
inom citattecken söktes som fraser; blocken nedan förenades med
\texttt{AND} och deras ord med \texttt{OR}.
% Ragged right: the items are runs of long quoted query phrases.
\begingroup\raggedright
\begin{description}
  \item[CLI] \enquote{command line}, \enquote{command-line},
    \enquote{command language(s)}, \enquote{command-based},
    \enquote{command-driven}, \enquote{command interface(s)},
    \enquote{typed commands}, shell, Unix.
  \item[CMD] command, commands: de nakna orden, använda bara bredvid det
    snäva blocket DM, eftersom de ensamma också betyder en militär order
    eller ett röstkommando.
  \item[GUI] \enquote{graphical user interface}, \enquote{graphical
    interface}, GUI(s), \enquote{direct manipulation}, menu(s),
    \enquote{menu-driven}, \enquote{point-and-click}, icon(s), iconic.
  \item[DM] \enquote{direct manipulation}, iconic, \enquote{menu
    selection}, \enquote{menu-driven}, \enquote{menu-based},
    \enquote{menu interface(s)}, \enquote{graphical user
    interface(s)}.
  \item[EFF] \enquote{task performance}, \enquote{user performance},
    \enquote{completion time}, \enquote{task time}, speed, efficiency,
    errors, \enquote{error rate}, productivity.
  \item[EXP] expert(s), expertise, \enquote{experienced users},
    \enquote{frequent users}, \enquote{power users}, repetitive,
    \enquote{repeated tasks}, \enquote{batch processing}, bulk.
  \item[SCR] script(s), scripted, scripting.
  \item[REP] reproducible, reproducibility, replicable, replicability,
    \enquote{reproducible research}; i C7 också replication.
  \item[MAN] \enquote{point-and-click}, \enquote{graphical user
    interface}, \enquote{graphical interface}, GUI, \enquote{by hand},
    \enquote{manual steps}, \enquote{manual data}, \enquote{manual
    processing}.
  \item[ADV] \enquote{best practices}, \enquote{good practices},
    \enquote{good enough practices}, \enquote{simple rules},
    \enquote{ten simple rules}.
  \item[RES] \enquote{scientific computing}, \enquote{research
    computing}, \enquote{computational research}, \enquote{computational
    science}, \enquote{research software}; i S6 också \enquote{data
    analysis} och bioinformatics.
  \item[AUT] \enquote{command line}, \enquote{command-line}, shell,
    script(s), scripted, scripting, automate(d), automation, automating.
  \item[NOV] novice(s), beginner(s), \enquote{new users},
    \enquote{casual users}, \enquote{occasional users},
    \enquote{inexperienced users}, learnability, \enquote{ease of
    learning}, \enquote{ease of use}.
  \item[DIF] difficult, difficulty, difficulties, errors, mistakes,
    misconception(s), barrier(s), \enquote{hard to learn}.
  \item[STYLE] fältets egna namn på begreppet:
    \enquote{interaction style(s)}, \enquote{dialogue style},
    \enquote{dialog style}, \enquote{interface style(s)}; i C2 och C3
    också WIMP bredvid GUI och DM.
  \item[NEG] \enquote{no significant difference}, \enquote{no
    difference}, superior, outperform(ed), preferred, preference,
    satisfaction, \enquote{easier to use}, \enquote{easier to learn}.
  \item[FAIL] \enquote{not reproducible}, irreproducible,
    \enquote{re-execution}, \enquote{re-execute}, \enquote{failed to
    run}, \enquote{fail to run}, \enquote{execution errors},
    \enquote{run successfully}, \enquote{executed successfully}; i C7
    också \enquote{re-executed}, \enquote{without error(s)}.
  \item[GREP] \enquote{graphical user interface}, \enquote{graphical
    interface}, GUI, \enquote{point-and-click}, \enquote{workflow
    system(s)}, \enquote{workflow management}.
  \item[PROV] provenance, \enquote{without programming}, \enquote{no
    programming}, \enquote{non-programmers}, \enquote{without coding}.
\end{description}
\par\endgroup
De tre titelfrågorna T1--T3 använder vidare formuleringar som skulle
dränka en sökning i sammanfattningar: i T1 command(s), \enquote{command
line}, \enquote{command-line}, \enquote{command-based}, \enquote{command
language}, \enquote{character-based} och \enquote{text-based} bredvid
menu(s), \enquote{direct manipulation}, graphical, graphic, GUI, iconic,
\enquote{point-and-click} och WIMP; i T2 blocket CLI med \enquote{Unix
shell}, \enquote{shell commands}, \enquote{shell scripting} och bash i
stället för det nakna ordet shell, bredvid novice(s), beginner(s),
learning, learn, learners, teaching, teach, students, errors, mistakes,
difficulties och literacy; i T3 \enquote{command name(s)},
\enquote{command syntax}, CUI och \enquote{character user interface}
bredvid orden i T2:s andra block, i T1:s andra block eller desktop.
\end{onlyin}

% Author's note (verification and reproduction): see the block above the
% hit-list \input at the end of the chapter.
\begin{table}[htbp]
  \begin{sidecaption}{\ensv{Queries and hits per database, 18 November
    2026, with the tool \texttt{scholar}.  S = support query, C = counter
    query, T = query in titles only; no source was retrieved as a known
    item.  The blocks are defined in the text.  Databases: OpenAlex
    (OA; not open access), Web of Science (WoS), IEEE Xplore (IEEE),
    Scopus, DBLP.  The searches have no upper limit: the numbers are
    all hits.  OpenAlex was searched in title and abstract (S, C) or in
    titles (T): its full-text search returned between 31\,405 and
    2\,478\,176 records per query, too many to screen.  Web of Science
    got the keywords as \texttt{TS=(\dots)} (T: \texttt{TI=(\dots)}),
    Scopus as \texttt{TITLE-ABS-KEY(\dots)} (T: \texttt{TITLE(\dots)}),
    IEEE as Boolean free text (T: every word as a \enquote{Document
    Title} field).  \enquote{--} means that the search failed: DBLP
    answered every query with a page that stops automated clients, so
    DBLP could not be searched.  OpenAlex answered S3 at the third
    attempt and S6 at the second (rate limit), IEEE answered C7 at the
    second (timeout).  Semantic Scholar, not in the table, answered S1
    with 0 hits: its search takes no Boolean operators, so it was not
    used further.}%
    {Frågor och träffar per databas, 18 november 2026, med verktyget
    \texttt{scholar}.  S = stödfråga, C = motfråga, T = fråga bara i
    titlar; ingen källa hämtades som känd titel.  Blocken definieras i
    texten.  Databaser: OpenAlex (OA; inte open access), Web of Science
    (WoS), IEEE Xplore (IEEE), Scopus, DBLP.  Sökningarna har ingen övre
    gräns: talen är alla träffar.  OpenAlex söktes i titel och
    sammanfattning (S, C) eller i titlar (T): dess fulltextsökning gav
    mellan 31\,405 och 2\,478\,176 poster per fråga, för många att
    gallra.  Web of Science fick nyckelorden som \texttt{TS=(\dots)} (T:
    \texttt{TI=(\dots)}), Scopus som \texttt{TITLE-ABS-KEY(\dots)} (T:
    \texttt{TITLE(\dots)}), IEEE som boolesk fritext (T: varje ord som
    ett fält \enquote{Document Title}).  \enquote{--} betyder att
    sökningen misslyckades: DBLP svarade på varje fråga med en sida som
    stoppar automatiska klienter, så DBLP kunde inte genomsökas.
    OpenAlex svarade på S3 vid tredje försöket och på S6 vid det andra
    (hastighetsgräns), IEEE svarade på C7 vid det andra (tidsgräns).
    Semantic Scholar, som inte står i tabellen, svarade på S1 med 0
    träffar: dess sökning tar inga booleska operatorer, så den användes
    inte vidare.}% % sync english=ba6a77 swedish=282400
  }[tab:cli-efficiency-queries]
    \centering\footnotesize
    \begin{tabular}{@{}llrrrrr@{}}
      \toprule
      & \ensv{Blocks}{Block} & OA & WoS & IEEE & Scopus & DBLP \\ % sync english=1cd5a6 swedish=211d0b
      \midrule
      S1 & CLI, GUI, EFF & 869 & 209 & 214 & 407 & -- \\
      S2 & CMD, DM, EFF & 820 & 205 & 472 & 407 & -- \\
      S3 & CLI, GUI, EXP & 639 & 198 & 133 & 351 & -- \\
      S4 & CMD, DM, EXP & 476 & 121 & 131 & 207 & -- \\
      S5 & \ensv{CLI or SCR, REP, MAN}{CLI eller SCR, REP, MAN} % sync english=38a04e swedish=be1a35
         & 893 & 157 & 25 & 212 & -- \\
      S6 & ADV, RES, AUT & 611 & 104 & 222 & 278 & -- \\
      S7 & spreadsheet(s), REP, error(s) & 348 & 37 & 7 & 68 & -- \\
      S8 & ADV, RES & 554 & 133 & 535 & 234 & -- \\
      S9 & \ensv{spreadsheet(s) or Excel, errors}%
           {spreadsheet(s) eller Excel, errors} % sync english=0270a4 swedish=d31b9e
         & 138 & 22 & 1 & 56 & -- \\
      C1 & CLI, NOV, DIF & 1035 & 84 & 64 & 143 & -- \\
      C2 & \ensv{CLI or STYLE, GUI, NOV or NEG}%
           {CLI eller STYLE, GUI, NOV eller NEG} % sync english=d807ca swedish=ce2ec0
         & 852 & 162 & 119 & 413 & -- \\
      C3 & \ensv{CMD or STYLE, DM, NOV or NEG}%
           {CMD eller STYLE, DM, NOV eller NEG} % sync english=b0f883 swedish=aa1ea1
         & 858 & 150 & 207 & 376 & -- \\
      C4 & \ensv{CLI or SCR, REP, FAIL}{CLI eller SCR, REP, FAIL} % sync english=5ea64c swedish=103da9
         & 251 & 18 & 5 & 27 & -- \\
      C5 & GREP, REP, PROV & 293 & 105 & 59 & 134 & -- \\
      C6 & ADV, RES, AUT, criticism & 197 & 28 & 54 & 53 & -- \\
      C7 & \ensv{CLI, SCR or code, REP, FAIL}%
           {CLI, SCR eller code, REP, FAIL} % sync english=4ebab5 swedish=611218
         & 811 & 114 & 57 & 145 & -- \\
      T1 & \ensv{command words, graphical words}%
           {kommandoord, grafiska ord} % sync english=5329ca swedish=644d44
         & 260 & 109 & 35 & 161 & -- \\
      T2 & \ensv{CLI words, learning words}{CLI-ord, inlärningsord} % sync english=1ad9ac swedish=fede24
         & 221 & 64 & 12 & 81 & -- \\
      T3 & \ensv{command names, CUI}{kommandonamn, CUI} % sync english=69763d swedish=ef3fdb
         & 65 & 15 & 4 & 37 & -- \\
      \bottomrule
    \end{tabular}
  \end{sidecaption}
\end{table}

\paragraph{\ensv{Queries.}{Frågor.}} % sync english=060671 swedish=26c3ae
\begin{onlyin}{english} % sync english=b27c31 swedish=8630af
\Cref{tab:cli-efficiency-queries} lists the queries.  S1--S4 search the
efficiency side, comparisons of command and graphical interfaces in
general (S1, S2) and for experts and repetitive work (S3, S4); S1 and S3
use the phrases for a command line, S2 and S4 the bare word command
beside the narrower block DM, as older studies say \enquote{command
interface} or \enquote{commands vs.\ direct manipulation}.  S5--S9
search the reproducibility side: scripted against manual or graphical
work (S5), best-practice advice for research computing (S6, S8) and
errors in spreadsheets (S7, S9).  C1--C7 are the counter queries.  They
reuse every concept block of the support queries and add the field's
own names for interaction styles: novices' difficulties with command
lines (C1), studies where the graphical side won, where no difference
was found, or where novices fared differently (C2, C3), scripts that
fail to run again (C4, C7), graphical tools that make work
reproducible (C5) and criticism of the best-practice advice (C6).  The
title queries T1--T3 are described below, under deviations.
\end{onlyin}
\begin{onlyin}{swedish}
\Cref{tab:cli-efficiency-queries} listar frågorna.  S1--S4 söker på
effektivitetssidan, jämförelser av kommandogränssnitt och grafiska
gränssnitt i allmänhet (S1, S2) och för experter och upprepat arbete
(S3, S4); S1 och S3 använder fraserna för en kommandorad, S2 och S4 det
nakna ordet command bredvid det snävare blocket DM, eftersom äldre
studier säger \enquote{command interface} eller \enquote{commands vs.\
direct manipulation}.  S5--S9 söker på upprepbarhetssidan: skriptat mot
manuellt eller grafiskt arbete (S5), råd om bästa praxis för
datoranvändning i forskning (S6, S8) och fel i kalkylblad (S7, S9).
C1--C7 är motfrågorna.  De återanvänder varje begreppsblock i
stödfrågorna och lägger till fältets egna namn på interaktionsstilar:
nybörjares svårigheter med kommandorader (C1), studier där den grafiska
sidan vann, där ingen skillnad hittades eller där nybörjare klarade sig
annorlunda (C2, C3), skript som inte går att köra igen (C4, C7),
grafiska verktyg som gör arbete upprepbart (C5) och kritik av råden om
bästa praxis (C6).  Titelfrågorna T1--T3 beskrivs nedan, under
avvikelser.
\end{onlyin}

\paragraph{\ensv{Studies the queries had to return.}%
  {Studier som frågorna måste ge.}} % sync english=50f120 swedish=6aaaf2
\begin{onlyin}{english} % sync english=504d13 swedish=ed63b4
A first run of S1--S4 and C1--C3 returned, among others, the
comparison of command, menu and iconic interfaces by
\textcite{Whiteside1985User}, the file-manipulation experiment of
\textcite{Margono1995Study} and the comparison of command, menu and
natural-language interfaces by \textcite{Hauptmann1983Comparison}; a
first run of S5--S6 returned the rules of \textcite{Sandve2013Ten} and
the best practices of \textcite{Wilson2014Best}.  To these we added
three studies we expected on the reproducibility side before relying on
the search: the follow-up \enquote{good enough} practices
\autocite{Wilson2017Good}, the audit of spreadsheet errors in genomics
\autocite{Ziemann2016Gene} and a re-execution study of research code
\autocite{Trisovic2022Large}.  The first run did not return these
three; the repaired queries described next do, and all eight are in the
final record.
\end{onlyin}
\begin{onlyin}{swedish}
En första körning av S1--S4 och C1--C3 gav bland annat jämförelsen av
kommando-, meny- och ikongränssnitt av \textcite{Whiteside1985User},
experimentet med filhantering av \textcite{Margono1995Study} och
jämförelsen av kommando-, meny- och naturligt språk-gränssnitt av
\textcite{Hauptmann1983Comparison}; en första körning av S5--S6 gav
reglerna av \textcite{Sandve2013Ten} och de bästa metoderna av
\textcite{Wilson2014Best}.  Till dessa lade vi tre studier som vi
väntade oss på upprepbarhetssidan innan vi litade på sökningen: den
uppföljande \enquote{good enough}-praxisen\autocite{Wilson2017Good},
granskningen av kalkylbladsfel i genomik\autocite{Ziemann2016Gene} och
en studie av omkörning av forskningskod\autocite{Trisovic2022Large}.
Den första körningen gav inte de tre; de reparerade frågor som beskrivs
härnäst gör det, och alla åtta finns i den slutliga uppteckningen.
\end{onlyin}

\paragraph{\ensv{Check by meaning and deviations.}%
  {Kontroll efter betydelse och avvikelser.}} % sync english=db4b61 swedish=2baf6e
\begin{onlyin}{english} % sync english=0f2980 swedish=1b5782
After the keyword searches we asked OpenAlex's search by meaning one
question per concept, seven in all, and looked up each of the 50 works
it returned per question in our record.  The check found gaps: older
interface comparisons whose abstracts lack our outcome words, or whose
record has no abstract at all, and studies that name the command side
\enquote{character-based}, \enquote{text-based} or \enquote{CUI}.  We
repaired them by concept, not by title: S6 was joined by S8 without
the automation words, S7 by S9, C4 by C7 with the words code,
replication and \enquote{without error(s)}, and we added the title
queries T1--T3, which can use wide wordings because titles are short.
S9's phrases \enquote{supplementary files} and \enquote{scientific
literature} come from the abstract of \textcite{Ziemann2016Gene} and
are in that sense fitted to it.  A second check by meaning, of the
three concepts on interfaces and novices, found the gaps closed except
for studies outside the claim (input devices, menu design, Unix
textbooks) and two on it that have no abstract in OpenAlex and that no
keyword query returned: an experiment on icons against text and direct
manipulation against menus, and a comparison of a textual and a
graphic interface for novice and expert searchers.  We could not
obtain either.
\end{onlyin}
\begin{onlyin}{swedish}
Efter nyckelordssökningarna ställde vi en fråga per begrepp till
OpenAlex sökning efter betydelse, sju i allt, och slog upp vart och ett
av de 50 verk som den gav per fråga i vår uppteckning.  Kontrollen fann
luckor: äldre jämförelser av gränssnitt vars sammanfattningar saknar
våra utfallsord, eller vars post inte har någon sammanfattning alls, och
studier som kallar kommandosidan \enquote{character-based},
\enquote{text-based} eller \enquote{CUI}.  Vi reparerade dem efter
begrepp, inte efter titel: S6 fick sällskap av S8 utan
automatiseringsorden, S7 av S9, C4 av C7 med orden code, replication och
\enquote{without error(s)}, och vi lade till titelfrågorna T1--T3, som
kan använda vida formuleringar eftersom titlar är korta.  S9:s fraser
\enquote{supplementary files} och \enquote{scientific literature} kommer
från sammanfattningen av \textcite{Ziemann2016Gene} och är i den
meningen anpassade efter den.  En andra kontroll efter betydelse, av de
tre begreppen om gränssnitt och nybörjare, fann luckorna stängda utom
för studier utanför påståendet (inmatningsenheter, menydesign, läroböcker
om Unix) och två om det som saknar sammanfattning i OpenAlex och som
ingen nyckelordsfråga gav: ett experiment om ikoner mot text och direkt
manipulation mot menyer, och en jämförelse av ett textbaserat och ett
grafiskt gränssnitt för nybörjare och experter som söker.  Vi kunde inte
få tag på någon av dem.
\end{onlyin}

\paragraph{\ensv{Screening.}{Gallring.}} % sync english=156955 swedish=fe7361
\begin{onlyin}{english} % sync english=d49134 swedish=2d4618
The searches returned 10\,405 records, 9\,744 after records of the same
work from different databases were merged.  A record is judged on its
abstract, so a record without one is excluded unless one can be found.
After the searches 1\,955 records had no abstract.  Three of them were
cited sources whose duplicate record held the abstract, which we
copied, and five were duplicates of cited sources.  For the other
1\,947 we looked up the abstract by DOI, in Semantic Scholar and then
in Crossref: 464 gained one (393 from Semantic Scholar, 71 from
Crossref), and the remaining 1\,483, of which 557 had no DOI, were
excluded as \emph{no abstract}.  All other records were classified
against a written statement of the research question by a language
model (Claude Sonnet~5.5) from title and abstract, as supporting the
claim, qualifying it, an adjacent subject or off topic, and every
excluded record was given one of the labelled criteria below.  We then
read the titles of all records the model kept as supporting or
qualifying and of those it excluded with a confidence below 0.7, and
the abstracts of those that might bear on the claim.  Then we read the
abstract of every record still labelled supporting or qualifying and
not yet cited, 154 works: 69 were confirmed and are cited below, 77
were relabelled and excluded, 6 were further records of cited works,
and two whose authors could not be verified are not cited.  In all we
recorded 272 decisions of our own besides the \emph{no abstract} rule:
the 121 cited sources were kept with the theme under which they are
cited, and 151 records were excluded, among them 26 duplicates of
cited sources.  Of the 10\,405 records, 210 were kept (66 supporting
the claim, 66 qualifying it, 78 on an adjacent subject) and 10\,195
excluded: 2\,670 as \emph{not interfaces}, 2\,668 as \emph{tool
only}, 1\,483 as \emph{no abstract}, 1\,327 as \emph{other question},
1\,218 as \emph{no evidence}, 803 as \emph{dataset record} and 26 as
duplicates of cited sources.  The criteria are these: a record is
excluded if it is one of the following.
\begin{description}
  \item[No abstract] It has no abstract, and none could be found from
    its DOI.
  \item[Not interfaces] It is not about people using computer
    interfaces or research computing: another meaning of shell, menu,
    command, script, terminal or icon, or another science.
  \item[Tool only] It presents a software tool, system or interface
    design and offers no evidence on users' efficiency, learning or the
    reproducibility of their work.
  \item[No evidence] It is an opinion, tutorial, textbook, editorial or
    overview without data or a stated method, and not best-practice
    advice on reproducibility.
  \item[Other question] It is human--computer interaction or research
    computing on a different question, such as menu layout, input
    devices, voice or natural-language input, or code quality in
    general.
  \item[Dataset record] It is a dataset, supplementary file,
    replication package or erratum rather than a study.
\end{description}
\end{onlyin}
\begin{onlyin}{swedish}
Sökningarna gav 10\,405 poster, 9\,744 efter att poster för samma verk
från olika databaser slagits samman.  En post bedöms på sin
sammanfattning, så en post utan sammanfattning utesluts om ingen kan
hittas.  Efter sökningarna saknade 1\,955 poster sammanfattning.  Tre av
dem var citerade källor vars dubblettpost hade sammanfattningen, som vi
kopierade, och fem var dubbletter av citerade källor.  För de övriga
1\,947 slog vi upp sammanfattningen via DOI, i Semantic Scholar och
sedan i Crossref: 464 fick en (393 från Semantic Scholar, 71 från
Crossref), och de återstående 1\,483, varav 557 saknade DOI, uteslöts
som \emph{ingen sammanfattning}.  Alla andra poster klassades mot en
skriven formulering av forskningsfrågan av en språkmodell (Claude
Sonnet~5.5) utifrån titel och sammanfattning, som stödjande,
kvalificerande, angränsande ämne eller utanför ämnet, och varje utesluten
post fick ett av de namngivna kriterierna nedan.  Sedan läste vi
titlarna på alla poster som modellen behöll som stödjande eller
kvalificerande och på dem som den uteslöt med en konfidens under 0,7,
och sammanfattningarna av dem som kunde röra påståendet.  Därefter läste
vi sammanfattningen av varje post som fortfarande var märkt stödjande
eller kvalificerande och ännu inte citerad, 154 verk: 69 bekräftades och
citeras nedan, 77 märktes om och uteslöts, sex var ytterligare poster
för citerade verk, och två vars författare inte kunde kontrolleras
citeras inte.  Sammanlagt registrerade vi 272 egna beslut utöver regeln
om \emph{ingen sammanfattning}: de 121 citerade källorna behölls med det
tema som de citeras under, och 151 poster uteslöts, bland dem 26
dubbletter av citerade källor.  Av de 10\,405 posterna behölls 210 (66
stödjande, 66 kvalificerande, 78 om ett angränsande ämne) och 10\,195
uteslöts: 2\,670 som \emph{inte gränssnitt}, 2\,668 som \emph{bara
verktyg}, 1\,483 som \emph{ingen sammanfattning}, 1\,327 som \emph{annan
fråga}, 1\,218 som \emph{inga belägg}, 803 som \emph{datamängdspost} och
26 som dubbletter av citerade källor.  Kriterierna är dessa: en post
utesluts om den är något av följande.
\begin{description}
  \item[Ingen sammanfattning] Den saknar sammanfattning, och ingen
    kunde hittas via dess DOI.
  \item[Inte gränssnitt] Den handlar inte om människor som använder
    datorgränssnitt eller om datoranvändning i forskning: en annan
    betydelse av shell, menu, command, script, terminal eller icon,
    eller en annan vetenskap.
  \item[Bara verktyg] Den presenterar ett programverktyg, ett system
    eller en gränssnittsdesign och ger inga belägg om användarnas
    effektivitet, inlärning eller hur upprepbart deras arbete är.
  \item[Inga belägg] Den är en åsikt, en handledning, en lärobok, en
    ledare eller en översikt utan data eller angiven metod, och inte råd
    om bästa praxis för upprepbarhet.
  \item[Annan fråga] Den gäller människa--datorinteraktion eller
    datoranvändning i forskning men en annan fråga, som menylayout,
    inmatningsenheter, röst- eller naturligt språk-inmatning eller
    kodkvalitet i allmänhet.
  \item[Datamängdspost] Den är en datamängd, en bilagefil, ett
    replikationspaket eller en rättelse snarare än en studie.
\end{description}
\end{onlyin}

\paragraph{\ensv{Reading.}{Läsning.}} % sync english=0b0b47 swedish=47095f
\begin{onlyin}{english} % sync english=896837 swedish=05e628
We read in full every cited source that we could obtain: the report of
\textcite{Margono1995Study}, the paper of \textcite{Chen2007Comparing},
the
survey of \textcite{Voronkov2019System}, the guides of
\textcite{Sandve2013Ten}, \textcite{Wilson2014Best} and
\textcite{Wilson2017Good}, the audits of \textcite{Ziemann2016Gene},
\textcite{Abeysooriya2021Gene}, \textcite{Trisovic2022Large} and
\textcite{Saju2025Computational}, the comparison of
\textcite{Magar2025Data} and the interview study of
\textcite{Doyle2007Why}.  Most interface studies of the 1980s and
1990s appeared with publishers that refused the university's network
access or stopped it with a bot check; every other cited source was
read in its abstract only.
\end{onlyin}
\begin{onlyin}{swedish}
Vi läste i sin helhet varje citerad källa som vi kunde få tag på:
rapporten av \textcite{Margono1995Study}, artikeln av
\textcite{Chen2007Comparing}, enkäten av \textcite{Voronkov2019System},
guiderna av \textcite{Sandve2013Ten}, \textcite{Wilson2014Best} och
\textcite{Wilson2017Good}, granskningarna av \textcite{Ziemann2016Gene},
\textcite{Abeysooriya2021Gene}, \textcite{Trisovic2022Large} och
\textcite{Saju2025Computational}, jämförelsen av
\textcite{Magar2025Data} och intervjustudien av \textcite{Doyle2007Why}.
De flesta gränssnittsstudierna från 1980- och 1990-talen kom ut hos
förlag som nekade universitetets nätverk åtkomst eller stoppade det med
en kontroll mot robotar; varje annan citerad källa lästes bara i sin
sammanfattning.
\end{onlyin}

\section{\ensv{Results}{Resultat}} % sync english=219c4a swedish=6e45f9
\label{sec:cli-efficiency-results}

\paragraph{\ensv{Efficiency.}{Effektivitet.}} % sync english=e29307 swedish=05a6b6
\begin{onlyin}{english} % sync english=423fb1 swedish=b2439b
The experiments that compare command interfaces with menu, iconic or
direct-manipulation interfaces do not show a general advantage for
either side.  The broadest of them measured 76 users with seven systems
and concluded that \textcquote{Whiteside1985User}{interface style is
not related to performance or preference (but careful design is)}.
\Textcite{Hauptmann1983Comparison} gave one program a command, a menu
and a natural-language interface and found \textcquote
{Hauptmann1983Comparison}{\textins*{n}o significant differences
\textelp{} for time, error and attitude measures}.  No difference is
also the result for business students \autocite{Lindner1998Command},
for the workload of a drawing task except under the heaviest secondary task
\autocite{Eberleh1992Thinking}, and for a computer-aided design system
with and without menus \autocite{Roy1992Evaluation}.  Where a side won,
it was sometimes the command side.  In a computer-aided design system
\textcquote{Antin1988Empirical}{command entry produced performance
that was superior to the other two modes}, for experienced and novice
users alike, although users preferred a mode that combined commands and
menus; students who wrote essays with a command-driven word processor
got better grades than those with a mouse and pull-down menus
\autocite{Mahach1995Effects}.  More often it was the graphical side, by
\textcquote{Karat1987Evaluating}{large differences in performance} in
one experiment, and mainly for learners and newcomers: novices
manipulated files faster and with fewer errors on a direct-manipulation
interface than with typed commands \autocite{Margono1995Study};
learners on a Macintosh \textcquote{Davis1993Training}{performed
substantially better than those using the command-based interface},
DOS, in two experiments \autocite{Davis1992Experimental}; secretaries
learned a word processor more easily by direct manipulation
\autocite{Jordan1992Claims}; graduate students did better with the
menu version of a statistics package than with its command version
\autocite{Schultz1998Effects}; users with a command language alone did
worse than those who also had menus \autocite{Taylor1987Experimental};
and 98 nurses were faster, made fewer errors and learned faster with a
new graphical interface than with a legacy text-based one
\autocite{Staggers2000Comparing}, as were physicians with a graphical
patient record \autocite{Rodriguez2002Study}.

Further experiments, read in their abstracts, point the same mixed way.
Typed commands did better in some: they beat menu and hybrid modes on a
space-station interface for experienced and inexperienced users alike
\autocite{Antin1986Evaluation}, physicians entered orders faster with a
command-line pathway \autocite{Lovis2001Evaluation}, users of a command
interface performed better in a simulation task
\autocite{Schar1997Effect}, students needed less training with native
commands than with a menu front end to a bibliographic database, which
did not change what they retrieved \autocite{Sullivan1990End}, and
students querying in SQL were faster though they found it harder than a
graphical query language \autocite{Jones1996Articulating}.  The
graphical side did better in others: menus and choice systems were
easier than commands in a decision-support system
\autocite{Waring1992Factors}, a diagrammatic query language was easier
than SQL \autocite{Catarci1995Diagrammatic}, direct manipulation was
faster in a laboratory experiment \autocite{Lee2008Effects}, a GUI
helped even skilled programmers of a virtual-reality tool
\autocite{Fellmann2007Command}, and operators of milling machines
preferred direct manipulation \autocite{TorresChazaro1992Design}.  In
yet others the result was a trade-off: preference and performance went
together and depended on experience \autocite{Kissel1997Influence},
better usability did not mean better performance
\autocite{Chhanabhai2004First}, and a command language proved flexible
but hard to debug where direct manipulation gave a faster first
prototype \autocite{Kjeldskov2008Interaction}.

Expertise changes the picture, but not in one direction.  In a dental
record system \textcquote{Chen2007Comparing}{the GUI interface was not
better than the TUI for expert users. GUI interface was better for
novice users}, and experienced users of a character-based word
processor worked as fast as users of its graphical versions, mostly
through function keys \autocite{dYdewalle1995Graphical}.  Against this,
a desktop graphical interface beat a character-based menu interface
for beginners and experts alike, and experts with the graphical
interface needed 51~per~cent less time
\autocite{Rauterberg1992Empirical}, though that comparison is of menus,
not of typed commands.  How much a command line asks of its
user also depends on the user: experts with a weak associative memory
performed with commands like novices \autocite{Westerman1997Individual}.
The one experiment that follows the change with practice found the
pattern the claim assumes, for typed commands in general rather than a
shell.  Its authors note that \textcquote{Remington2016With}{\textins*
{i}t is widely believed that a graphical user interface (GUI) is
superior to a command line interface (CLI) for novice users, but less
efficient than the CLI after practice. However, there appears to be no
detailed study of the crossover interaction that this implies}; in
their experiment, keystroke commands overtook menu selection
\textcquote{Remington2016With}{after approximately 200 responses}.
Inside graphical programs, typed keyboard shortcuts are likewise the
most efficient way to give a command, and those who use them finish
tasks faster \autocite{Jardina2009Keyboard}, yet even experienced users
seldom take them up \autocite{Lane2005Hidden,Tak2013Satisficing}: an
efficiency that exists is not the same as one that is used.

Further studies agree that typed commands pay off once known and that
few make the switch: function-key commands sped up menu access
\autocite{Hammontree1991Communicating}, grouped hotkeys were among the
fastest ways to pick a command \autocite{Odell2004Toolglasses}, and
users led towards hotkeys did better \autocite{Scarr2011Dips}; yet many
users, experienced ones too, keep to menus
\autocite{Giannisakis2019Promoting}, keyboard and mouse entry did not
differ when the menu stayed visible \autocite{Davies1989Effects},
shortcuts can slow users down \autocite{Howes2000Trouble}, and what
users prefer predicts their performance poorly
\autocite{Heasly1988Comparative}.

None of these experiments tests the case of \cref{sec:cli-both}, a task
repeated over many files and saved as a script.  One recent comparison
of design workflows reports that \textcquote{Saddi2026Human}
{automation, scalability, and repeatability are more likely to be
achieved with CLI-based workflows}, but its abstract gives neither
sample nor numbers.  Beyond it we found only what practitioners
report.  System administrators, the users who
do such work daily, mostly use command lines although graphical tools
are offered, and \textcquote{Takayama2006Trust}{many of these power
users perceive CLIs as more effective than GUIs}; interviews and a
survey of 331 administrators describe a
\textcquote{Mahendiran2014Understanding}{culture of command-line
interface (CLI) use}.  In a survey of 300 administrators of
firewalls, \textcquote[8]{Voronkov2019System}{\textins*{a}ll use cases
associated with entering many similar rules on one computer or bulk
changes on several computers at once require the use of a CLI}, and
their remarks on the strengths of command lines named automation 64
times out of 309; yet in the same survey
\textcquote[1]{Voronkov2019System}{only 32\% of the respondents prefer
CLIs for managing firewalls, while this figure reaches 60\% for GUIs}.
A further survey sums up the trade-off in the administrators' own
terms: command-line software \textcquote{Murillo2014Empowering}{is
efficient but error prone}.  These are statements of experience and
preference, not measurements.

In clinical genomics, bioinformaticians preferred command lines where
geneticists and counsellors wanted a graphical layout
\autocite{Shyr2015Dynamic}; in flow cytometry, scripted analysis is
reported to take a minute where manual analysis takes an expert hours
\autocite{Brinkman2019Improving}.
\end{onlyin}
\begin{onlyin}{swedish}
De experiment som jämför kommandogränssnitt med meny-, ikon- eller
direktmanipulationsgränssnitt visar ingen allmän fördel för någon sida.
Det bredaste av dem mätte 76 användare med sju system och drog
slutsatsen att \foreigntextcquote{english}{Whiteside1985User}{interface
style is not related to performance or preference (but careful design
is)}.  \Textcite{Hauptmann1983Comparison} gav ett program ett
kommando-, ett meny- och ett naturligt språk-gränssnitt och fann
\foreigntextcquote{english}{Hauptmann1983Comparison}{\textins*{n}o
significant differences \textelp{} for time, error and attitude
measures}.  Ingen skillnad är också resultatet för ekonomistudenter
\autocite{Lindner1998Command}, för arbetsbelastningen i en rituppgift
utom under den tyngsta sidouppgiften \autocite{Eberleh1992Thinking}, och
för ett system för datorstödd konstruktion med och utan menyer
\autocite{Roy1992Evaluation}.  Där en sida vann var det ibland
kommandosidan.  I ett system för datorstödd konstruktion
\foreigntextcquote{english}{Antin1988Empirical}{command entry produced
performance that was superior to the other two modes}, för erfarna
användare och nybörjare lika, fastän användarna föredrog ett läge som
förenade kommandon och menyer; studenter som skrev uppsatser med en
kommandostyrd ordbehandlare fick bättre betyg än de som hade mus och
rullgardinsmenyer \autocite{Mahach1995Effects}.  Oftare var det den
grafiska sidan, med
\foreigntextcquote{english}{Karat1987Evaluating}{large differences in
performance} i ett experiment, och främst för de som lär sig och för
nykomlingar: nybörjare hanterade filer snabbare och med färre fel i ett
gränssnitt för direkt manipulation än med skrivna kommandon
\autocite{Margono1995Study}; de som lärde sig på en Macintosh
\foreigntextcquote{english}{Davis1993Training}{performed substantially
better than those using the command-based interface}, DOS, i två
experiment \autocite{Davis1992Experimental}; sekreterare lärde sig en
ordbehandlare lättare genom direkt manipulation
\autocite{Jordan1992Claims}; forskarstudenter klarade sig bättre med
menyversionen av ett statistikpaket än med dess kommandoversion
\autocite{Schultz1998Effects}; användare med bara ett kommandospråk
klarade sig sämre än de som också hade menyer
\autocite{Taylor1987Experimental}; och 98 sjuksköterskor var snabbare,
gjorde färre fel och lärde sig snabbare med ett nytt grafiskt gränssnitt
än med ett äldre textbaserat \autocite{Staggers2000Comparing}, liksom
läkare med en grafisk patientjournal \autocite{Rodriguez2002Study}.

Ytterligare experiment, lästa i sina sammanfattningar, pekar åt samma
blandade håll.  Skrivna kommandon klarade sig bättre i några: de slog
meny- och hybridlägen i ett gränssnitt för en rymdstation för både
erfarna och oerfarna användare \autocite{Antin1986Evaluation}, läkare
förde in ordinationer snabbare med en kommandoradsväg
\autocite{Lovis2001Evaluation}, användare av ett kommandogränssnitt
presterade bättre i en simuleringsuppgift \autocite{Schar1997Effect},
studenter behövde mindre träning med de egna kommandona än med ett
menygränssnitt framför en bibliografisk databas, vilket inte ändrade vad
de hittade \autocite{Sullivan1990End}, och studenter som ställde frågor
i SQL var snabbare fastän de tyckte att det var svårare än ett grafiskt
frågespråk \autocite{Jones1996Articulating}.  Den grafiska sidan klarade
sig bättre i andra: menyer och valsystem var lättare än kommandon i ett
beslutsstödsystem \autocite{Waring1992Factors}, ett diagrambaserat
frågespråk var lättare än SQL \autocite{Catarci1995Diagrammatic}, direkt
manipulation var snabbare i ett laboratorieexperiment
\autocite{Lee2008Effects}, ett GUI hjälpte även skickliga programmerare
av ett verktyg för virtuell verklighet \autocite{Fellmann2007Command},
och operatörer av fräsmaskiner föredrog direkt manipulation
\autocite{TorresChazaro1992Design}.  I ytterligare andra blev resultatet
en avvägning: preferens och prestation gick hand i hand och berodde på
erfarenhet \autocite{Kissel1997Influence}, bättre användbarhet innebar
inte bättre prestation \autocite{Chhanabhai2004First}, och ett
kommandospråk visade sig flexibelt men svårt att felsöka där direkt
manipulation gav en snabbare första prototyp
\autocite{Kjeldskov2008Interaction}.

Expertis ändrar bilden, men inte åt ett håll.  I ett journalsystem för
tandvård \foreigntextcquote{english}{Chen2007Comparing}{the GUI
interface was not better than the TUI for expert users. GUI interface
was better for novice users}, och erfarna användare av en teckenbaserad
ordbehandlare arbetade lika snabbt som användare av dess grafiska
versioner, mest genom funktionstangenter
\autocite{dYdewalle1995Graphical}.  Mot detta står att ett grafiskt
skrivbordsgränssnitt slog ett teckenbaserat menygränssnitt för både
nybörjare och experter, och experter med det grafiska gränssnittet
behövde 51~procent kortare tid \autocite{Rauterberg1992Empirical},
fastän den jämförelsen gäller menyer, inte skrivna kommandon.  Hur
mycket en kommandorad kräver av sin användare beror också på användaren:
experter med svagt associativt minne presterade med kommandon som
nybörjare \autocite{Westerman1997Individual}.  Det enda experimentet som
följer förändringen med övning fann det mönster som påståendet förutsätter,
för skrivna kommandon i allmänhet snarare än ett skal.  Dess författare
noterar att \foreigntextcquote{english}{Remington2016With}{\textins*{i}t
is widely believed that a graphical user interface (GUI) is superior to
a command line interface (CLI) for novice users, but less efficient than
the CLI after practice. However, there appears to be no detailed study
of the crossover interaction that this implies}; i deras experiment gick
tangentkommandon förbi menyval
\foreigntextcquote{english}{Remington2016With}{after approximately 200
responses}.  Inuti grafiska program är skrivna kortkommandon på samma
sätt det effektivaste sättet att ge ett kommando, och de som använder
dem gör uppgifter snabbare \autocite{Jardina2009Keyboard}, men även
erfarna användare tar sällan till dem
\autocite{Lane2005Hidden,Tak2013Satisficing}: en effektivitet som finns
är inte detsamma som en som används.

Ytterligare studier är eniga om att skrivna kommandon lönar sig när man
kan dem och att få byter: kommandon på funktionstangenter snabbade upp
åtkomsten till menyer \autocite{Hammontree1991Communicating}, grupperade
snabbtangenter var bland de snabbaste sätten att välja ett kommando
\autocite{Odell2004Toolglasses}, och användare som leddes mot
snabbtangenter klarade sig bättre \autocite{Scarr2011Dips}; men många
användare, också erfarna, håller sig till menyer
\autocite{Giannisakis2019Promoting}, inmatning med tangentbord och mus
skilde sig inte åt när menyn förblev synlig \autocite{Davies1989Effects},
kortkommandon kan göra användare långsammare \autocite{Howes2000Trouble},
och vad användare föredrar förutsäger deras prestation dåligt
\autocite{Heasly1988Comparative}.

Inget av experimenten prövar fallet i \cref{sec:cli-both}, en uppgift
som upprepas över många filer och sparas som ett skript.  En nyligen
gjord jämförelse av designarbetsflöden rapporterar att
\foreigntextcquote{english}{Saddi2026Human}{automation, scalability,
and repeatability are more likely to be achieved with CLI-based
workflows}, men dess sammanfattning ger varken urval eller siffror.
Utöver den hittade vi bara vad yrkesverksamma rapporterar.
Systemadministratörer, de användare som gör sådant arbete dagligen,
använder mest kommandorader fastän grafiska verktyg erbjuds, och
\foreigntextcquote{english}{Takayama2006Trust}{many of these power users
perceive CLIs as more effective than GUIs}; intervjuer och en enkät med
331 administratörer beskriver en
\foreigntextcquote{english}{Mahendiran2014Understanding}{culture of
command-line interface (CLI) use}.  I en enkät med 300 administratörer
av brandväggar
\foreigntextcquote{english}[8]{Voronkov2019System}{\textins*{a}ll use
cases associated with entering many similar rules on one computer or
bulk changes on several computers at once require the use of a CLI},
och deras kommentarer om kommandoradernas styrkor nämnde automatisering
64 gånger av 309; ändå
\foreigntextcquote{english}[1]{Voronkov2019System}{only 32\% of the
respondents prefer CLIs for managing firewalls, while this figure
reaches 60\% for GUIs}.  En ytterligare enkät sammanfattar avvägningen i
administratörernas egna ord: kommandoradsprogram
\foreigntextcquote{english}{Murillo2014Empowering}{is efficient but
error prone}.  Det är uttalanden om erfarenhet och preferens, inte
mätningar.

Inom klinisk genomik föredrog bioinformatiker kommandorader där
genetiker och genetiska vägledare ville ha en grafisk layout
\autocite{Shyr2015Dynamic}; inom flödescytometri rapporteras skriptad
analys ta en minut där manuell analys tar en expert timmar
\autocite{Brinkman2019Improving}.
\end{onlyin}

\paragraph{\ensv{Reproducibility.}{Upprepbarhet.}} % sync english=32663b swedish=188478
\begin{onlyin}{english} % sync english=214354 swedish=06b48e
The advice is clear and widely repeated.  Among rules for reproducible
computational research, \textcquote[2]{Sandve2013Ten}{\textins*{w}henever
possible, rely on the execution of programs instead of manual
procedures to modify data. Such manual procedures are not only
inefficient and error-prone, they are also difficult to reproduce. If
working at the UNIX command line, manual modification of files can
usually be replaced by the use of standard UNIX commands or small
custom scripts.}  Best practices for scientific computing advise
scientists to \textcquote[2]{Wilson2014Best}{make the computer repeat
tasks (2a) and save recent commands in a file for re-use (2b)}, since
those who repeat commands or clicks by hand will
\textcquote[2]{Wilson2014Best}{sooner or later} make mistakes; the
\enquote{good enough} practices that followed advise you to
\textcquote[5]{Wilson2017Good}{write scripts for every stage of data
processing}.  The advice rests on experience, and says
so: the best practices cite two practitioner textbooks for the remark
that the command line's \enquote{do this again} \textcquote[2]
{Wilson2014Best}{saves time and reduces errors}, and the \enquote{good
enough} practices draw on \textcquote[2]{Wilson2017Good}{our personal
experience, and \textelp{} the experiences of the thousands of people
who have taken part in Software Carpentry and Data Carpentry
workshops}.  The counter query on the advice (C6) returned no
evaluation of it.

The same advice recurs across fields.  In ecology, scripts replicate an
analysis exactly where point-and-click packages do not
\autocite{Borregaard2016Towards}; in clinical and laboratory research,
manual cleaning, point-and-click analysis and spreadsheets are advised
against in favour of scripts
\autocite{Kaizer2026Statistical,Obstfeld2026Reproducible,%
Simons2019Reproducible}; in genomics, analyses are to run as automated,
non-interactive pipelines
\autocite{Vos2016Ten,Ziemann2023Recipe,Satagopam2025Best}; data
processing is to be built from Unix tools with logging
\autocite{Spinellis2024Engineering}; and further guides recommend code
that records and regenerates the whole analysis
\autocite{Kelley2008R,Valle2012Computer,RodriguezSanchez2016Ciencia,%
Hutton2016Most,Petyuk2019Reproducibility,Kalenkovich2021Reproducible,%
Finley2026Scripts,Schlogl2026Overcoming}.  One pilot converted ten
psychology papers to single scripts and judged it a success
\autocite{Baker2024Reproduceme}.

Two kinds of evidence go beyond advice, and they point in different
directions.  Manual work in a graphical data tool does introduce errors
at scale: in leading genomics journals \textcquote{Ziemann2016Gene}
{approximately one-fifth of papers with supplementary Excel gene lists
contain erroneous gene name conversions}, made by the spreadsheet's
default settings, and five years later the share had grown to
\textcquote{Abeysooriya2021Gene}{30.9\% (3,436/11,117) of articles with
supplementary Excel gene lists}.  But a script is no guarantee either:
of more than 9000 R files published with research data,
\textcquote{Trisovic2022Large}{74\% of R files failed to complete
without error in the initial execution, while 56\% failed when code
cleaning was applied}.  Audits in other fields agree: a quarter of the
R scripts shared on one platform ran without error even after their
environments were rebuilt \autocite{Saju2025Computational}, 7~per~cent
of ecology papers provided code that ran \autocite{Kellner2024Functional},
the results of 15~per~cent of statistics papers on brain imaging data
could be reproduced \autocite{Xiong2022State}, of computational physics
papers whose code was obtained two thirds could not easily be
regenerated \autocite{Stodden2018Enabling}, and in hydrology
1.6~per~cent of the tested articles could \autocite{Stagge2019Assessing}.
And a graphical tool can keep the record that makes work repeatable:
the \enquote{good enough} practices name data-cleaning tools that
\textcquote[5]{Wilson2017Good}{provide a graphical user interface but
also automatically keep track of each step in the process}.  The only
comparison of the same analysis done in a spreadsheet and by script
that we found is a preprint in which its author did one analysis of
one data set both ways and found \textcquote{Magar2025Data}{that Excel
is only efficient and error-free to the extent of basic analysis and
simple outputs}.  We found no study in which someone else redid work
first done by script or by hand in a graphical program.

Further studies point both ways here too.  Excel's corruption of
supplementary data continues year after year
\autocite{DeRijk2018Opentsv}, a forensic reanalysis traced errors in
published gene signatures to simple data management
\autocite{Master2010Getting}, and none of ten hydrological models built
through a graphical interface could easily be reproduced
\autocite{Chawanda2020User}.  But code fails too: two thirds of the
example code on model cards did not run \autocite{Yao2026Understanding},
two of 19 notebooks from published papers were reproducible
\autocite{Kashiwa2026Study}, few benchmark submissions could be installed
as published \autocite{Lyu2026Code}, archived retrieval runs no longer
replicated after four years \autocite{Lin2020Reproducibility}, and
reproductions by a student, a hackathon, a review of geographic
information science and an audit of published tools succeeded only in
part \autocite{Oneka2018Evaluating,Willems2026Insights,%
Nust2022Infrastructures,Nong2023Open}.
\end{onlyin}
\begin{onlyin}{swedish}
Rådet är tydligt och upprepas ofta.  Bland regler för upprepbar
beräkningsforskning står
\foreigntextcquote{english}[2]{Sandve2013Ten}{\textins*{w}henever
possible, rely on the execution of programs instead of manual procedures
to modify data. Such manual procedures are not only inefficient and
error-prone, they are also difficult to reproduce. If working at the
UNIX command line, manual modification of files can usually be replaced
by the use of standard UNIX commands or small custom scripts.}  Bästa
praxis för vetenskapliga beräkningar råder forskare att
\foreigntextcquote{english}[2]{Wilson2014Best}{make the computer repeat
tasks (2a) and save recent commands in a file for re-use (2b)},
eftersom de som upprepar kommandon eller klick för hand
\foreigntextcquote{english}[2]{Wilson2014Best}{sooner or later} gör fel;
den \enquote{good enough}-praxis som följde råder dig att
\foreigntextcquote{english}[5]{Wilson2017Good}{write scripts for every
stage of data processing}.  Rådet vilar på erfarenhet, och säger det:
bästa praxis hänvisar till två läroböcker av yrkesverksamma för
anmärkningen att kommandoradens \enquote{gör det här igen}
\foreigntextcquote{english}[2]{Wilson2014Best}{saves time and reduces
errors}, och \enquote{good enough}-praxisen bygger på
\foreigntextcquote{english}[2]{Wilson2017Good}{our personal experience,
and \textelp{} the experiences of the thousands of people who have taken
part in Software Carpentry and Data Carpentry workshops}.  Motfrågan om
rådet (C6) gav ingen utvärdering av det.

Samma råd återkommer inom många fält.  Inom ekologi upprepar skript en
analys exakt där peka-och-klicka-paket inte gör det
\autocite{Borregaard2016Towards}; inom klinisk forskning och
laboratorieforskning avråds manuell rensning, peka-och-klicka-analys och
kalkylblad till förmån för skript
\autocite{Kaizer2026Statistical,Obstfeld2026Reproducible,%
Simons2019Reproducible}; inom genomik ska analyser köras som
automatiserade, icke-interaktiva flöden
\autocite{Vos2016Ten,Ziemann2023Recipe,Satagopam2025Best};
databearbetning ska byggas av Unixverktyg med loggning
\autocite{Spinellis2024Engineering}; och ytterligare guider
rekommenderar kod som registrerar och återskapar hela analysen
\autocite{Kelley2008R,Valle2012Computer,RodriguezSanchez2016Ciencia,%
Hutton2016Most,Petyuk2019Reproducibility,Kalenkovich2021Reproducible,%
Finley2026Scripts,Schlogl2026Overcoming}.  Ett pilotprojekt gjorde om
tio psykologiartiklar till var sitt skript och bedömde det som en
framgång \autocite{Baker2024Reproduceme}.

Två sorters belägg går utöver råd, och de pekar åt olika håll.  Manuellt
arbete i ett grafiskt dataverktyg ger fel i stor skala: i ledande
tidskrifter i genomik
\foreigntextcquote{english}{Ziemann2016Gene}{approximately one-fifth of
papers with supplementary Excel gene lists contain erroneous gene name
conversions}, orsakade av kalkylprogrammets standardinställningar, och
fem år senare hade andelen vuxit till
\foreigntextcquote{english}{Abeysooriya2021Gene}{30.9\%
(3,436/11,117) of articles with supplementary Excel gene lists}.  Men
ett skript är ingen garanti heller: av mer än 9000 R-filer som
publicerats med forskningsdata
\foreigntextcquote{english}{Trisovic2022Large}{74\% of R files failed to
complete without error in the initial execution, while 56\% failed when
code cleaning was applied}.  Granskningar inom andra fält håller med:
en fjärdedel av de R-skript som delats på en plattform gick att köra
utan fel även sedan deras miljöer byggts upp på nytt
\autocite{Saju2025Computational}, 7~procent av artiklarna i ekologi
tillhandahöll kod som gick att köra \autocite{Kellner2024Functional},
resultaten i 15~procent av statistikartiklarna om hjärnavbildningsdata
gick att återskapa \autocite{Xiong2022State}, av artiklarna i
beräkningsfysik vars kod kunde fås gick två tredjedelar inte lätt att
återskapa \autocite{Stodden2018Enabling}, och inom hydrologi gick det
för 1,6~procent av de prövade artiklarna
\autocite{Stagge2019Assessing}.  Och ett grafiskt verktyg kan föra den
logg som gör arbete upprepbart: \enquote{good enough}-praxisen nämner
verktyg för datarensning som
\foreigntextcquote{english}[5]{Wilson2017Good}{provide a graphical user
interface but also automatically keep track of each step in the
process}.  Den enda jämförelse av samma analys gjord i ett kalkylblad och
med skript som vi hittade är en förhandspublikation där författaren
gjorde en analys av en datamängd på båda sätten och fann
\foreigntextcquote{english}{Magar2025Data}{that Excel is only efficient
and error-free to the extent of basic analysis and simple outputs}.  Vi
hittade ingen studie där någon annan gjorde om arbete som först gjorts
med skript eller för hand i ett grafiskt program.

Ytterligare studier pekar åt båda hållen även här.  Excels förvanskning
av bilagedata fortsätter år efter år \autocite{DeRijk2018Opentsv}, en
kriminalteknisk omanalys spårade fel i publicerade gensignaturer till
enkel datahantering \autocite{Master2010Getting}, och ingen av tio
hydrologiska modeller som byggts genom ett grafiskt gränssnitt gick lätt
att återskapa \autocite{Chawanda2020User}.  Men kod fallerar också: två
tredjedelar av exempelkoden på modellkort gick inte att köra
\autocite{Yao2026Understanding}, två av 19 notebooks från publicerade
artiklar gick att återskapa \autocite{Kashiwa2026Study}, få bidrag till
jämförelsetester gick att installera som de publicerats
\autocite{Lyu2026Code}, arkiverade sökkörningar gick inte längre att
upprepa efter fyra år \autocite{Lin2020Reproducibility}, och
återskapanden gjorda av en student, ett hackathon, en översikt inom
geografisk informationsvetenskap och en granskning av publicerade
verktyg lyckades bara delvis \autocite{Oneka2018Evaluating,%
Willems2026Insights,Nust2022Infrastructures,Nong2023Open}.
\end{onlyin}

\paragraph{\ensv{Novices.}{Nybörjare.}} % sync english=418c40 swedish=545f0f
\begin{onlyin}{english} % sync english=6ff930 swedish=efe6f1
That newcomers find a command line harder at first is supported, but
not without exceptions.  Novices made more errors and took longer with
typed commands than with direct manipulation, and their errors were
those of a typed language: in the words of the experimenters,
\textcquote{Margono1995Study}{\textins*{f}orgetting the punctuation,
adding unnecessary spaces, and omitting the parameters in such commands
as copy and rename were among the common problems the subjects had on
the IBM PC}.  Composing complex UNIX commands strained novices so that
\textcquote{Sohn1997Cognitive}{prompting novices with required
knowledge is not sufficient to obtain correct performance; they must
receive additional assistance that reduces working memory demands}.
The newer graphical interface was also the faster one to learn for the
nurses \autocite{Staggers2000Comparing}, the learners on the Macintosh
did better than those on DOS without finding it easier
\autocite{Davis1993Training}, and the menus helped the graduate
students to find their errors \autocite{Schultz1998Effects}; 12
undergraduates preferred a graphical library catalogue to a
command-line one \autocite{Meadows2011Graphical}; and in a first
programming course, visual tools \textcquote{Dillon2012Studying}{could
provide a lower learning curve for novices, while command line tools
could broaden their understanding of programming}.  Against this stand
studies in which the command side held up.  High-school students with
no computer experience made fewer errors learning the UNIX command
language than an English-like one, but asked for help more often a
week later, so that the authors conclude that \textcquote
{DeLeon1983Trouble}{even though the UNIX command language may not be
harder for novice users to learn, it is probably more difficult for
them to use}.  People learning UNIX command sequences did better with a
command line than with a menu, \textcquote{Durham1998Learning}
{especially during retention testing}, while for a milling machine the
opposite held: a week's break slowed operators on the command language
but not on direct manipulation \autocite{TorresChazaro1992Comparison}.
In one study learning a direct-manipulation interface took longer than
learning a command language, which repaid it on complex tasks
\autocite{Reinhard1991Effect}, and in two first programming courses
more students finished a task with a command-line editor than with a
web environment, though they felt less at ease with it
\autocite{Cooper2020Exploring}.  How commands are named matters more
than whether they are natural words: names that told operations with
different syntax apart \textcquote{Landauer1983Natural}{greatly reduced
difficulty}.  That many students arrive without any experience of a
command line is reported by teachers \autocite{Kendon2016Unix}, and
among computing teachers \textcquote{Doyle2007Why}{Unix as a command
line interface superior to GUIs} is itself one of their conceptions of
Unix, a less developed one than Unix as a tool for solving problems.

Further learning studies, read in their abstracts, mostly agree.
Direct manipulation was easier to learn or more motivating than a
command style
\autocite{Davis1994Effects,Davis2001Mediating,Davis1998Effect,%
Wiedenbeck1997Influence}, a menu-to-command transition helped newcomers
to UNIX \autocite{Badre1991Effects}, newcomers to a command interface
met errors with no sign of being close \autocite{Lewis2003Menus},
prompts did not suffice for composing complex UNIX commands
\autocite{Doane1992Prompt}, students tended to learn more with a command
line given graphical support \autocite{Herrera2005Training},
computer-naive users did best with menus first and moved to commands
after 16 to 20 hours \autocite{Gilfoil1982Warming}, a guided system was
more efficient than a command line for novices
\autocite{Westendorp1997Guided}, a GUI gave authors of tutoring systems
a lower bar for entry \autocite{Devasani2012Evaluation}, a
command-language database interface led students to skip decision
factors \autocite{Chen2001Investigation}, learners found a graphical
library database easier \autocite{Gaff1994Gui}, and users of
command-line computing clusters valued graphical desktops for getting in
\autocite{Thota2019Research}.  Against these, a command interface built
from novices' own commands came to recognise most of them
\autocite{Good1984Building}, students who began with a command-line
editor moved more easily to a graphical one than the reverse
\autocite{Dillon2012Mental}, the extremes of assisting features were a
problem in either kind of environment \autocite{Dillon2012Comparing}, a
rich web interface rather than the simplest point-and-click one was
preferred \autocite{MarquezBarja2018Assessing}, and some of novices'
apparent misconceptions belong to the representation rather than to
them \autocite{Masson1988Misconceived}.
\end{onlyin}
\begin{onlyin}{swedish}
Att nykomlingar till en början tycker att en kommandorad är svårare har
stöd, men inte utan undantag.  Nybörjare gjorde fler fel och tog längre
tid med skrivna kommandon än med direkt manipulation, och deras fel var
ett skrivet språks: med försöksledarnas ord
\foreigntextcquote{english}{Margono1995Study}{\textins*{f}orgetting the
punctuation, adding unnecessary spaces, and omitting the parameters in
such commands as copy and rename were among the common problems the
subjects had on the IBM PC}.  Att sätta ihop komplexa UNIX-kommandon
belastade nybörjare så att
\foreigntextcquote{english}{Sohn1997Cognitive}{prompting novices with
required knowledge is not sufficient to obtain correct performance; they
must receive additional assistance that reduces working memory demands}.
Det nyare grafiska gränssnittet var också det som gick snabbast att lära
sig för sjuksköterskorna \autocite{Staggers2000Comparing}, de som lärde
sig på Macintosh klarade sig bättre än de på DOS utan att tycka att det
var lättare \autocite{Davis1993Training}, och menyerna hjälpte
forskarstudenterna att hitta sina fel \autocite{Schultz1998Effects}; 12
studenter föredrog en grafisk bibliotekskatalog framför en på
kommandorad \autocite{Meadows2011Graphical}; och i en första
programmeringskurs fann man att visuella verktyg
\foreigntextcquote{english}{Dillon2012Studying}{could provide a lower
learning curve for novices, while command line tools could broaden their
understanding of programming}.  Mot detta står
studier där kommandosidan höll.  Gymnasieelever utan datorerfarenhet
gjorde färre fel när de lärde sig UNIX kommandospråk än ett
engelskliknande, men bad oftare om hjälp en vecka senare, så att
författarna drar slutsatsen att
\foreigntextcquote{english}{DeLeon1983Trouble}{even though the UNIX
command language may not be harder for novice users to learn, it is
probably more difficult for them to use}.  De som lärde sig följder av
UNIX-kommandon klarade sig bättre med en kommandorad än med en meny,
\foreigntextcquote{english}{Durham1998Learning}{especially during
retention testing}, medan motsatsen gällde för en fräsmaskin: en veckas
uppehåll gjorde operatörerna långsammare med kommandospråket men inte
med direkt manipulation \autocite{TorresChazaro1992Comparison}.  I en
studie tog det längre tid att lära sig ett gränssnitt för direkt
manipulation än ett kommandospråk, som betalade sig på komplexa uppgifter
\autocite{Reinhard1991Effect}, och i två första programmeringskurser
gjorde fler studenter klart en uppgift med en kommandoradseditor än med
en webbmiljö, fastän de kände sig mindre hemma i den
\autocite{Cooper2020Exploring}.  Hur kommandon heter spelar större roll
än om de är naturliga ord: namn som skilde operationer med olika syntax
åt \foreigntextcquote{english}{Landauer1983Natural}{greatly reduced
difficulty}.  Att många studenter kommer utan någon erfarenhet av en
kommandorad rapporteras av lärare \autocite{Kendon2016Unix}, och bland
datavetenskapslärare är
\foreigntextcquote{english}{Doyle2007Why}{Unix as a command line
interface superior to GUIs} själv en av deras uppfattningar om Unix, en
mindre utvecklad än Unix som ett verktyg för att lösa problem.

Ytterligare studier av inlärning, lästa i sina sammanfattningar, håller
mest med.  Direkt manipulation var lättare att lära sig eller mer
motiverande än en kommandostil
\autocite{Davis1994Effects,Davis2001Mediating,Davis1998Effect,%
Wiedenbeck1997Influence}, en övergång från menyer till kommandon
hjälpte nykomlingar till UNIX \autocite{Badre1991Effects}, nykomlingar
till ett kommandogränssnitt mötte fel utan tecken på att vara nära
\autocite{Lewis2003Menus}, uppmaningar räckte inte för att sätta ihop
komplexa UNIX-kommandon \autocite{Doane1992Prompt}, studenter tenderade
att lära sig mer med en kommandorad som fick grafiskt stöd
\autocite{Herrera2005Training}, datorovana användare klarade sig bäst
med menyer först och gick över till kommandon efter 16 till 20 timmar
\autocite{Gilfoil1982Warming}, ett guidat system var effektivare än en
kommandorad för nybörjare \autocite{Westendorp1997Guided}, ett GUI gav
författare av handledningssystem en lägre tröskel
\autocite{Devasani2012Evaluation}, ett databasgränssnitt med
kommandospråk fick studenter att hoppa över beslutsfaktorer
\autocite{Chen2001Investigation}, de som lärde sig tyckte att en
grafisk biblioteksdatabas var lättare \autocite{Gaff1994Gui}, och
användare av beräkningskluster med kommandorad värdesatte grafiska
skrivbord för att komma in \autocite{Thota2019Research}.  Mot dessa står
att ett kommandogränssnitt byggt av nybörjarnas egna kommandon kom att
känna igen de flesta av dem \autocite{Good1984Building}, att studenter
som började med en kommandoradseditor gick lättare över till en grafisk
än tvärtom \autocite{Dillon2012Mental}, att ytterligheterna av
hjälpfunktioner var ett problem i båda slagen av miljö
\autocite{Dillon2012Comparing}, att ett rikt webbgränssnitt snarare än
det enklaste peka-och-klicka-gränssnittet föredrogs
\autocite{MarquezBarja2018Assessing}, och att en del av nybörjares
skenbara missuppfattningar hör till representationen snarare än till dem
\autocite{Masson1988Misconceived}.
\end{onlyin}

\section{\ensv{Discussion and limitations}{Diskussion och begränsningar}} % sync english=634094 swedish=5d16d0
\label{sec:cli-efficiency-discussion}

\begin{onlyin}{english} % sync english=100379 swedish=7b57ba
The interface studies are old: most date from 1983 to 2000, when the
graphical interface was new and the command lines compared were those
of MS-DOS, early UNIX or special systems.  Several compare a legacy
interface with a new one \autocite{Staggers2000Comparing}, so that age
and care of design are confounded with interface style; within each
style, \textcite{Whiteside1985User} found large differences between
systems.  Several of the
\enquote{text-based} or \enquote{character-based} interfaces are
keyboard menus or forms rather than command lines
\autocite{Rauterberg1992Empirical,dYdewalle1995Graphical,%
  Chen2007Comparing},
and the shortcut studies, the crossover experiment among them, concern
keys inside graphical programs
\autocite{Remington2016With,Lane2005Hidden,Tak2013Satisficing}; we
cite them because they test the same argument, that typed commands pay
off for those who know them, not because they test a shell.  The
samples are small, from eight experts \autocite{Chen2007Comparing} to
98 nurses \autocite{Staggers2000Comparing}, and the practitioners'
reports are surveys of preference and opinion
\autocite{Takayama2006Trust,Voronkov2019System,Murillo2014Empowering}.
The audits of research code measure whether published scripts run
again, mostly in R; they do not compare scripts with manual work.

Our access to the sources limits the reading.  Most of the interface
studies are behind publishers that refused the university's network or
stopped it with a bot check, so we read them in abstract only, and an
abstract reports a finding without its conditions.  The rule that a
record without an abstract is excluded weighs most on the oldest
studies: 62 records that the model had kept from their titles alone
were excluded for want of an abstract.  Of these, 24 bear on the claim
by their titles and could not be obtained, most of them from 1982 to
1998; among them are studies of the errors Unix users make, of
learning UNIX, of the transition from menus to commands, of command
reuse in the shell, of interface style and problem solving, of keyboard
shortcuts, and 17 further comparisons of command and graphical
interfaces.  To these come the two studies that only the check by
meaning found.  The counts in \cref{tab:cli-efficiency-queries} carry
the limits of the databases.  DBLP could not be searched at all.
OpenAlex was searched in titles and abstracts, not in full text, and its
abstracts are incomplete, so it does not equal the same search in Web
of Science.  The screening was done by a language model, mostly from
abstracts, and we read only the titles of the records it kept or was
unsure of, and abstracts where a
title suggested a bearing; a relevant study it excluded with confidence
would be missed.  Of the 124 works that \cref{tab:sok-cli-efficiency}
lists as bearing on the claim, we cite 121; 52 were read in full or in
their abstracts with a quote for each claim they carry, and the other
69, the further studies of the results, only in their abstracts, each
with one quote that we checked against the abstract word for word.  The
three that are not cited are an earlier comparison of performance and
preference in the two kinds of environment whose abstract states no
result, a book chapter whose authors the metadata does not give, and a
preprint whose author is given only by a first name.  The study of
firewall administrators is cited in its first arXiv version, which is
the one we read; its published version carries a different title.
\end{onlyin}
\begin{onlyin}{swedish}
Gränssnittsstudierna är gamla: de flesta är från 1983 till 2000, när
det grafiska gränssnittet var nytt och de kommandorader som jämfördes
var MS-DOS, tidig UNIX eller specialsystem.  Flera jämför ett äldre
gränssnitt med ett nytt \autocite{Staggers2000Comparing}, så att ålder
och omsorg i utformningen blandas ihop med gränssnittsstil; inom varje
stil fann \textcite{Whiteside1985User} stora skillnader mellan system.
Flera av de \enquote{textbaserade} eller \enquote{teckenbaserade}
gränssnitten är tangentbordsmenyer eller formulär snarare än
kommandorader
\autocite{Rauterberg1992Empirical,dYdewalle1995Graphical,%
  Chen2007Comparing},
och studierna av kortkommandon, korsningsexperimentet bland dem, gäller
tangenter inuti grafiska program
\autocite{Remington2016With,Lane2005Hidden,Tak2013Satisficing}; vi
citerar dem eftersom de prövar samma argument, att skrivna kommandon
lönar sig för dem som kan dem, inte för att de prövar ett skal.  Urvalen
är små, från åtta experter \autocite{Chen2007Comparing} till 98
sjuksköterskor \autocite{Staggers2000Comparing}, och de yrkesverksammas
rapporter är enkäter om preferens och åsikt
\autocite{Takayama2006Trust,Voronkov2019System,Murillo2014Empowering}.
Granskningarna av forskningskod mäter om publicerade skript går att köra
igen, mest i R; de jämför inte skript med manuellt arbete.

Vår åtkomst till källorna begränsar läsningen.  De flesta
gränssnittsstudierna ligger hos förlag som nekade universitetets nätverk
eller stoppade det med en kontroll mot robotar, så vi läste dem bara i
sammanfattning, och en sammanfattning redovisar ett fynd utan dess
villkor.  Regeln att en post utan sammanfattning utesluts väger tyngst
på de äldsta studierna: 62 poster som modellen hade behållit enbart på
sina titlar uteslöts i brist på sammanfattning.  Av dem rör 24
påståendet enligt sina titlar och gick inte att få tag på, de flesta
från 1982 till 1998; bland dem finns studier av de fel som
Unixanvändare gör, av att lära sig UNIX, av övergången från menyer till
kommandon, av återanvändning av kommandon i skalet, av gränssnittsstil
och problemlösning, av kortkommandon, och 17 ytterligare jämförelser av
kommandogränssnitt och grafiska gränssnitt.  Till dem kommer de två
studier som bara kontrollen efter betydelse hittade.  Talen i
\cref{tab:cli-efficiency-queries} bär databasernas gränser.  DBLP kunde
inte genomsökas alls.  OpenAlex söktes i titlar och sammanfattningar,
inte i fulltext, och dess sammanfattningar är ofullständiga, så den är
inte likvärdig med samma sökning i Web of Science.  Gallringen gjordes
av en språkmodell, mest utifrån sammanfattningar, och vi läste bara
titlarna på de poster som den behöll eller var osäker på, och
sammanfattningar där en titel antydde att posten rörde påståendet; en
relevant studie som den uteslöt med säkerhet skulle missas.  Av de 124
verk som \cref{tab:sok-cli-efficiency} listar som relevanta för
påståendet citerar vi 121; 52 lästes i sin helhet eller i sina
sammanfattningar med ett citat för varje påstående de bär, och de andra
69, resultatens ytterligare studier, bara i sina sammanfattningar, var
och en med ett citat som vi kontrollerade mot sammanfattningen ord för
ord.  De tre som inte citeras är en tidigare jämförelse av prestation
och preferens i de två slagen av miljö vars sammanfattning inte anger
något resultat, ett bokkapitel vars författare metadata inte anger, och
en förhandspublikation vars författare bara anges med förnamn.  Studien
av brandväggsadministratörer citeras i sin första arXiv-version, som är
den vi läste; dess publicerade version har en annan titel.
\end{onlyin}

\section{\ensv{Conclusion}{Slutsats}} % sync english=89db2a swedish=c3cc6f
\label{sec:cli-efficiency-conclusion}

\begin{onlyin}{english} % sync english=fe9248 swedish=939c7e
No, not as a general rule: the evidence supports a weaker claim than
the one we tested.  For single tasks, some 45 experiments found no
consistent efficiency advantage for command interfaces over graphical
ones: several found no difference or a trade-off, several favoured
commands, and more favoured the graphical side, mostly for learners and
newcomers.  That
typed commands become faster than pointing with practice is supported
by the one experiment that measured it, for keystroke commands rather
than a shell.  For repeated or bulk tasks saved as scripts we found no
controlled experiment, only the consistent report of system
administrators that bulk changes need a command line; the strength of
command lines they named most often was flexibility and efficiency of
use, and many (64 of 309 coded remarks) named automation.  That
scripted work is easier to
reproduce is advice, resting on experience, beside evidence that manual
work in a spreadsheet corrupts data at scale and that most published
research scripts do not run again as they stand.  That command lines are
harder for novices holds for the start in most studies, not in all, and
not as a measured size.

For the deck this means that the main text may show what the reports
task shows, that a line of commands can be written down and run again
on any number of files while the clicks in a file manager are gone,
since that is a fact about the tools.  It may say that people who do
bulk work, such as system administrators, report that bulk changes need
a command line, and that guides for research computing advise scripts
instead of manual steps because manual steps are error-prone and hard
to redo.  It may not say that a command line is faster or more
efficient than a graphical interface, for experts or in general, nor
that a script makes work reproducible by itself.
\end{onlyin}
\begin{onlyin}{swedish}
Nej, inte som allmän regel: beläggen stöder ett svagare påstående än det
vi prövade.  För enskilda uppgifter fann ett 45-tal experiment ingen
genomgående fördel i effektivitet för kommandogränssnitt framför
grafiska: flera fann ingen skillnad eller en avvägning, flera gynnade
kommandon, och fler gynnade den grafiska sidan, mest för de som lär sig
och för nykomlingar.  Att skrivna kommandon med övning blir snabbare än
att peka stöds av det enda experiment som mätte det, för
tangentkommandon snarare än ett skal.  För upprepade uppgifter och
massuppgifter som sparas som skript hittade vi inget kontrollerat
experiment, bara systemadministratörernas samstämmiga rapport att
massändringar kräver en kommandorad; den styrka hos kommandorader som de
oftast nämnde var flexibilitet och effektivitet i användningen, och
många (64 av 309 kodade kommentarer) nämnde automatisering.  Att
skriptat arbete är lättare att upprepa är ett råd,
som vilar på erfarenhet, bredvid belägg för att manuellt arbete i ett
kalkylblad förvanskar data i stor skala och att de flesta publicerade
forskningsskript inte går att köra igen som de är.  Att kommandorader
är svårare för nybörjare gäller början i de flesta studier, inte i
alla, och inte som en uppmätt storlek.

För föreläsningen betyder det att huvudtexten får visa det som
rapportuppgiften visar, att en rad kommandon kan skrivas ner och köras
igen på hur många filer som helst medan klicken i en filhanterare är
borta, eftersom det är ett faktum om verktygen.  Den får säga att de
som gör massarbete, som systemadministratörer, rapporterar att
massändringar kräver en kommandorad, och att guider för
datoranvändning i forskning råder till skript i stället för manuella
steg eftersom manuella steg ger fel och är svåra att göra om.  Den får
inte säga att en kommandorad är snabbare eller effektivare än ett
grafiskt gränssnitt, för experter eller i allmänhet, och inte heller att
ett skript i sig gör arbete upprepbart.
\end{onlyin}

\section{\ensv{Further work}{Fortsatt arbete}} % sync english=a44494 swedish=969e69
\label{sec:cli-efficiency-further}

\begin{onlyin}{english} % sync english=2d5c82 swedish=7bb856
This search left work undone.  DBLP could not be searched, and the
older interface studies were read in abstract only; the full texts of
\textcite{Whiteside1985User}, \textcite{Antin1988Empirical},
\textcite{Durham1998Learning} and \textcite{Staggers2000Comparing}
would show the conditions under which each result held, and the full
text of \textcite{Remington2016With} how far its keystroke commands
resemble a shell.  The 26 studies mentioned in
\cref{sec:cli-efficiency-discussion} could not be obtained at all; with
some two dozen comparisons of the same kind unread, the balance of the
efficiency results could still shift.  We
searched OpenAlex in titles and abstracts only, and did not follow
citations from the studies we read, which in a field whose
classic studies have few abstracts would likely find more.

The field lacks the study that would settle the claim as the deck
makes it: a controlled comparison of the same bulk task, such as
selecting and joining hundreds of files, done in a file manager and at
a command line, by novices after a short introduction and by
practised users, measuring time, errors and whether the work can be
redone a week later.  A study of reproducibility would compare the same
analysis done by script and by hand in a graphical program, redone by
someone else.  If such studies favoured the command line, the deck
could state its advantage for bulk work as measured; if they did not,
the reports task would remain an illustration of what a command line
makes possible, not of what it makes faster.
\end{onlyin}
\begin{onlyin}{swedish}
Sökningen lämnade arbete ogjort.  DBLP kunde inte genomsökas, och de
äldre gränssnittsstudierna lästes bara i sammanfattning; fulltexterna
av \textcite{Whiteside1985User}, \textcite{Antin1988Empirical},
\textcite{Durham1998Learning} och \textcite{Staggers2000Comparing}
skulle visa under vilka villkor varje resultat gällde, och fulltexten av
\textcite{Remington2016With} hur mycket dess tangentkommandon liknar ett
skal.  De 26 studier som nämns i \cref{sec:cli-efficiency-discussion}
gick inte att få tag på alls; med ett par dussin jämförelser av samma
slag olästa kan balansen i effektivitetsresultaten fortfarande skifta.
Vi sökte i OpenAlex bara i titlar och sammanfattningar, och följde inte
citeringar från de studier vi läste, vilket i ett fält vars klassiska
studier har få sammanfattningar troligen skulle hitta fler.

Fältet saknar den studie som skulle avgöra påståendet så som
föreläsningen gör det: en kontrollerad jämförelse av samma massuppgift,
som att välja ut och foga samman hundratals filer, gjord i en
filhanterare och på en kommandorad, av nybörjare efter en kort
introduktion och av vana användare, som mäter tid, fel och om arbetet
kan göras om en vecka senare.  En studie av upprepbarhet skulle jämföra
samma analys gjord med skript och för hand i ett grafiskt program, och
omgjord av någon annan.  Gynnade sådana studier kommandoraden kunde
föreläsningen ange dess fördel för massarbete som uppmätt; gjorde de det
inte skulle rapportuppgiften förbli en illustration av vad en
kommandorad gör möjligt, inte av vad den gör snabbare.
\end{onlyin}

\section{\ensv{Hits bearing on the claim}{Träffar som rör påståendet}} % sync english=b678a6 swedish=eaf510
\label{sec:cli-efficiency-hits}

% Author's note (verification and reproduction):
%   scholar session : cli-efficiency
%   exports         : litteratursokning/cli-efficiency.csv (the session,
%                     every record and its decision; the abstract column
%                     is dropped for size, and every record keeps its DOI
%                     and URL, so its abstract can be looked up) and the
%                     -full hit-list tables
%   searched        : 2026-11-18; providers openalex wos ieee scopus
%                     (dblp: failed, anti-bot page HTTP 202 text/html, every
%                     query; s2: searched with the key renewed on
%                     2026-11-18, but it takes no Boolean operators, so S1's
%                     0 is no real negative and s2 was not used for the
%                     other queries; s2 and Crossref filled in abstracts);
%                     openalex via the backing-claims skill's
%                     scripts/openalex_title_abstract.py (S, C) and a
%                     title.search variant of it (T); wos via TS=/TI=,
%                     scopus via TITLE-ABS-KEY()/TITLE(), ieee free text /
%                     "Document Title": fields; the queries are in the
%                     Method section and tab:cli-efficiency-queries
%   semantic check  : the backing-claims skill's scripts/semantic_check.py,
%                     7 questions before the repair, 3 after
%   classification  : scholar llm context + llm classify --no-examples
%                     --no-enrich -n 100, model anthropic/claude-sonnet-5-5;
%                     272 human decisions in all via scholar sessions
%                     decide (see Screening); read: titles of all kept
%                     supports/qualifies rows and of the discarded rows
%                     with confidence < 0.7, abstracts where relevant, and
%                     the abstracts of the 154 works of the corroboration
%                     round
%   no-abstract rule: abstracts looked up by DOI in S2, then Crossref; 464
%                     rescued and reclassified, 1483 tagged no-abstract
%   table caption   : hand-edited after export in both languages (scholar
%                     calls uncategorised discards "duplicates or other
%                     versions")
%   hit-list tables : scholar sessions export cli-efficiency -f table
%                     --bearing-only --lang en|sv --label sok-cli-efficiency
%                     --track "Claim B"|"Påstående B" --theme ... -o
%                     litteratursokning/cli-efficiency-full(-sv)
%   provenance      : ltnotes.bib (CLAIM/FOUND-VIA/QUOTE/VERIFIED)
\begin{onlyin}{english} % sync english=609530 swedish=e63cc8
% --- scholar LaTeX fragment --------------------------------------------
% NOTE (hand edit after export, 2026-11-18): scholar labels every
% uncategorised discard "duplicates or other versions of cited sources";
% the 1252 are mostly the no-abstract exclusions, so the caption was
% reworded by hand.
% Generated by: scholar sessions export -f table
% Session: cli-efficiency
% \input this file from the body of a document whose
% preamble loads:
%   \usepackage{longtable}
%   \usepackage{booktabs}
%   \usepackage{array}
% ----------------------------------------------------------------------

\begingroup\footnotesize
\begin{longtable}{@{}%
  >{\raggedright\arraybackslash}p{0.44\textwidth}c%
  >{\raggedright\arraybackslash}p{0.13\textwidth}%
  >{\raggedright\arraybackslash}p{0.26\textwidth}@{}}
\caption{Records bearing on the claim, Claim B (121 cited, 1 supporting and 2 qualifying or contradicting it, out of 9744 unique hits; the 2400 adjacent and 5965 false hits, 1252 records excluded without a category (those without an abstract and the duplicates of cited sources), 3 excluded although classified as bearing on the claim are not listed). The last column gives the reason for each inclusion; for machine-classified rows the model's confidence is given. Databases: OpenAlex (OA; not open access), IEEE Xplore (IEEE), Web of Science (WoS), Scopus, Semantic Scholar (S2).}\label{tab:sok-cli-efficiency}\\
\toprule
Title & Year & Provider & Reason \\
\midrule
\endfirsthead
\multicolumn{4}{@{}l}{\emph{\tablename~\thetable{} (cont.)}}\\
\toprule
Title & Year & Provider & Reason \\
\midrule
\endhead
\midrule \multicolumn{4}{r@{}}{\emph{continued on next page}}\\
\endfoot
\bottomrule
\endlastfoot
A command line interface versus a graphical user interface in coding VR systems & 2007 & WoS, OA & \emph{cited}: experiments comparing command and graphical interfaces \\
A comparison of command, menu-selection and natural-language computer programs & 1983 & OA & \emph{cited}: experiments comparing command and graphical interfaces \\
A large-scale study on research code quality and execution & 2022 & OA & \emph{cited}: errors and re-execution in research data and code \\
A recipe for extremely reproducible enrichment analysis v2 & 2023 & OA & \emph{cited}: best-practice advice for research computing \\
A Reproducible MEEG Data Analysis Workflow with conda, Snakemake, and R Markdown & 2021 & OA & \emph{cited}: best-practice advice for research computing \\
A Study of File Manipulation by Novices Using Commands vs. Direct Manipulation & 1995 & OA & \emph{cited}: experiments comparing command and graphical interfaces \\
A study of physicians interaction with text-based and graphical-based electronic patient record systems & 2002 & WoS, OA & \emph{cited}: experiments comparing command and graphical interfaces \\
A Study of Scientific Computational Notebook Quality & 2026 & OA & \emph{cited}: errors and re-execution in research data and code \\
AN EMPIRICAL-COMPARISON OF MENU SELECTION, COMMAND ENTRY AND COMBINED MODES OF COMPUTER CONTROL & 1988 & WoS, OA & \emph{cited}: experiments comparing command and graphical interfaces \\
AN EMPIRICAL-COMPARISON OF MENU-SELECTION (CUI) AND DESK-TOP (GUI) COMPUTER-PROGRAMS CARRIED OUT BY BEGINNERS AND EXPERTS & 1992 & WoS, OA & \emph{cited}: experiments comparing command and graphical interfaces \\
AN EVALUATION OF COMMAND LINE AND MENU INTERFACES IN A CAD ENVIRONMENT & 1992 & WoS, OA & \emph{cited}: experiments comparing command and graphical interfaces \\
An Evaluation of Menu Systems for Space Station Interfaces & 1986 & OA & \emph{cited}: experiments comparing command and graphical interfaces \\
An experimental investigation of the roles of the computer interface and individual characteristics in the learning of computer systems & 1992 & OA & \emph{cited}: experiments comparing command and graphical interfaces \\
AN EXPERIMENTAL STUDY OF USER SATISFACTION WITH MULTIPLE DIALOG MODES & 1987 & OA & \emph{cited}: experiments comparing command and graphical interfaces \\
An Investigation of Students' Thought Processes in Solving Business Problems with a Database Application. & 2001 & OA & \emph{cited}: experiments comparing command and graphical interfaces \\
Articulating accounting database queries: An analysis of actual and perceived effort & 1996 & WoS & \emph{cited}: experiments comparing command and graphical interfaces \\
Assessing data availability and research reproducibility in hydrology and water resources & 2019 & OA & \emph{cited}: errors and re-execution in research data and code \\
Assessing the Impact of User Interface Abstraction on Online Telecommunications Course Laboratories & 2018 & IEEE & \emph{cited}: novices and command lines \\
Best Practices for Reproducible Microbial Genomics Analyses & 2025 & OA & \emph{cited}: best-practice advice for research computing \\
Best Practices for Scientific Computing & 2014 & OA & \emph{cited}: best-practice advice for research computing \\
Building a user-derived interface & 1984 & OA & \emph{cited}: novices and command lines \\
Ciencia reproducible: ¿qué, por qué, cómo? & 2016 & OA & \emph{cited}: best-practice advice for research computing \\
Claims for Direct Manipulation Interfaces Investigated & 1992 & OA & \emph{cited}: experiments comparing command and graphical interfaces \\
Cognitive constraints on computer problem-solving skills & 1997 & WoS, Scopus & \emph{cited}: novices and command lines \\
Command-language versus graphic-user interface types: Examining their associations with user performance and perceived system control & 1998 & WoS & \emph{cited}: experiments comparing command and graphical interfaces \\
Communicating With Graphic User Interfaces: A Comparison of Menu Selection and Menu Bypass Techniques & 1991 & OA & \emph{cited}: typed commands in graphical programs \\
Comparative analysis of iconic representation of select word processing functions with command string counterparts as a function of experience & 1988 & OA & \emph{cited}: typed commands in graphical programs \\
Comparing feature sets within visual and command line environments and their effect on novice programming (abstract only) & 2012 & OA & \emph{cited}: novices and command lines \\
Comparing Mental Models of Novice Programmers when Using Visual and Command Line Environments & 2012 & WoS, OA & \emph{cited}: novices and command lines \\
Comparing response time, errors, and satisfaction between text-based and graphical user interfaces during nursing order tasks & 2000 & WoS, OA & \emph{cited}: experiments comparing command and graphical interfaces \\
Comparing Text-based and Graphic User Interfaces for novice and expert users. & 2007 & OA & \emph{cited}: experiments comparing command and graphical interfaces \\
COMPARISON OF COMMAND-LANGUAGE AND DIRECT MANIPULATION INTERFACES FOR CNC MILLING MACHINES & 1992 & WoS, OA & \emph{cited}: experiments comparing command and graphical interfaces \\
Computational Reproducibility of R Code Supplements on OSF & 2025 & OA & \emph{cited}: errors and re-execution in research data and code \\
Computer Programming Skills for Environmental Sciences & 2012 & OA & \emph{cited}: best-practice advice for research computing \\
Data Analysis in Excel and R: A Comparative Evaluation & 2025 & OA & \emph{cited}: errors and re-execution in research data and code \\
Design and evaluation of CNC-user interfaces & 1992 & OA & \emph{cited}: experiments comparing command and graphical interfaces \\
Diagrammatic A Experiment & 1995 & OA & \emph{cited}: experiments comparing command and graphical interfaces \\
Dips and Ceilings: Understanding and Supporting Transitions to Expertise in User Interfaces & 2011 & WoS & \emph{cited}: typed commands in graphical programs \\
Dynamic software design for clinical exome and genome analyses: insights from bioinformaticians, clinical geneticists, and genetic counselors & 2015 & OA & \emph{cited}: practitioners' reports on bulk work \\
Effects of computer interfaces on computer-based statistical analysis & 1998 & WoS, OA & \emph{cited}: experiments comparing command and graphical interfaces \\
Empowering Interfaces for System Administrators: Keeping the Command Line in Mind when Designing GUIs & 2014 & OA & \emph{cited}: practitioners' reports on bulk work \\
Enabling the Verification of Computational Results: An Empirical Evaluation of Computational Reproducibility & 2018 & OA & \emph{cited}: errors and re-execution in research data and code \\
End-users, mediated searches, and front-end assistance programs on Dialog: A comparison of learning, performance, and satisfaction & 1990 & OA & \emph{cited}: experiments comparing command and graphical interfaces \\
Engineering Data Processing Workflows & 2024 & IEEE & \emph{cited}: best-practice advice for research computing \\
Evaluating Reproducibility in Computational Biology Research & 2018 & OA & \emph{cited}: errors and re-execution in research data and code \\
Evaluating User Interface Complexity & 1987 & OA & \emph{cited}: experiments comparing command and graphical interfaces \\
Evaluation of a command-line parser-based order entry pathway for the Department of Veterans Affairs electronic patient record & 2001 & WoS, OA & \emph{cited}: experiments comparing command and graphical interfaces \\
Evaluation of two intelligent tutoring system authoring tool paradigms graphical user interface-based and text-based & 2012 & Scopus, OA & \emph{cited}: experiments comparing command and graphical interfaces \\
Exploring the Impact of Exposing CS Majors to Programming Concepts using IDE Programming vs. non-IDE Programming in the Classroom & 2020 & OA & \emph{cited}: novices and command lines \\
Factors which contribute to ease of use of decision support systems and the impact of ease of use on decision effectiveness & 1992 & OA & \emph{cited}: experiments comparing command and graphical interfaces \\
First impressions of operating system styles affect usability & 2004 & OA & \emph{cited}: experiments comparing command and graphical interfaces \\
Functional R code is rare in species distribution and abundance papers & 2024 & OA & \emph{cited}: errors and re-execution in research data and code \\
Gene name errors are widespread in the scientific literature & 2016 & OA & \emph{cited}: errors and re-execution in research data and code \\
Gene name errors: Lessons not learned & 2021 & OA & \emph{cited}: errors and re-execution in research data and code \\
Getting the Details Right: Gene Signatures for Cancer Therapy & 2010 & OA & \emph{cited}: errors and re-execution in research data and code \\
Good enough practices in scientific computing & 2017 & OA & \emph{cited}: best-practice advice for research computing \\
GRAPHICAL VERSUS CHARACTER-BASED WORD-PROCESSORS - AN ANALYSIS OF USER PERFORMANCE & 1995 & WoS, OA & \emph{cited}: experiments comparing command and graphical interfaces \\
Graphical vs. Command User Interface: An Exploratory Study of Undergraduate Students' OPAC Interface Preferences & 2011 & OA & \emph{cited}: experiments comparing command and graphical interfaces \\
GUI vs CUI: Individual personality types and the experience of learning to use library databases & 1994 & OA & \emph{cited}: novices and command lines \\
Guided systems versus interactive online help & 1997 & IEEE, Scopus & \emph{cited}: experiments comparing command and graphical interfaces \\
Hidden costs of graphical user interfaces:: Failure to make the transition from menus and icon toolbars to keyboard shortcuts & 2005 & WoS, OA & \emph{cited}: typed commands in graphical programs \\
Human-in-the-Loop Design Automation: A Comparative Study of Graphical and Command-Line Interfaces in CAD and EDA Workflows & 2026 & OA & \emph{cited}: practitioners' reports on bulk work \\
Improving the Rigor and Reproducibility of Flow CytometryBased Clinical Research and Trials Through Automated Data Analysis & 2019 & OA & \emph{cited}: practitioners' reports on bulk work \\
Individual differences in the use of command line and menu computer interfaces & 1997 & WoS, OA & \emph{cited}: experiments comparing command and graphical interfaces \\
Infrastructures and Practices for Reproducible Research in Geography, Geosciences, and GIScience & 2022 & OA & \emph{cited}: errors and re-execution in research data and code \\
Insights from the SwissRN Computational Reproducibility Hackathon 2025 & 2026 & Scopus & \emph{cited}: errors and re-execution in research data and code \\
Interaction styles in tools for developing virtual environments & 2008 & OA & \emph{cited}: experiments comparing command and graphical interfaces \\
Is there really trouble with UNIX? & 1983 & OA & \emph{cited}: novices and command lines \\
Keyboard Shortcut Users: They Are Faster at More than Just Typing & 2009 & OA & \emph{cited}: typed commands in graphical programs \\
Learning and retention with a menu and a command line interface & 1998 & WoS, OA & \emph{cited}: experiments comparing command and graphical interfaces \\
Misconceived misconceptions? & 1988 & OA & \emph{cited}: novices and command lines \\
Most computational hydrology is not reproducible, so is it really science? & 2016 & OA & \emph{cited}: best-practice advice for research computing \\
NATURAL COMMAND NAMES AND INITIAL LEARNING - A STUDY OF TEXT-EDITING TERMS & 1983 & WoS, OA & \emph{cited}: novices and command lines \\
Open Science in Software Engineering: A Study on Deep Learning-Based Vulnerability Detection & 2023 & IEEE & \emph{cited}: errors and re-execution in research data and code \\
opentsv prevents the corruption of scientific data by Excel & 2018 & OA & \emph{cited}: errors and re-execution in research data and code \\
Overcoming barriers to reproducibility in geoscientific data analysis: Challenges and practical implementation strategies & 2026 & OA & \emph{cited}: best-practice advice for research computing \\
Promoting and characterizing the menu to keyboard shortcuts transition & 2019 & OA & \emph{cited}: typed commands in graphical programs \\
PROMPT COMPREHENSION IN UNIX COMMAND PRODUCTION & 1992 & WoS, OA & \emph{cited}: novices and command lines \\
ReproduceMe: Lessons from a pilot project on computational reproducibility & 2024 & OA & \emph{cited}: best-practice advice for research computing \\
Reproducibility and Transparency by Design & 2019 & OA & \emph{cited}: best-practice advice for research computing \\
Reproducibility is a Process, Not an Achievement: The Replicability of IR Reproducibility Experiments & 2020 & OA & \emph{cited}: errors and re-execution in research data and code \\
Reproducible Data Analysis With R in Laboratory Hematology & 2026 & OA & \emph{cited}: best-practice advice for research computing \\
Reproducible Research and Reports with R & 2019 & OA & \emph{cited}: best-practice advice for research computing \\
Research computing desktops: Demystifying research computing for non-Linux users & 2019 & Scopus & \emph{cited}: novices and command lines \\
Satisficing and the Use of Keyboard Shortcuts: Being Good Enough Is Enough? & 2013 & IEEE & \emph{cited}: typed commands in graphical programs \\
Scripts and reproducible workflows & 2026 & OA & \emph{cited}: best-practice advice for research computing \\
Statistical Considerations for Research Reproducibility: Strategies for Building Transparent Analytic Workflows & 2026 & OA & \emph{cited}: best-practice advice for research computing \\
Studying the Novice's Perception of Visual Vs. Command Line Programming Tools in CS1 & 2012 & OA & \emph{cited}: novices and command lines \\
System Administrators Are Parting Ways With Command Line Interfaces: An Exploratory Study of Firewall Interfaces. & 2019 & OA & \emph{cited}: practitioners' reports on bulk work \\
Ten simple rules for managing high-throughput nucleotide sequencing data & 2016 & OA & \emph{cited}: best-practice advice for research computing \\
Ten Simple Rules for Reproducible Computational Research & 2013 & OA & \emph{cited}: best-practice advice for research computing \\
The effect of interaction style and training method on end user learning of software packages & 1998 & WoS, Scopus & \emph{cited}: novices and command lines \\
The effect of task complexity on user interfaces: a comparison of command language interface and direct manipulation interfaces. & 1991 & OA & \emph{cited}: experiments comparing command and graphical interfaces \\
The Effect of the Interface on Learning Style in a Simulation-Based Learning Situation & 1997 & OA & \emph{cited}: experiments comparing command and graphical interfaces \\
The effects of an intrinsically motivating instructional environment on software learning and acceptance & 1994 & OA & \emph{cited}: novices and command lines \\
The Effects of Manipulating Recognition Variables on the Transition from Menus to Commands & 1991 & OA & \emph{cited}: novices and command lines \\
The effects of mice and pulldown menus versus commanddriven interfaces on writing & 1995 & OA & \emph{cited}: experiments comparing command and graphical interfaces \\
THE EFFECTS OF THE AVAILABILITY OF MENU INFORMATION DURING COMMAND LEARNING IN A WORD-PROCESSING APPLICATION & 1989 & WoS, OA & \emph{cited}: typed commands in graphical programs \\
The effects of using icons and direct manipulation interfaces : an empirical study & 2008 & OA & \emph{cited}: experiments comparing command and graphical interfaces \\
The Influence of Computer Experience on the Agreement between User Preference and Performance Rankings in Software Usability Testing & 1997 & OA & \emph{cited}: experiments comparing command and graphical interfaces \\
The influence of interaction style and experience on user perceptions of software packages & 1997 & WoS, OA & \emph{cited}: experiments comparing command and graphical interfaces \\
The mediating effects of intrinsic motivation, ease of use and usefulness perceptions on performance in first-time and subsequent computer users & 2001 & WoS, OA & \emph{cited}: novices and command lines \\
The State of Play of Reproducibility in Statistics: An Empirical Analysis & 2022 & OA & \emph{cited}: errors and re-execution in research data and code \\
The trouble with shortcuts & 2000 & OA & \emph{cited}: typed commands in graphical programs \\
Thinking or acting? Mental workload and subjective preferences for a command code and a direct manipulation interaction style & 1992 & OA & \emph{cited}: experiments comparing command and graphical interfaces \\
Toolglasses, marking menus, and hotkeys: A comparison of one and two-handed command selection techniques & 2004 & WoS, OA & \emph{cited}: typed commands in graphical programs \\
Towards a more reproducible ecology & 2016 & OA & \emph{cited}: best-practice advice for research computing \\
TRAINING END USERS - AN EXPERIMENTAL INVESTIGATION OF THE ROLES OF THE COMPUTER-INTERFACE AND TRAINING METHODS & 1993 & WoS, OA & \emph{cited}: experiments comparing command and graphical interfaces \\
Training Wheels for the Command Line & 2005 & OA & \emph{cited}: novices and command lines \\
Trust as an underlying factor of system administrator interface choice & 2006 & OA & \emph{cited}: practitioners' reports on bulk work \\
UNDERSTANDING SYSTEM ADMINISTRATORS' WORK PRACTICES AND THE ROLE FOR ENHANCED VISUALIZATIONS IN THEIR TOOLS & 2014 & OA & \emph{cited}: practitioners' reports on bulk work \\
Understanding the Executability of Example Code in Hugging Face Model Cards: An Empirical Study & 2026 & Scopus & \emph{cited}: errors and re-execution in research data and code \\
Unix Literacy for First-Year Computer Science Students & 2016 & OA & \emph{cited}: novices and command lines \\
User performance with command, menu, and iconic interfaces & 1985 & OA & \emph{cited}: experiments comparing command and graphical interfaces \\
User-friendly workflows for catchment modelling: Towards reproducible SWAT+ model studies & 2020 & OA & \emph{cited}: errors and re-execution in research data and code \\
Using menus to learn & 1989 & Scopus, OA & \emph{cited}: novices and command lines \\
Using R for Data Analysis: A Best Practice for Research & 2008 & OA & \emph{cited}: best-practice advice for research computing \\
Warming up to computers: A study of cognitive and affective interaction over time & 1982 & OA & \emph{cited}: experiments comparing command and graphical interfaces \\
When code does not run: reproducibility challenges in materials machine learning benchmarks & 2026 & OA & \emph{cited}: errors and re-execution in research data and code \\
Why teach unix? & 2007 & Scopus, OA & \emph{cited}: novices and command lines \\
With practice, keyboard shortcuts become faster than menu selection: A crossover interaction. & 2015 & OA & \emph{cited}: typed commands in graphical programs \\
Utilizing the keyboard to issue commands: Do self-efficacy, goal orientation, and computer anxiety impact usage of efficient techniques with software? & 2016 & OA & supports the claim (0.70) \\
Growth in a Time of Debt, Replicated: A Transparent Re-examination of the Reinhart-Rogoff 90-Percent Debt Threshold & 2026 & OA & qualifies or contradicts the claim (0.40) \\
WHEN DOES A MOUSE BECOME A RAT OR ... COMPARING PERFORMANCE AND PREFERENCES IN DIRECT MANIPULATION AND COMMAND LINE ENVIRONMENT & 1991 & WoS, OA & qualifies or contradicts the claim (0.75) \\
\end{longtable}
\endgroup

\end{onlyin}
\begin{onlyin}{swedish}
% --- scholar LaTeX fragment --------------------------------------------
% NOTE (hand edit after export, 2026-11-18): scholar labels every
% uncategorised discard "dubbletter eller andra versioner av citerade
% källor"; the 1252 are mostly the no-abstract exclusions, so the caption
% was reworded by hand, as in the English table.
% Generated by: scholar sessions export -f table
% Session: cli-efficiency
% \input this file from the body of a document whose
% preamble loads:
%   \usepackage{longtable}
%   \usepackage{booktabs}
%   \usepackage{array}
% ----------------------------------------------------------------------

\begingroup\footnotesize
\begin{longtable}{@{}%
  >{\raggedright\arraybackslash}p{0.44\textwidth}c%
  >{\raggedright\arraybackslash}p{0.13\textwidth}%
  >{\raggedright\arraybackslash}p{0.26\textwidth}@{}}
\caption{Träffar som rör påståendet, Påstående B (121 citerade, 1 som stöder påståendet och 2 som kvalificerar eller motsäger det, av 9744 unika träffar; de 2400 angränsande och 5965 felträffarna, 1252 poster uteslutna utan kategori (de som saknar sammanfattning och dubbletterna av citerade källor), 3 uteslutna trots att de klassats som relevanta för påståendet listas inte). Skälet till varje inkludering står i sista kolumnen; för maskinklassade rader anges modellens konfidens. Databaser: OpenAlex (OA; inte open access), IEEE Xplore (IEEE), Web of Science (WoS), Scopus, Semantic Scholar (S2).}\label{tab:sok-cli-efficiency}\\
\toprule
Titel & År & Leverantör & Skäl \\
\midrule
\endfirsthead
\multicolumn{4}{@{}l}{\emph{\tablename~\thetable{} (forts.)}}\\
\toprule
Titel & År & Leverantör & Skäl \\
\midrule
\endhead
\midrule \multicolumn{4}{r@{}}{\emph{forts.\ på nästa sida}}\\
\endfoot
\bottomrule
\endlastfoot
A command line interface versus a graphical user interface in coding VR systems & 2007 & WoS, OA & \emph{citerad}: experiment som jämför kommando- och grafiska gränssnitt \\
A comparison of command, menu-selection and natural-language computer programs & 1983 & OA & \emph{citerad}: experiment som jämför kommando- och grafiska gränssnitt \\
A large-scale study on research code quality and execution & 2022 & OA & \emph{citerad}: fel och omkörning i forskningsdata och kod \\
A recipe for extremely reproducible enrichment analysis v2 & 2023 & OA & \emph{citerad}: råd om datoranvändning i forskning \\
A Reproducible MEEG Data Analysis Workflow with conda, Snakemake, and R Markdown & 2021 & OA & \emph{citerad}: råd om datoranvändning i forskning \\
A Study of File Manipulation by Novices Using Commands vs. Direct Manipulation & 1995 & OA & \emph{citerad}: experiment som jämför kommando- och grafiska gränssnitt \\
A study of physicians interaction with text-based and graphical-based electronic patient record systems & 2002 & WoS, OA & \emph{citerad}: experiment som jämför kommando- och grafiska gränssnitt \\
A Study of Scientific Computational Notebook Quality & 2026 & OA & \emph{citerad}: fel och omkörning i forskningsdata och kod \\
AN EMPIRICAL-COMPARISON OF MENU SELECTION, COMMAND ENTRY AND COMBINED MODES OF COMPUTER CONTROL & 1988 & WoS, OA & \emph{citerad}: experiment som jämför kommando- och grafiska gränssnitt \\
AN EMPIRICAL-COMPARISON OF MENU-SELECTION (CUI) AND DESK-TOP (GUI) COMPUTER-PROGRAMS CARRIED OUT BY BEGINNERS AND EXPERTS & 1992 & WoS, OA & \emph{citerad}: experiment som jämför kommando- och grafiska gränssnitt \\
AN EVALUATION OF COMMAND LINE AND MENU INTERFACES IN A CAD ENVIRONMENT & 1992 & WoS, OA & \emph{citerad}: experiment som jämför kommando- och grafiska gränssnitt \\
An Evaluation of Menu Systems for Space Station Interfaces & 1986 & OA & \emph{citerad}: experiment som jämför kommando- och grafiska gränssnitt \\
An experimental investigation of the roles of the computer interface and individual characteristics in the learning of computer systems & 1992 & OA & \emph{citerad}: experiment som jämför kommando- och grafiska gränssnitt \\
AN EXPERIMENTAL STUDY OF USER SATISFACTION WITH MULTIPLE DIALOG MODES & 1987 & OA & \emph{citerad}: experiment som jämför kommando- och grafiska gränssnitt \\
An Investigation of Students' Thought Processes in Solving Business Problems with a Database Application. & 2001 & OA & \emph{citerad}: experiment som jämför kommando- och grafiska gränssnitt \\
Articulating accounting database queries: An analysis of actual and perceived effort & 1996 & WoS & \emph{citerad}: experiment som jämför kommando- och grafiska gränssnitt \\
Assessing data availability and research reproducibility in hydrology and water resources & 2019 & OA & \emph{citerad}: fel och omkörning i forskningsdata och kod \\
Assessing the Impact of User Interface Abstraction on Online Telecommunications Course Laboratories & 2018 & IEEE & \emph{citerad}: nybörjare och kommandorader \\
Best Practices for Reproducible Microbial Genomics Analyses & 2025 & OA & \emph{citerad}: råd om datoranvändning i forskning \\
Best Practices for Scientific Computing & 2014 & OA & \emph{citerad}: råd om datoranvändning i forskning \\
Building a user-derived interface & 1984 & OA & \emph{citerad}: nybörjare och kommandorader \\
Ciencia reproducible: ¿qué, por qué, cómo? & 2016 & OA & \emph{citerad}: råd om datoranvändning i forskning \\
Claims for Direct Manipulation Interfaces Investigated & 1992 & OA & \emph{citerad}: experiment som jämför kommando- och grafiska gränssnitt \\
Cognitive constraints on computer problem-solving skills & 1997 & WoS, Scopus & \emph{citerad}: nybörjare och kommandorader \\
Command-language versus graphic-user interface types: Examining their associations with user performance and perceived system control & 1998 & WoS & \emph{citerad}: experiment som jämför kommando- och grafiska gränssnitt \\
Communicating With Graphic User Interfaces: A Comparison of Menu Selection and Menu Bypass Techniques & 1991 & OA & \emph{citerad}: skrivna kommandon i grafiska program \\
Comparative analysis of iconic representation of select word processing functions with command string counterparts as a function of experience & 1988 & OA & \emph{citerad}: skrivna kommandon i grafiska program \\
Comparing feature sets within visual and command line environments and their effect on novice programming (abstract only) & 2012 & OA & \emph{citerad}: nybörjare och kommandorader \\
Comparing Mental Models of Novice Programmers when Using Visual and Command Line Environments & 2012 & WoS, OA & \emph{citerad}: nybörjare och kommandorader \\
Comparing response time, errors, and satisfaction between text-based and graphical user interfaces during nursing order tasks & 2000 & WoS, OA & \emph{citerad}: experiment som jämför kommando- och grafiska gränssnitt \\
Comparing Text-based and Graphic User Interfaces for novice and expert users. & 2007 & OA & \emph{citerad}: experiment som jämför kommando- och grafiska gränssnitt \\
COMPARISON OF COMMAND-LANGUAGE AND DIRECT MANIPULATION INTERFACES FOR CNC MILLING MACHINES & 1992 & WoS, OA & \emph{citerad}: experiment som jämför kommando- och grafiska gränssnitt \\
Computational Reproducibility of R Code Supplements on OSF & 2025 & OA & \emph{citerad}: fel och omkörning i forskningsdata och kod \\
Computer Programming Skills for Environmental Sciences & 2012 & OA & \emph{citerad}: råd om datoranvändning i forskning \\
Data Analysis in Excel and R: A Comparative Evaluation & 2025 & OA & \emph{citerad}: fel och omkörning i forskningsdata och kod \\
Design and evaluation of CNC-user interfaces & 1992 & OA & \emph{citerad}: experiment som jämför kommando- och grafiska gränssnitt \\
Diagrammatic A Experiment & 1995 & OA & \emph{citerad}: experiment som jämför kommando- och grafiska gränssnitt \\
Dips and Ceilings: Understanding and Supporting Transitions to Expertise in User Interfaces & 2011 & WoS & \emph{citerad}: skrivna kommandon i grafiska program \\
Dynamic software design for clinical exome and genome analyses: insights from bioinformaticians, clinical geneticists, and genetic counselors & 2015 & OA & \emph{citerad}: yrkesverksammas rapporter om massarbete \\
Effects of computer interfaces on computer-based statistical analysis & 1998 & WoS, OA & \emph{citerad}: experiment som jämför kommando- och grafiska gränssnitt \\
Empowering Interfaces for System Administrators: Keeping the Command Line in Mind when Designing GUIs & 2014 & OA & \emph{citerad}: yrkesverksammas rapporter om massarbete \\
Enabling the Verification of Computational Results: An Empirical Evaluation of Computational Reproducibility & 2018 & OA & \emph{citerad}: fel och omkörning i forskningsdata och kod \\
End-users, mediated searches, and front-end assistance programs on Dialog: A comparison of learning, performance, and satisfaction & 1990 & OA & \emph{citerad}: experiment som jämför kommando- och grafiska gränssnitt \\
Engineering Data Processing Workflows & 2024 & IEEE & \emph{citerad}: råd om datoranvändning i forskning \\
Evaluating Reproducibility in Computational Biology Research & 2018 & OA & \emph{citerad}: fel och omkörning i forskningsdata och kod \\
Evaluating User Interface Complexity & 1987 & OA & \emph{citerad}: experiment som jämför kommando- och grafiska gränssnitt \\
Evaluation of a command-line parser-based order entry pathway for the Department of Veterans Affairs electronic patient record & 2001 & WoS, OA & \emph{citerad}: experiment som jämför kommando- och grafiska gränssnitt \\
Evaluation of two intelligent tutoring system authoring tool paradigms graphical user interface-based and text-based & 2012 & Scopus, OA & \emph{citerad}: experiment som jämför kommando- och grafiska gränssnitt \\
Exploring the Impact of Exposing CS Majors to Programming Concepts using IDE Programming vs. non-IDE Programming in the Classroom & 2020 & OA & \emph{citerad}: nybörjare och kommandorader \\
Factors which contribute to ease of use of decision support systems and the impact of ease of use on decision effectiveness & 1992 & OA & \emph{citerad}: experiment som jämför kommando- och grafiska gränssnitt \\
First impressions of operating system styles affect usability & 2004 & OA & \emph{citerad}: experiment som jämför kommando- och grafiska gränssnitt \\
Functional R code is rare in species distribution and abundance papers & 2024 & OA & \emph{citerad}: fel och omkörning i forskningsdata och kod \\
Gene name errors are widespread in the scientific literature & 2016 & OA & \emph{citerad}: fel och omkörning i forskningsdata och kod \\
Gene name errors: Lessons not learned & 2021 & OA & \emph{citerad}: fel och omkörning i forskningsdata och kod \\
Getting the Details Right: Gene Signatures for Cancer Therapy & 2010 & OA & \emph{citerad}: fel och omkörning i forskningsdata och kod \\
Good enough practices in scientific computing & 2017 & OA & \emph{citerad}: råd om datoranvändning i forskning \\
GRAPHICAL VERSUS CHARACTER-BASED WORD-PROCESSORS - AN ANALYSIS OF USER PERFORMANCE & 1995 & WoS, OA & \emph{citerad}: experiment som jämför kommando- och grafiska gränssnitt \\
Graphical vs. Command User Interface: An Exploratory Study of Undergraduate Students' OPAC Interface Preferences & 2011 & OA & \emph{citerad}: experiment som jämför kommando- och grafiska gränssnitt \\
GUI vs CUI: Individual personality types and the experience of learning to use library databases & 1994 & OA & \emph{citerad}: nybörjare och kommandorader \\
Guided systems versus interactive online help & 1997 & IEEE, Scopus & \emph{citerad}: experiment som jämför kommando- och grafiska gränssnitt \\
Hidden costs of graphical user interfaces:: Failure to make the transition from menus and icon toolbars to keyboard shortcuts & 2005 & WoS, OA & \emph{citerad}: skrivna kommandon i grafiska program \\
Human-in-the-Loop Design Automation: A Comparative Study of Graphical and Command-Line Interfaces in CAD and EDA Workflows & 2026 & OA & \emph{citerad}: yrkesverksammas rapporter om massarbete \\
Improving the Rigor and Reproducibility of Flow CytometryBased Clinical Research and Trials Through Automated Data Analysis & 2019 & OA & \emph{citerad}: yrkesverksammas rapporter om massarbete \\
Individual differences in the use of command line and menu computer interfaces & 1997 & WoS, OA & \emph{citerad}: experiment som jämför kommando- och grafiska gränssnitt \\
Infrastructures and Practices for Reproducible Research in Geography, Geosciences, and GIScience & 2022 & OA & \emph{citerad}: fel och omkörning i forskningsdata och kod \\
Insights from the SwissRN Computational Reproducibility Hackathon 2025 & 2026 & Scopus & \emph{citerad}: fel och omkörning i forskningsdata och kod \\
Interaction styles in tools for developing virtual environments & 2008 & OA & \emph{citerad}: experiment som jämför kommando- och grafiska gränssnitt \\
Is there really trouble with UNIX? & 1983 & OA & \emph{citerad}: nybörjare och kommandorader \\
Keyboard Shortcut Users: They Are Faster at More than Just Typing & 2009 & OA & \emph{citerad}: skrivna kommandon i grafiska program \\
Learning and retention with a menu and a command line interface & 1998 & WoS, OA & \emph{citerad}: experiment som jämför kommando- och grafiska gränssnitt \\
Misconceived misconceptions? & 1988 & OA & \emph{citerad}: nybörjare och kommandorader \\
Most computational hydrology is not reproducible, so is it really science? & 2016 & OA & \emph{citerad}: råd om datoranvändning i forskning \\
NATURAL COMMAND NAMES AND INITIAL LEARNING - A STUDY OF TEXT-EDITING TERMS & 1983 & WoS, OA & \emph{citerad}: nybörjare och kommandorader \\
Open Science in Software Engineering: A Study on Deep Learning-Based Vulnerability Detection & 2023 & IEEE & \emph{citerad}: fel och omkörning i forskningsdata och kod \\
opentsv prevents the corruption of scientific data by Excel & 2018 & OA & \emph{citerad}: fel och omkörning i forskningsdata och kod \\
Overcoming barriers to reproducibility in geoscientific data analysis: Challenges and practical implementation strategies & 2026 & OA & \emph{citerad}: råd om datoranvändning i forskning \\
Promoting and characterizing the menu to keyboard shortcuts transition & 2019 & OA & \emph{citerad}: skrivna kommandon i grafiska program \\
PROMPT COMPREHENSION IN UNIX COMMAND PRODUCTION & 1992 & WoS, OA & \emph{citerad}: nybörjare och kommandorader \\
ReproduceMe: Lessons from a pilot project on computational reproducibility & 2024 & OA & \emph{citerad}: råd om datoranvändning i forskning \\
Reproducibility and Transparency by Design & 2019 & OA & \emph{citerad}: råd om datoranvändning i forskning \\
Reproducibility is a Process, Not an Achievement: The Replicability of IR Reproducibility Experiments & 2020 & OA & \emph{citerad}: fel och omkörning i forskningsdata och kod \\
Reproducible Data Analysis With R in Laboratory Hematology & 2026 & OA & \emph{citerad}: råd om datoranvändning i forskning \\
Reproducible Research and Reports with R & 2019 & OA & \emph{citerad}: råd om datoranvändning i forskning \\
Research computing desktops: Demystifying research computing for non-Linux users & 2019 & Scopus & \emph{citerad}: nybörjare och kommandorader \\
Satisficing and the Use of Keyboard Shortcuts: Being Good Enough Is Enough? & 2013 & IEEE & \emph{citerad}: skrivna kommandon i grafiska program \\
Scripts and reproducible workflows & 2026 & OA & \emph{citerad}: råd om datoranvändning i forskning \\
Statistical Considerations for Research Reproducibility: Strategies for Building Transparent Analytic Workflows & 2026 & OA & \emph{citerad}: råd om datoranvändning i forskning \\
Studying the Novice's Perception of Visual Vs. Command Line Programming Tools in CS1 & 2012 & OA & \emph{citerad}: nybörjare och kommandorader \\
System Administrators Are Parting Ways With Command Line Interfaces: An Exploratory Study of Firewall Interfaces. & 2019 & OA & \emph{citerad}: yrkesverksammas rapporter om massarbete \\
Ten simple rules for managing high-throughput nucleotide sequencing data & 2016 & OA & \emph{citerad}: råd om datoranvändning i forskning \\
Ten Simple Rules for Reproducible Computational Research & 2013 & OA & \emph{citerad}: råd om datoranvändning i forskning \\
The effect of interaction style and training method on end user learning of software packages & 1998 & WoS, Scopus & \emph{citerad}: nybörjare och kommandorader \\
The effect of task complexity on user interfaces: a comparison of command language interface and direct manipulation interfaces. & 1991 & OA & \emph{citerad}: experiment som jämför kommando- och grafiska gränssnitt \\
The Effect of the Interface on Learning Style in a Simulation-Based Learning Situation & 1997 & OA & \emph{citerad}: experiment som jämför kommando- och grafiska gränssnitt \\
The effects of an intrinsically motivating instructional environment on software learning and acceptance & 1994 & OA & \emph{citerad}: nybörjare och kommandorader \\
The Effects of Manipulating Recognition Variables on the Transition from Menus to Commands & 1991 & OA & \emph{citerad}: nybörjare och kommandorader \\
The effects of mice and pulldown menus versus commanddriven interfaces on writing & 1995 & OA & \emph{citerad}: experiment som jämför kommando- och grafiska gränssnitt \\
THE EFFECTS OF THE AVAILABILITY OF MENU INFORMATION DURING COMMAND LEARNING IN A WORD-PROCESSING APPLICATION & 1989 & WoS, OA & \emph{citerad}: skrivna kommandon i grafiska program \\
The effects of using icons and direct manipulation interfaces : an empirical study & 2008 & OA & \emph{citerad}: experiment som jämför kommando- och grafiska gränssnitt \\
The Influence of Computer Experience on the Agreement between User Preference and Performance Rankings in Software Usability Testing & 1997 & OA & \emph{citerad}: experiment som jämför kommando- och grafiska gränssnitt \\
The influence of interaction style and experience on user perceptions of software packages & 1997 & WoS, OA & \emph{citerad}: experiment som jämför kommando- och grafiska gränssnitt \\
The mediating effects of intrinsic motivation, ease of use and usefulness perceptions on performance in first-time and subsequent computer users & 2001 & WoS, OA & \emph{citerad}: nybörjare och kommandorader \\
The State of Play of Reproducibility in Statistics: An Empirical Analysis & 2022 & OA & \emph{citerad}: fel och omkörning i forskningsdata och kod \\
The trouble with shortcuts & 2000 & OA & \emph{citerad}: skrivna kommandon i grafiska program \\
Thinking or acting? Mental workload and subjective preferences for a command code and a direct manipulation interaction style & 1992 & OA & \emph{citerad}: experiment som jämför kommando- och grafiska gränssnitt \\
Toolglasses, marking menus, and hotkeys: A comparison of one and two-handed command selection techniques & 2004 & WoS, OA & \emph{citerad}: skrivna kommandon i grafiska program \\
Towards a more reproducible ecology & 2016 & OA & \emph{citerad}: råd om datoranvändning i forskning \\
TRAINING END USERS - AN EXPERIMENTAL INVESTIGATION OF THE ROLES OF THE COMPUTER-INTERFACE AND TRAINING METHODS & 1993 & WoS, OA & \emph{citerad}: experiment som jämför kommando- och grafiska gränssnitt \\
Training Wheels for the Command Line & 2005 & OA & \emph{citerad}: nybörjare och kommandorader \\
Trust as an underlying factor of system administrator interface choice & 2006 & OA & \emph{citerad}: yrkesverksammas rapporter om massarbete \\
UNDERSTANDING SYSTEM ADMINISTRATORS' WORK PRACTICES AND THE ROLE FOR ENHANCED VISUALIZATIONS IN THEIR TOOLS & 2014 & OA & \emph{citerad}: yrkesverksammas rapporter om massarbete \\
Understanding the Executability of Example Code in Hugging Face Model Cards: An Empirical Study & 2026 & Scopus & \emph{citerad}: fel och omkörning i forskningsdata och kod \\
Unix Literacy for First-Year Computer Science Students & 2016 & OA & \emph{citerad}: nybörjare och kommandorader \\
User performance with command, menu, and iconic interfaces & 1985 & OA & \emph{citerad}: experiment som jämför kommando- och grafiska gränssnitt \\
User-friendly workflows for catchment modelling: Towards reproducible SWAT+ model studies & 2020 & OA & \emph{citerad}: fel och omkörning i forskningsdata och kod \\
Using menus to learn & 1989 & Scopus, OA & \emph{citerad}: nybörjare och kommandorader \\
Using R for Data Analysis: A Best Practice for Research & 2008 & OA & \emph{citerad}: råd om datoranvändning i forskning \\
Warming up to computers: A study of cognitive and affective interaction over time & 1982 & OA & \emph{citerad}: experiment som jämför kommando- och grafiska gränssnitt \\
When code does not run: reproducibility challenges in materials machine learning benchmarks & 2026 & OA & \emph{citerad}: fel och omkörning i forskningsdata och kod \\
Why teach unix? & 2007 & Scopus, OA & \emph{citerad}: nybörjare och kommandorader \\
With practice, keyboard shortcuts become faster than menu selection: A crossover interaction. & 2015 & OA & \emph{citerad}: skrivna kommandon i grafiska program \\
Utilizing the keyboard to issue commands: Do self-efficacy, goal orientation, and computer anxiety impact usage of efficient techniques with software? & 2016 & OA & stöder påståendet (0.70) \\
Growth in a Time of Debt, Replicated: A Transparent Re-examination of the Reinhart-Rogoff 90-Percent Debt Threshold & 2026 & OA & kvalificerar eller motsäger påståendet (0.40) \\
WHEN DOES A MOUSE BECOME A RAT OR ... COMPARING PERFORMANCE AND PREFERENCES IN DIRECT MANIPULATION AND COMMAND LINE ENVIRONMENT & 1991 & WoS, OA & kvalificerar eller motsäger påståendet (0.75) \\
\end{longtable}
\endgroup

\end{onlyin}

% --- Appendix C: claim C.  A stub until the chapter (hpc-cli) is final.
% TODO(claim C): the chapter and its bounds for sec:where-cli.
\chapter{\ensv{Are computing clusters reached through a command line?}%
  {Nås beräkningskluster genom en kommandorad?}} % sync english=91c6fc swedish=864ccd
\label{app:hpc-cli}

\chapterprecis{\ensv{The author has not yet reviewed the results in this
  appendix in full.}%
  {Författaren har ännu inte granskat resultaten i den här bilagan i sin
  helhet.}} % sync english=21ea16 swedish=60ed7a

\begin{onlyin}{english} % sync english=a3bd8c swedish=d35af1
This appendix is being written.  It will report how the computing
clusters used in research are reached: through a command line, a
graphical remote desktop or a web portal.
\end{onlyin}
\begin{onlyin}{swedish}
Den här bilagan skrivs just nu.  Den kommer att redovisa hur de
beräkningskluster som används i forskningen nås: genom en kommandorad,
ett grafiskt fjärrskrivbord eller en webbportal.
\end{onlyin}

% --- Appendix D: the Unix lineage, from the search of 2026-11-18
% (scholar session unix-lineage).
\chapter{\ensv{Which of today's systems descend from Unix?}%
  {Vilka av dagens system härstammar från Unix?}} % sync english=9c8ed7 swedish=dff1b1
\label{app:unix-lineage}

\chapterprecis{\ensv{The author has not yet reviewed the results in this
  appendix in full.}%
  {Författaren har ännu inte granskat resultaten i den här bilagan i sin
  helhet.}} % sync english=21ea16 swedish=60ed7a

\begin{onlyin}{english} % sync english=9ad102 swedish=022606
In \cref{sec:where-cli} we say that the command line you meet on Linux
is much alike on macOS and the BSD systems, because the BSD systems and
macOS descend from Unix while Linux imitates it, and that Android is
built on the Linux kernel.
Unix is the operating system that Ritchie and Thompson described at Bell
Laboratories\autocite{RitchieThompson1974}.
This appendix asks which of today's systems descend from it, and it
separates three relations that everyday speech runs together.
A system \emph{descends} from Unix if its code has a history that leads
back to the Bell Labs system.
A system is \emph{Unix-like} if it imitates Unix's interfaces and
behaviour without having its code.
A system is a \emph{certified UNIX} if The Open Group has registered it
and allows it to carry the trademark.
The claim we test is the following.
\begin{quote}
  Linux systems are Unix-like: they imitate Unix without descending
  from its code.
  The BSD systems descend from Unix.
  macOS descends from Unix through BSD and is a certified UNIX.
  Android is built on the Linux kernel.
\end{quote}
The question is therefore about the status of the \emph{fact}: it asks
for evidence of descent, of imitation and of certification, one system
at a time, and the answer is a different relation for each.
\end{onlyin}
\begin{onlyin}{swedish}
I \cref{sec:where-cli} säger vi att den kommandorad som du möter på
Linux är mycket lik den på macOS och BSD-systemen, eftersom BSD-systemen
och macOS härstammar från Unix medan Linux efterliknar det, och att
Android är byggt på Linuxkärnan.
Unix är det operativsystem som Ritchie och Thompson beskrev vid Bell
Laboratories\autocite{RitchieThompson1974}.
Den här bilagan frågar vilka av dagens system som härstammar från det,
och den skiljer på tre förhållanden som vardagligt tal blandar ihop.
Ett system \emph{härstammar} från Unix om dess kod har en historia som
leder tillbaka till systemet från Bell Labs.
Ett system är \emph{Unixliknande} om det efterliknar Unix gränssnitt och
beteende utan att ha dess kod.
Ett system är ett \emph{certifierat UNIX} om The Open Group har
registrerat det och tillåter det att bära varumärket.
Påståendet vi prövar är följande.
\begin{quote}
  Linuxsystem är Unixliknande: de efterliknar Unix utan att härstamma
  från dess kod.
  BSD-systemen härstammar från Unix.
  macOS härstammar från Unix genom BSD och är ett certifierat UNIX.
  Android är byggt på Linuxkärnan.
\end{quote}
Frågan gäller alltså \emph{sakförhållandet}: den ber om belägg för
härstamning, för efterlikning och för certifiering, ett system i taget,
och svaret är ett olika förhållande för vart och ett.
\end{onlyin}

\section{\ensv{Method}{Metod}} % sync english=52a0f9 swedish=f41b92

\paragraph{\ensv{Known items.}{Kända titlar.}} % sync english=c6cf33 swedish=7be4e1
\begin{onlyin}{english} % sync english=9d88bf swedish=c94841
Most of the evidence for this claim is a primary or authoritative
document that we know by name, so we retrieved it directly (rows K1 to K7
of \cref{tab:unix-lineage-queries}) and verified it by opening it:
Ritchie and Thompson's paper for the origin of Unix (K1);
McKusick's first-hand history of the Berkeley systems (K2), an open
book chapter that no article database indexes;
Torvalds's own announcement of Linux (K3) and the kernel's own
documentation (K4);
The Open Group's register of UNIX certified products (K5), which is the
only authority for the word \enquote{certified};
Apple's developer documentation of Darwin (K6);
and the Android Open Source Project's kernel documentation (K7).
Each is the party that makes the statement, so each is authoritative for
what \emph{it} says and silent on how outsiders classify the system; the
literature search supplies the outside view.
The register's pages and the web documents were read on 18 November 2026;
Ritchie and Thompson's paper was read in the electronic copy that
Eric A.\ Brewer recreated from the 1974 text, because the publisher's
own PDF refused the request.
\end{onlyin}
\begin{onlyin}{swedish}
Det mesta av beläggen för påståendet är primära eller auktoritativa
dokument som vi känner till vid namn, så vi hämtade dem direkt (raderna
K1 till K7 i \cref{tab:unix-lineage-queries}) och kontrollerade dem
genom att öppna dem:
Ritchies och Thompsons artikel för Unix ursprung (K1);
McKusicks förstahandsberättelse om Berkeleysystemen (K2), ett öppet
bokkapitel som ingen artikeldatabas indexerar;
Torvalds eget tillkännagivande av Linux (K3) och kärnans egen
dokumentation (K4);
The Open Groups register över UNIX-certifierade produkter (K5), som är
den enda auktoriteten för ordet \enquote{certifierad};
Apples utvecklardokumentation av Darwin (K6);
och Android Open Source Projects dokumentation av kärnan (K7).
Var och en är den part som gör uttalandet, så var och en är auktoritativ
för vad \emph{den} säger och tyst om hur utomstående klassar systemet;
litteratursökningen ger synen utifrån.
Registrets sidor och webbdokumenten lästes den 18 november 2026;
Ritchies och Thompsons artikel lästes i den elektroniska kopia som
Eric A.\ Brewer återskapade från texten från 1974, eftersom förlagets
egen PDF nekade förfrågan.
\end{onlyin}

\paragraph{\ensv{Literature search.}{Litteratursökning.}} % sync english=e81424 swedish=5ed9eb
\begin{onlyin}{english} % sync english=3f7b28 swedish=c1fb28
On 18 November 2026 we searched with the tool \texttt{scholar} in four
databases: OpenAlex (OA; not open access), IEEE Xplore (IEEE),
Web of Science (WoS) and Scopus.
Semantic Scholar rejected our key (HTTP 403) and DBLP answered every
request with a page that stops automated clients, so neither could be
searched.
IEEE Xplore stopped answering during the search because its daily quota
of 200 requests was spent: it answered for the queries S1 to S4 and
failed for all later ones, which \cref{tab:unix-lineage-queries} marks.
OpenAlex's default search matches full texts and returned thousands of
records for each query, so we searched OpenAlex on title and abstract,
which is the field that Web of Science and Scopus search as well; its
abstracts are sometimes shortened or missing, so it is not the same search
as in Web of Science.
No search had an upper limit.

We built the queries from the concepts of the claim, each in several
wordings, and ran every query with the same keywords on every database
(Web of Science as \texttt{TS=(\dots)}, Scopus as
\texttt{TITLE-ABS-KEY(\dots)}, the others as Boolean free text).
% Ragged right: the items are runs of long quoted query phrases.
\begingroup\raggedright
\begin{description}
  \item[Descent and lineage] \enquote{lineage}, \enquote{genealogy},
    \enquote{family tree}, \enquote{descendants}, \enquote{Unix history},
    \enquote{Unix evolution} (S1); BSD, FreeBSD, NetBSD, OpenBSD,
    \enquote{Berkeley Software Distribution} with Unix and
    history, descendant, derived or origins (S2).
  \item[Imitation] \enquote{Unix-like}, \enquote{Unix like},
    \enquote{Unix clone}, \enquote{clone of Unix}, \enquote{Linux
    history} (S3); the same Unix-like terms with definition, taxonomy or
    classification (C3).
  \item[macOS] macOS, \enquote{Mac OS X}, \enquote{OS X}, XNU, Darwin and
    Leopard with BSD, UNIX, certification, derived or compliance (S4,
    S4b).
  \item[Android] Android with \enquote{Linux kernel} and \enquote{based
    on}, \enquote{built on}, \enquote{derived from} or
    \enquote{Linux-based} (S5); Android with Linux or GNU/Linux and
    adapting, differences, comparison or modifications (S5b).
  \item[Counter-search] Linux with Unix and trademark, certification,
    \enquote{UNIX 03} or \enquote{not Unix} (C1); the certification
    vocabulary alone: \enquote{Single UNIX Specification},
    \enquote{UNIX 03}, \enquote{UNIX certification}, \enquote{UNIX
    trademark}, \enquote{POSIX compliance} (C2); and Android with
    differs, \enquote{not Linux}, \enquote{GNU/Linux} or
    \enquote{Linux distribution} (C4).
\end{description}
\par\endgroup
After the keyword search we asked a search by meaning, one plain-language
question per concept, in OpenAlex's semantic search, and looked up every
work it returned in our records.
It returned works with \enquote{Leopard} or \enquote{Linux for Android}
in the title that the first wordings of the macOS and Android queries
had missed, so we reworded those two concepts (S4b, S5b) and reran them;
these are the two queries that were fitted to what the check showed.
Old books and encyclopaedia entries without an abstract (for example
\emph{The Unix story}) cannot be found by a search on title and abstract,
and the check named several: this is a limit of the method that rewording
does not remove.
\end{onlyin}
\begin{onlyin}{swedish}
Den 18 november 2026 sökte vi med verktyget \texttt{scholar} i fyra
databaser: OpenAlex (OA; inte open access), IEEE Xplore (IEEE),
Web of Science (WoS) och Scopus.
Semantic Scholar avvisade vår nyckel (HTTP 403) och DBLP svarade på
varje förfrågan med en sida som stoppar automatiska klienter, så ingen
av dem kunde genomsökas.
IEEE Xplore slutade svara under sökningen eftersom dess dagskvot på 200
förfrågningar var förbrukad: den svarade på frågorna S1 till S4 och
misslyckades för alla senare, vilket \cref{tab:unix-lineage-queries}
markerar.
OpenAlex standardsökning matchar fulltexter och gav tusentals poster för
varje fråga, så vi sökte i OpenAlex på titel och sammanfattning, som är
det fält som Web of Science och Scopus också söker i; dess
sammanfattningar är ibland förkortade eller saknas, så det är inte samma
sökning som i Web of Science.
Ingen sökning hade någon övre gräns.

Vi byggde frågorna av påståendets begrepp, vart och ett i flera
formuleringar, och körde varje fråga med samma nyckelord i varje databas
(Web of Science som \texttt{TS=(\dots)}, Scopus som
\texttt{TITLE-ABS-KEY(\dots)}, de andra som boolesk fritext).
% Ragged right: the items are runs of long quoted query phrases.
\begingroup\raggedright
\begin{description}
  \item[Släktlinje] \enquote{lineage},
    \enquote{genealogy}, \enquote{family tree}, \enquote{descendants},
    \enquote{Unix history}, \enquote{Unix evolution} (S1); BSD, FreeBSD,
    NetBSD, OpenBSD, \enquote{Berkeley Software Distribution} med Unix
    och history, descendant, derived eller origins (S2).
  \item[Efterlikning] \enquote{Unix-like}, \enquote{Unix like},
    \enquote{Unix clone}, \enquote{clone of Unix}, \enquote{Linux
    history} (S3); samma Unixliknande termer med definition, taxonomy
    eller classification (C3).
  \item[macOS] macOS, \enquote{Mac OS X}, \enquote{OS X}, XNU, Darwin
    och Leopard med BSD, UNIX, certification, derived eller compliance
    (S4, S4b).
  \item[Android] Android med \enquote{Linux kernel} och \enquote{based
    on}, \enquote{built on}, \enquote{derived from} eller
    \enquote{Linux-based} (S5); Android med Linux eller GNU/Linux och
    adapting, differences, comparison eller modifications (S5b).
  \item[Motsökning] Linux med Unix och trademark, certification,
    \enquote{UNIX 03} eller \enquote{not Unix} (C1); certifieringens
    ordförråd för sig: \enquote{Single UNIX Specification},
    \enquote{UNIX 03}, \enquote{UNIX certification}, \enquote{UNIX
    trademark}, \enquote{POSIX compliance} (C2); och Android med
    differs, \enquote{not Linux}, \enquote{GNU/Linux} eller
    \enquote{Linux distribution} (C4).
\end{description}
\par\endgroup
Efter nyckelordssökningen gjorde vi en sökning efter betydelse, en fråga
på vanligt språk per begrepp, i OpenAlex semantiska sökning, och slog
upp varje verk som den gav i våra poster.
Den gav verk med \enquote{Leopard} eller \enquote{Linux for Android} i
titeln som de första formuleringarna av frågorna om macOS och Android
hade missat, så vi formulerade om de två begreppen (S4b, S5b) och körde
dem igen; det är de två frågor som anpassades efter vad kontrollen
visade.
Gamla böcker och uppslagsartiklar utan sammanfattning (till exempel
\emph{The Unix story}) kan inte hittas av en sökning på titel och
sammanfattning, och kontrollen nämnde flera: det är en gräns hos metoden
som omformuleringar inte tar bort.
\end{onlyin}

\begin{table}[htbp]
  % Author's note (verification and reproduction): see the block above
  % the hit-list \input at the end of the chapter.
  \begin{sidecaption}{\ensv{Queries and number of hits per database,
    18 November 2026, tool \texttt{scholar}.
    S = support query, C = counter query, K = known item.
    Databases: OpenAlex (OA; not open access), IEEE Xplore (IEEE),
    Web of Science (WoS), Scopus.
    The searches have no upper limit: the numbers are all hits.
    \enquote{--} means that the search failed (IEEE Xplore's daily
    quota was spent after S4), so the database was not searched for that
    query.
    OpenAlex was searched on title and abstract.
    S4b and S5b reword S4 and S5.
    Semantic Scholar and DBLP could not be searched.}%
    {Frågor och antal träffar per databas, 18 november 2026, verktyget
    \texttt{scholar}.
    S = stödfråga, C = motfråga, K = känd titel.
    Databaser: OpenAlex (OA; inte open access), IEEE Xplore (IEEE),
    Web of Science (WoS), Scopus.
    Sökningarna har ingen övre gräns: talen är alla träffar.
    \enquote{--} betyder att sökningen misslyckades (IEEE Xplores
    dagskvot var förbrukad efter S4), så databasen genomsöktes inte för
    den frågan.
    OpenAlex söktes på titel och sammanfattning.
    S4b och S5b formulerar om S4 och S5.
    Semantic Scholar och DBLP kunde inte genomsökas.}% % sync english=59bdb4 swedish=d56275
  }[tab:unix-lineage-queries]
    \centering\footnotesize
    \begin{tabular}{@{}lp{0.43\linewidth}rrrr@{}}
      \toprule
      & \ensv{Query concept}{Frågans begrepp} & OA & IEEE & WoS & Scopus \\ % sync english=beb164 swedish=80790f
      \midrule
      S1 & \ensv{Unix with lineage, genealogy, family tree, descendants,
           Unix history}{Unix med lineage, genealogy, family tree,
           descendants, Unix history} & 66 & 12 & 18 & 22 \\ % sync english=385d49 swedish=4162d3
      S2 & \ensv{BSD systems with Unix and history, descendant, derived,
           origins}{BSD-systemen med Unix och history, descendant,
           derived, origins} & 71 & 10 & 13 & 22 \\ % sync english=114744 swedish=47d2da
      S3 & \ensv{Linux with Unix-like, Unix clone, Linux history}%
           {Linux med Unix-like, Unix clone, Linux history} % sync english=2dea17 swedish=d071b2
           & 258 & 45 & 66 & 117 \\
      S4 & \ensv{macOS, Mac OS X, XNU with BSD, UNIX 03, certification}%
           {macOS, Mac OS X, XNU med BSD, UNIX 03, certification} % sync english=71d8c2 swedish=acd6d6
           & 187 & 9 & 47 & 58 \\
      S4b & \ensv{S4 reworded: OS X, Darwin, Leopard; derived,
           compliance}{S4 omformulerad: OS X, Darwin, Leopard; derived,
           compliance} & 76 & -- & 22 & 28 \\ % sync english=3e5cd4 swedish=5b1290
      S5 & \ensv{Android with Linux kernel and based on, built on,
           derived}{Android med Linux kernel och based on, built on,
           derived} & 168 & -- & 62 & 102 \\ % sync english=45e1bf swedish=4933e8
      S5b & \ensv{S5 reworded: Linux or GNU/Linux; adapting,
           differences}{S5 omformulerad: Linux eller GNU/Linux;
           adapting, differences} & 136 & -- & 48 & 55 \\ % sync english=9abb0c swedish=d50a81
      C1 & \ensv{Linux and Unix with trademark, certification, not Unix}%
           {Linux och Unix med trademark, certification, not Unix} % sync english=0edf1b swedish=a174a3
           & 22 & -- & 0 & 2 \\
      C2 & \ensv{Certification vocabulary: Single UNIX Specification,
           UNIX 03, POSIX compliance}{Certifieringens ordförråd: Single
           UNIX Specification, UNIX 03, POSIX compliance} % sync english=9ffa88 swedish=2020eb
           & 245 & -- & 101 & 143 \\
      C3 & \ensv{Unix-like with definition, taxonomy, classification}%
           {Unix-like med definition, taxonomy, classification} % sync english=e2d9db swedish=1865fc
           & 72 & -- & 33 & 28 \\
      C4 & \ensv{Android with differs, not Linux, GNU/Linux}%
           {Android med differs, not Linux, GNU/Linux} % sync english=d82e05 swedish=26ccea
           & 16 & -- & 2 & 11 \\
      \midrule
      K1 & \ensv{Ritchie and Thompson 1974}{Ritchie och Thompson 1974} % sync english=1d74ce swedish=e1940e
         & \multicolumn{4}{c}{\ensv{known item}{känd titel}} \\ % sync english=95335c swedish=f439fe
      K2 & \ensv{McKusick on Berkeley Unix}{McKusick om Berkeley Unix} % sync english=c8cb82 swedish=81b9a7
         & \multicolumn{4}{c}{\ensv{known item}{känd titel}} \\ % sync english=95335c swedish=f439fe
      K3 & \ensv{Torvalds's announcement}{Torvalds tillkännagivande} % sync english=bec467 swedish=400270
         & \multicolumn{4}{c}{\ensv{known item}{känd titel}} \\ % sync english=95335c swedish=f439fe
      K4 & \ensv{Kernel documentation}{Kärnans dokumentation} % sync english=7774ff swedish=0cd16f
         & \multicolumn{4}{c}{\ensv{known item}{känd titel}} \\ % sync english=95335c swedish=f439fe
      K5 & \ensv{The Open Group register}{The Open Groups register} % sync english=fe0ef7 swedish=8e370b
         & \multicolumn{4}{c}{\ensv{known item}{känd titel}} \\ % sync english=95335c swedish=f439fe
      K6 & \ensv{Apple's Darwin documentation}%
           {Apples dokumentation av Darwin} % sync english=04dd6b swedish=8d4db7
         & \multicolumn{4}{c}{\ensv{known item}{känd titel}} \\ % sync english=95335c swedish=f439fe
      K7 & \ensv{Android kernel documentation}%
           {Androids dokumentation av kärnan} % sync english=36c1bf swedish=f862d1
         & \multicolumn{4}{c}{\ensv{known item}{känd titel}} \\ % sync english=95335c swedish=f439fe
      \bottomrule
    \end{tabular}
  \end{sidecaption}
\end{table}

\paragraph{\ensv{Studies the queries had to return.}%
  {Studier som frågorna måste ge.}} % sync english=50f120 swedish=6aaaf2
\begin{onlyin}{english} % sync english=7ac9a5 swedish=ba3587
Before relying on the search we named studies that we knew to bear on the
claim: Spinellis's repository of Unix history\autocite{Spinellis2016Repository}
and his architecture study of the FreeBSD lineage%
\autocite{SpinellisAvgeriou2021} for descent, and Takahashi and
Takamatsu's history of the licence dispute\autocite{Takahashi2013Unix}
for Linux.
The first two came back from S1 and S2 in at least three of the four
databases; Takahashi and Takamatsu came back from S3, in OpenAlex only.
The study of how Android adapts the Linux kernel%
\autocite{Khomh2012Adapting} came back only from the reworded S5b, which
is why the Android concept has two queries.
Ritchie and Thompson's paper and McKusick's chapter did not come back,
because their titles and abstracts contain none of the concept words;
that is the reason they are known items.
\end{onlyin}
\begin{onlyin}{swedish}
Innan vi litade på sökningen namngav vi studier som vi visste rör
påståendet: Spinellis arkiv över Unix
historia\autocite{Spinellis2016Repository} och hans arkitekturstudie av
FreeBSD:s släktlinje\autocite{SpinellisAvgeriou2021} för härstamningen,
och Takahashis och Takamatsus historia om licenstvisten%
\autocite{Takahashi2013Unix} för Linux.
De två första kom tillbaka från S1 och S2 i minst tre av de fyra
databaserna; Takahashi och Takamatsu kom tillbaka från S3, bara i
OpenAlex.
Studien av hur Android anpassar Linuxkärnan%
\autocite{Khomh2012Adapting} kom tillbaka bara från den omformulerade
S5b, och därför har begreppet Android två frågor.
Ritchies och Thompsons artikel och McKusicks kapitel kom inte tillbaka,
eftersom deras titlar och sammanfattningar inte innehåller något av
begreppsorden; det är skälet till att de är kända titlar.
\end{onlyin}

\paragraph{\ensv{Screening and exclusion criteria.}%
  {Gallring och uteslutningskriterier.}} % sync english=c0d3d3 swedish=e2519b
\begin{onlyin}{english} % sync english=24ba07 swedish=ae7b02
The searches returned 1437 records, 1267 of them distinct.
We classified all of them against the claim with a language model
(\texttt{claude-sonnet-5-5}, without examples), one category each:
supports the claim, qualifies it, adjacent subtopic, or off-topic false
hit.
We read the title of every record that the model labelled as supporting
(100) or qualifying (nine), and the abstract where there was one,
and decided ourselves; many have no abstract.
A record is excluded if it is one of the following.
\begin{description}
  \item[Tool use only] It uses Unix, Linux, Android or macOS only as the
    environment, platform or tool of the study (738 records).
  \item[App level only] It studies applications, malware, permissions or
    security on one of these systems and says nothing about how the
    system relates to Unix or Linux (162 records).
  \item[Kernel internals] It studies one mechanism inside a kernel
    (scheduling, drivers, memory, file systems) and says nothing about
    lineage or classification (275 records).
  \item[No evidence] It is an opinion, a user manual or an overview with
    no evidence of its own, no source and no stated method
    (229 records).
\end{description}
(A record can meet more than one criterion.)
Of these 109 records we kept 31, which are 26 distinct works; the table
at the end of the chapter lists them, and marks those that the text
cites.
The other 78, mostly user manuals and security books about Mac OS X and
compilations of encyclopaedia text, were excluded as \emph{no evidence} or
\emph{app level only}.
\end{onlyin}
\begin{onlyin}{swedish}
Sökningarna gav 1437 poster, varav 1267 olika.
Vi klassade alla mot påståendet med en språkmodell
(\texttt{claude-sonnet-5-5}, utan exempel), en kategori var: stöder
påståendet, kvalificerar det, angränsande delämne eller felträff utanför
ämnet.
Vi läste titeln på varje post som modellen märkte som stödjande (100)
eller kvalificerande (nio), och sammanfattningen där det fanns en, och
avgjorde själva; många saknar sammanfattning.
En post utesluts om den är något av följande.
\begin{description}
  \item[Bara verktyg] Den använder Unix, Linux, Android eller macOS bara
    som studiens miljö, plattform eller verktyg (738 poster).
  \item[Bara appnivå] Den studerar program, skadlig kod, behörigheter
    eller säkerhet i ett av systemen och säger inget om hur systemet
    förhåller sig till Unix eller Linux (162 poster).
  \item[Kärnans inre] Den studerar en mekanism inuti en kärna
    (schemaläggning, drivrutiner, minne, filsystem) och säger inget om
    släktlinje eller klassificering (275 poster).
  \item[Inga belägg] Den är en åsikt, en användarmanual eller en
    översikt utan egna belägg, utan källa och utan angiven metod (229
    poster).
\end{description}
(En post kan uppfylla mer än ett kriterium.)
Av de 109 posterna behöll vi 31, som är 26 olika verk; tabellen i slutet
av kapitlet listar dem, och markerar dem som texten citerar.
De övriga 78, mest användarmanualer och säkerhetsböcker om Mac OS X och
sammanställningar av uppslagstext, uteslöts som \emph{inga belägg} eller
\emph{bara appnivå}.
\end{onlyin}

\section{\ensv{Results}{Resultat}} % sync english=219c4a swedish=6e45f9

\paragraph{\ensv{Unix, the origin.}{Unix, ursprunget.}} % sync english=be7e8c swedish=94b3a4
\begin{onlyin}{english} % sync english=3ab957 swedish=092c08
Ritchie and Thompson describe the system at Bell Laboratories and give
its early history:
\textcquote{RitchieThompson1974}{There have been three versions of UNIX.}
They add that the system was first written in assembly language and
then \textcquote{RitchieThompson1974}{rewritten in C} in the summer of
1973, and it is this rewrite that made it portable.
\enquote{Descend from Unix} therefore means a code history that leads
back to this system.
\end{onlyin}
\begin{onlyin}{swedish}
Ritchie och Thompson beskriver systemet vid Bell Laboratories och ger
dess tidiga historia:
\foreigntextcquote{english}{RitchieThompson1974}{There have been three
versions of UNIX.}
De tillägger att systemet först skrevs i assembler och sedan
\foreigntextcquote{english}{RitchieThompson1974}{rewritten in C}
sommaren 1973, och det är den omskrivningen som gjorde det flyttbart.
\enquote{Härstamma från Unix} betyder därför en kodhistoria som leder
tillbaka till det här systemet.
\end{onlyin}

\paragraph{\ensv{The BSD systems descend from Unix.}%
  {BSD-systemen härstammar från Unix.}} % sync english=647f9c swedish=640f43
\begin{onlyin}{english} % sync english=3d863f swedish=db5649
McKusick, who worked on the Berkeley systems himself, tells how Berkeley
began with Bell Labs' Unix and added to it.
The first release for the VAX machine
\textcquote{McKusick1999Berkeley}{included the virtual memory kernel, the
standard 32/V utilities, and the additions from 2BSD}, and 32/V was
Bell Labs' own port of Unix to the VAX.
Spinellis's repository of Unix history has the same lineage as its
backbone: it was built from \textcquote{Spinellis2016Repository}{24
snapshots of systems developed at Bell Labs, the University of California
at Berkeley, and the 386BSD team, two legacy repositories, and the modern
repository of the open source FreeBSD system}, and it speaks of
\textcquote{Spinellis2016Repository}{BSD Unix descendants}, among them
FreeBSD, NetBSD and OpenBSD.
In the same author's architecture study, a footnote explains that
\enquote{Unix} stands for both the Bell Labs systems and for
\textcquote{SpinellisAvgeriou2021}{Unix-like systems, such as FreeBSD,
that descended from them}.
The conference version of the repository paper%
\autocite{Spinellis2015Repository}, a measurement of the similarity of
the BSD versions%
\autocite{Yamamoto2005Measuring} and an analysis that treats
FreeBSD, NetBSD, OpenBSD and others as one family of software
variants\autocite{Tenev2012Bioinformatics} corroborate that the BSD
systems form one lineage.

There is a qualification that matters for how we word it.
In 1994 the University of California released 4.4BSD-Lite after a
lawsuit with AT\&T's subsidiary.
McKusick recounts the earlier work of
\textcquote{McKusick1999Berkeley}{going over the entire distribution,
file by file, removing code that had originated in the 32/V release},
and he reports that after the settlement
\textcquote{McKusick1999Berkeley}{all the BSD groups that were doing
releases at that time, BSDI, NetBSD, and FreeBSD, had to restart their
code base with the 4.4BSD-Lite sources}.
So the BSD systems descend from Unix in the sense of a continuous code
history, but since 1994 they rest on a release that Unix System
Laboratories agreed not to sue over, and they are not certified: the
lists in The Open Group's register of UNIX certified products named no
BSD system when we read them.
\end{onlyin}
\begin{onlyin}{swedish}
McKusick, som själv arbetade med Berkeleysystemen, berättar hur Berkeley
började med Unix från Bell Labs och byggde vidare på det.
Den första utgåvan för VAX-maskinen
\foreigntextcquote{english}{McKusick1999Berkeley}{included the virtual
memory kernel, the standard 32/V utilities, and the additions from
2BSD}, och 32/V var Bell Labs egen överföring av Unix till VAX.
Spinellis arkiv över Unix historia har samma släktlinje som ryggrad: det
byggdes av \foreigntextcquote{english}{Spinellis2016Repository}{24
snapshots of systems developed at Bell Labs, the University of California
at Berkeley, and the 386BSD team, two legacy repositories, and the modern
repository of the open source FreeBSD system}, och det talar om
\foreigntextcquote{english}{Spinellis2016Repository}{BSD Unix
descendants}, bland dem FreeBSD, NetBSD och OpenBSD.
I samma författares arkitekturstudie förklarar en fotnot att
\enquote{Unix} står både för systemen från Bell Labs och för
\foreigntextcquote{english}{SpinellisAvgeriou2021}{Unix-like systems,
such as FreeBSD, that descended from them}.
Konferensversionen av artikeln om arkivet%
\autocite{Spinellis2015Repository}, en mätning av hur lika
BSD-versionerna är\autocite{Yamamoto2005Measuring} och en analys som
behandlar FreeBSD, NetBSD, OpenBSD och andra som en familj av
programvarianter\autocite{Tenev2012Bioinformatics} bekräftar att
BSD-systemen bildar en släktlinje.

Det finns en kvalifikation som spelar roll för hur vi formulerar det.
År 1994 släppte University of California 4.4BSD-Lite efter en
rättstvist med AT\&T:s dotterbolag.
McKusick berättar om det tidigare arbetet med att
\foreigntextcquote{english}{McKusick1999Berkeley}{going over the entire
distribution, file by file, removing code that had originated in the
32/V release}, och han berättar att efter förlikningen
\foreigntextcquote{english}{McKusick1999Berkeley}{all the BSD groups
that were doing releases at that time, BSDI, NetBSD, and FreeBSD, had to
restart their code base with the 4.4BSD-Lite sources}.
BSD-systemen härstammar alltså från Unix i betydelsen en obruten
kodhistoria, men sedan 1994 vilar de på en utgåva som Unix System
Laboratories gick med på att inte stämma för, och de är inte
certifierade: listorna i The Open Groups register över UNIX-certifierade
produkter nämnde inget BSD-system när vi läste dem.
\end{onlyin}

\paragraph{\ensv{Linux is Unix-like, written from scratch.}%
  {Linux är Unixliknande, skrivet från grunden.}} % sync english=13287c swedish=7fca7f
\begin{onlyin}{english} % sync english=3700b1 swedish=e497ad
Torvalds announced the system in August 1991 as
\textcquote{Torvalds1992History}{a (free) operating system (just a hobby,
won't be big and professional like gnu) for 386(486) AT clones}, and
added that it is \textcquote{Torvalds1992History}{free of any minix
code}.
The kernel's documentation still says what Linux is:
\textcquote{LinuxKernelDocs2026Readme}{a clone of the operating system
Unix, written from scratch by Linus Torvalds with assistance from a
loosely-knit team of hackers across the Net}, which
\textcquote{LinuxKernelDocs2026Readme}{aims towards POSIX and Single UNIX
Specification compliance}.
An independent business-history article agrees:
\textcquote{Takahashi2013Unix}{Torvalds created Linux from the ground
up as a student at the University of Helsinki, not infringing on the
AT\&T license}, in the years of the licence dispute over BSD.
That is the relation we call Unix-like: imitation of the interface, with
no code in common at the start.
\end{onlyin}
\begin{onlyin}{swedish}
Torvalds tillkännagav systemet i augusti 1991 som
\foreigntextcquote{english}{Torvalds1992History}{a (free) operating
system (just a hobby, won't be big and professional like gnu) for
386(486) AT clones}, och lade till att det är
\foreigntextcquote{english}{Torvalds1992History}{free of any minix
code}.
Kärnans dokumentation säger fortfarande vad Linux är:
\foreigntextcquote{english}{LinuxKernelDocs2026Readme}{a clone of the
operating system Unix, written from scratch by Linus Torvalds with
assistance from a loosely-knit team of hackers across the Net}, som
\foreigntextcquote{english}{LinuxKernelDocs2026Readme}{aims towards
POSIX and Single UNIX Specification compliance}.
En oberoende artikel i företagshistoria håller med:
\foreigntextcquote{english}{Takahashi2013Unix}{Torvalds created Linux
from the ground up as a student at the University of Helsinki, not
infringing on the AT\&T license}, under åren med licenstvisten om BSD.
Det är det förhållande som vi kallar Unixliknande: efterlikning av
gränssnittet, utan någon gemensam kod från början.
\end{onlyin}

\paragraph{\ensv{\enquote{Unix-like} is a loose word.}%
  {\enquote{Unixliknande} är ett löst ord.}} % sync english=2d419a swedish=9d034d
\begin{onlyin}{english} % sync english=910232 swedish=5e5e8c
The word has no definition that everyone uses.
A user group tried to set criteria for Unix-like systems already in
1982\autocite{Mason1982Criteria}.
Spinellis and Avgeriou use \enquote{Unix-like} for the descendants as
well, such as FreeBSD (above).
A textbook chapter says that BSD and Linux alike are
\textcquote{LaLau2021Unix}{systems that descend from or are inspired by
Unix},
and a book chapter calls BSD, Linux and Mac OS X
\textcquote{MacForLinux2009Comparison}{branches straight from a single tap
root: the UNIX operating system}.
An article on Linux distributions calls Linux
\textcquote{Jung2017Trend}{considered to be a subset of the UNIX
operating system} (its Korean abstract says, in our reading, that Linux
was regarded as an imitation of Unix), and a book on open source
opens with a chapter on \enquote{How UNIX Gave Birth to
Linux}\autocite{Meeker2008OpenSource}.
Where sources write that Linux descends from Unix they mean its design
inheritance, and the code history above does not support more.
The deck therefore says \enquote{descend} for the BSD systems and
\enquote{imitate} or \enquote{Unix-like} for Linux.
One study found code flowing between the systems after their separate
beginnings: code copied in one direction or the other between Linux and
FreeBSD and OpenBSD, in the direction their licences allow%
\autocite{German2009Siblings}.
That is not descent, but it shows why the systems feel alike.
\end{onlyin}
\begin{onlyin}{swedish}
Ordet har ingen definition som alla använder.
En användarförening försökte sätta upp kriterier för Unixliknande system
redan 1982\autocite{Mason1982Criteria}.
Spinellis och Avgeriou använder \enquote{Unixliknande} även för
ättlingarna, som FreeBSD (se ovan).
Ett läroboksavsnitt säger att BSD och Linux båda är
\foreigntextcquote{english}{LaLau2021Unix}{systems that descend from or
are inspired by Unix},
och ett bokkapitel kallar BSD, Linux och Mac OS X
\foreigntextcquote{english}{MacForLinux2009Comparison}{branches straight
from a single tap root: the UNIX operating system}.
En artikel om Linuxdistributioner säger att Linux
\foreigntextcquote{english}{Jung2017Trend}{considered to be a subset of
the UNIX operating system} (dess koreanska sammanfattning säger, i vår
läsning, att Linux sågs som en efterlikning av Unix), och en bok om
öppen källkod börjar med ett kapitel om \enquote{How UNIX Gave Birth to
Linux}\autocite{Meeker2008OpenSource}.
Där källor skriver att Linux härstammar från Unix menar de dess arv i
utformningen, och kodhistorien ovan stöder inte mer.
Texten säger därför \enquote{härstammar} om BSD-systemen och
\enquote{efterliknar} eller \enquote{Unixliknande} om Linux.
En studie fann kod som flyttats mellan systemen efter deras skilda
början: kod kopierad åt ena eller andra hållet mellan Linux och FreeBSD
och OpenBSD, åt det håll som deras licenser
tillåter\autocite{German2009Siblings}.
Det är inte härstamning, men det visar varför systemen känns lika.
\end{onlyin}

\paragraph{\ensv{macOS descends from Unix through BSD.}%
  {macOS härstammar från Unix genom BSD.}} % sync english=5a59eb swedish=ca72c6
\begin{onlyin}{english} % sync english=5f506b swedish=b3a668
Apple's documentation of the kernel says that
\textcquote{Apple2013KernelProgramming}{A large portion of this technology
is derived from FreeBSD, a version of 4.4BSD}.
The kernel of macOS is XNU, which combines Mach, a BSD layer and the
I/O Kit; a book on its internals describes it so%
\autocite{Levin2012Internals}.
The Apple page is archived and dated 2013, and describes OS X, so it
shows where the BSD layer came from, not what today's macOS is made of.
\end{onlyin}
\begin{onlyin}{swedish}
Apples dokumentation av kärnan säger att
\foreigntextcquote{english}{Apple2013KernelProgramming}{A large portion
of this technology is derived from FreeBSD, a version of 4.4BSD}.
Kärnan i macOS är XNU, som förenar Mach, ett BSD-lager och I/O Kit; en
bok om dess inre beskriver den så\autocite{Levin2012Internals}.
Apples sida är arkiverad och daterad 2013, och beskriver OS X, så den
visar var BSD-lagret kom ifrån, inte vad dagens macOS är gjort av.
\end{onlyin}

\paragraph{\ensv{macOS is a certified UNIX, with a caveat.}%
  {macOS är ett certifierat UNIX, med ett förbehåll.}} % sync english=f893c4 swedish=6e4a3a
\begin{onlyin}{english} % sync english=b1272c swedish=0fcb92
The Open Group states the rule:
\textcquote{OpenGroup2026UNIX}{Suppliers are granted a license to use the
UNIX trademark in connection with a product when: The product has
demonstrated passing the applicable certification test suites.}
Its register has a page for each registered product.
The pages that we opened name Apple Inc.\ as the company and macOS as the
product, with a registration date: version 10.12 Sierra on 26 September
2016, 10.13 High Sierra on 14 September 2017, 11.0 Big Sur on 6 October
2020, 12.0 Monterey on 12 October 2021, 13.0 Ventura on 9 August 2022,
14.0 Sonoma on 9 August 2023, 15.0 Sequoia on 12 September 2024 and
26.0 Tahoe, for Apple silicon and for Intel, on 29 August 2025%
\autocite{OpenGroup2026Register}.
The caveat is that the register's overview pages, which list the
certified products by standard and by company, did not show Apple on 18
November 2026, although the product pages are still served; we could
not find an explanation in a source that we could open.
A web search also returned an Italian-language article that claims that
Apple was removed from the list, and a commentary that says that the
tested configuration differs from a default installation; both are
secondary, we did not open them, and we do not rely on them.
So the claim that macOS is a certified UNIX is backed by the register's
own product pages, with the date of the latest one, and not by its
overview lists.
\end{onlyin}
\begin{onlyin}{swedish}
The Open Group anger regeln:
\foreigntextcquote{english}{OpenGroup2026UNIX}{Suppliers are granted a
license to use the UNIX trademark in connection with a product when: The
product has demonstrated passing the applicable certification test
suites.}
Dess register har en sida för varje registrerad produkt.
De sidor som vi öppnade anger Apple Inc.\ som företag och macOS som
produkt, med ett registreringsdatum: version 10.12 Sierra den 26
september 2016, 10.13 High Sierra den 14 september 2017, 11.0 Big Sur
den 6 oktober 2020, 12.0 Monterey den 12 oktober 2021, 13.0 Ventura den
9 augusti 2022, 14.0 Sonoma den 9 augusti 2023, 15.0 Sequoia den 12
september 2024 och 26.0 Tahoe, för Apple silicon och för Intel, den 29
augusti 2025\autocite{OpenGroup2026Register}.
Förbehållet är att registrets översiktssidor, som listar de
certifierade produkterna efter standard och efter företag, inte visade
Apple den 18 november 2026, fastän produktsidorna fortfarande finns; vi
kunde inte hitta någon förklaring i en källa som vi kunde öppna.
En webbsökning gav också en italienskspråkig artikel som påstår att
Apple har tagits bort från listan, och en kommentar som säger att den
testade konfigurationen skiljer sig från en standardinstallation; båda
är sekundära, vi öppnade dem inte, och vi bygger inte på dem.
Påståendet att macOS är ett certifierat UNIX stöds alltså av registrets
egna produktsidor, med datumet för den senaste, och inte av dess
översiktslistor.
\end{onlyin}

\paragraph{\ensv{Linux and certification.}{Linux och certifiering.}} % sync english=72411b swedish=eb4aa1
\begin{onlyin}{english} % sync english=027526 swedish=2f87ca
The Open Group's page on the UNIX standard says of the look-alikes that
\textcquote{OpenGroup2026UNIX}{Even a Linux distribution has become UNIX
certified continuing the trend and value of basing operating systems on the
UNIX standard}.
This qualifies a bare \enquote{Linux is not UNIX}: certification is given
to a product, not to a code history, and a Linux-based product has passed
the test suites.
The page does not name the distribution, and the product lists in the
register that we read contain no Linux distribution.
We therefore say that Linux systems are Unix-like, and that a registered
Linux-based product would be a certified UNIX, which we could not find
in the register.
\end{onlyin}
\begin{onlyin}{swedish}
The Open Groups sida om UNIX-standarden säger om efterliknarna att
\foreigntextcquote{english}{OpenGroup2026UNIX}{Even a Linux distribution
has become UNIX certified continuing the trend and value of basing
operating systems on the UNIX standard}.
Det kvalificerar ett naket \enquote{Linux är inte UNIX}: certifiering
ges till en produkt, inte till en kodhistoria, och en Linuxbaserad
produkt har klarat testsviterna.
Sidan namnger inte distributionen, och produktlistorna i registret som
vi läste innehåller ingen Linuxdistribution.
Vi säger därför att Linuxsystem är Unixliknande, och att en registrerad
Linuxbaserad produkt skulle vara ett certifierat UNIX, vilket vi inte
kunde hitta i registret.
\end{onlyin}

\paragraph{\ensv{Android is built on the Linux kernel, not on a Linux
  system.}{Android är byggt på Linuxkärnan, inte på ett Linuxsystem.}} % sync english=3a3117 swedish=b3c1f8
\begin{onlyin}{english} % sync english=e18a33 swedish=8668b8
The Android Open Source Project says that
\textcquote{Android2026Kernel}{The Android kernel is based on an upstream
Linux Long Term Supported (LTS) kernel}.
An empirical study from 2012 measured the adaptation:
\textcquote{Khomh2012Adapting}{only 0.7\% of files from the Linux kernel
are modified when reused for a mobile platform}.
A textbook chapter says that \textcquote{Tarkoma2014Mobile}{Android is
based on the Linux kernel, iOS on the Mac OS X kernel}, and a student
report repeats it\autocite{Goel2021Android}.
The sources that qualify the claim all concern the user space.
A student project reports that Android
\textcquote{Gupta2010DebianRunner}{runs atop the Linux kernel, yet lacks
the ability to natively execute standard Linux desktop applications};
a developer's guide to Android below the framework advertises that it
shows how Android differs from its Linux roots\autocite{Yaghmour2013Embedded};
and a paper proposes a thin layer to give Android applications POSIX.1
conformance\autocite{Nafees2013POSIX}.
We found no source that calls Android a UNIX, and it is on none of the
register's lists.
None of these sources says whether an Android device offers the command
line you learn here; they say only that the kernel is Linux's and that
the user space is not that of a Linux desktop.
\end{onlyin}
\begin{onlyin}{swedish}
Android Open Source Project säger att
\foreigntextcquote{english}{Android2026Kernel}{The Android kernel is
based on an upstream Linux Long Term Supported (LTS) kernel}.
En empirisk studie från 2012 mätte anpassningen:
\foreigntextcquote{english}{Khomh2012Adapting}{only 0.7\% of files from
the Linux kernel are modified when reused for a mobile platform}.
Ett läroboksavsnitt säger att
\foreigntextcquote{english}{Tarkoma2014Mobile}{Android is based on the
Linux kernel, iOS on the Mac OS X kernel}, och en studentrapport
upprepar det\autocite{Goel2021Android}.
De källor som kvalificerar påståendet gäller alla användarmiljön.
Ett studentprojekt rapporterar att Android
\foreigntextcquote{english}{Gupta2010DebianRunner}{runs atop the Linux
kernel, yet lacks the ability to natively execute standard Linux desktop
applications};
en utvecklarhandbok om Android under ramverket anger att den visar hur
Android skiljer sig från sina rötter i
Linux\autocite{Yaghmour2013Embedded};
och en artikel föreslår ett tunt lager som ger Androidprogram
POSIX.1-överensstämmelse\autocite{Nafees2013POSIX}.
Vi hittade ingen källa som kallar Android ett UNIX, och det finns inte
på någon av registrets listor.
Ingen av källorna säger om en Androidenhet erbjuder den kommandorad som
du lär dig här; de säger bara att kärnan är Linux och att
användarmiljön inte är ett Linuxskrivbords.
\end{onlyin}

\paragraph{\ensv{The family tree that the old deck linked.}%
  {Släktträdet som det gamla materialet länkade till.}} % sync english=a3dbd7 swedish=0b985c
\begin{onlyin}{english} % sync english=65d6c4 swedish=44ffaa
The previous version of this material linked to a web page with a family
tree of operating systems\autocite{Eylenburg2026FamilyTree}.
It is a personal compilation: the page names no author and gives no
sources, and says itself that it includes only
\textcquote{Eylenburg2026FamilyTree}{a few selected Linux/BSD/Solaris
distributions}.
It is useful as a picture of the landscape that you can explore, but it
is not evidence for any claim in this appendix, and the text does not
cite it as such.
\end{onlyin}
\begin{onlyin}{swedish}
Den tidigare versionen av det här materialet länkade till en webbsida
med ett släktträd över operativsystem\autocite{Eylenburg2026FamilyTree}.
Det är en personlig sammanställning: sidan namnger ingen författare och
anger inga källor, och säger själv att den bara tar med
\foreigntextcquote{english}{Eylenburg2026FamilyTree}{a few selected
Linux/BSD/Solaris distributions}.
Den är användbar som en bild av landskapet som du kan utforska, men den
är inget belägg för något påstående i den här bilagan, och texten
citerar den inte som det.
\end{onlyin}

\section{\ensv{Discussion and limitations}{Diskussion och begränsningar}} % sync english=634094 swedish=5d16d0

\begin{onlyin}{english} % sync english=baab3d swedish=c5239d
The evidence is of two kinds, and the limits differ.
The statements about each system's origin come from documents written by
the people involved or by the system's own project (Ritchie and Thompson,
McKusick, Torvalds, the kernel documentation, Apple, the Android Open
Source Project) and from three peer-reviewed studies; these documents are
authoritative on what they say, and interested in saying it.
The statements about certification come from The Open Group, which sells
the certification.
The outside literature that we found adds corroboration for the BSD
lineage and for the Android kernel, and it adds the observation that
authors use \enquote{Unix-like} and \enquote{descend} loosely.
It adds nothing on the certification of macOS.
We found no study that tests the claim as a whole, no source that
refutes any part of it, and the qualifying sources concern wording and the
Android user space.

The limits are facts about our data and tools.
The screening was at title level: we decided on the 109 records that the
model sorted as bearing on the claim mostly by their titles, reading an
abstract where one was available, so a relevant study with an
unrevealing title may have been excluded.
The searches did not go through the university's network, and Scopus
returns abstracts only there, so the Scopus records came without
abstracts; for those records both the model and we had the title alone.
Semantic Scholar and DBLP could not be searched, and IEEE Xplore was
searched only for the first four queries, so conference papers that only
IEEE Xplore indexes are possibly missing for the macOS, Android and
counter queries.
OpenAlex searched title and abstract only, and many older records and
books have no abstract; a search by meaning found works that the keywords
had missed, and the repair queries recovered those that bear on macOS and
Android.
Of the cited sources we read these in full: Ritchie and Thompson,
McKusick, Torvalds, the kernel documentation, the Open Group's pages,
Apple's and Android's pages, Spinellis's journal papers, Takahashi and
Takamatsu and the Android student report.
The following were read in abstract or publisher's description only: the
studies by Khomh and others, by German and others, by Mason, by Jung, by
Yamamoto and others, and by Tenev and Duszynski, the chapters by Tarkoma
and others and by La Lau, the two books on the internals of macOS and of
Android, whose publishers' descriptions we saw, the book on open source,
of which we saw only a table of contents, and the student report on
DebianRunner.
No page numbers were available for the web documents, and the quotation
from Ritchie and Thompson has no page because the copy that we read has
none of the original's page numbers.
The Apple page is archived, and the Open Group pages carry no revision
date.
The Open Group's register lists did not show macOS on the day we read
them, as described above; a later read may show something else.
One subject that this search did not settle: the search by meaning
returned works on a litigation between SCO and others (one is titled
\emph{Closed Source Fights Back: SCO vs.\ The World}), whose subject and
outcome we did not verify.
If it concerned whether Linux contains Unix code, it should be read
before the claim about Linux is repeated.
\end{onlyin}
\begin{onlyin}{swedish}
Beläggen är av två slag, och gränserna skiljer sig åt.
Uppgifterna om varje systems ursprung kommer från dokument som skrivits
av de inblandade eller av systemets eget projekt (Ritchie och Thompson,
McKusick, Torvalds, kärnans dokumentation, Apple, Android Open Source
Project) och från tre referentgranskade studier; dokumenten är
auktoritativa för vad de säger, och har ett intresse av att säga det.
Uppgifterna om certifiering kommer från The Open Group, som säljer
certifieringen.
Den utomstående litteratur som vi hittade bekräftar BSD-släktlinjen och
Androids kärna, och den tillför iakttagelsen att författare använder
\enquote{Unixliknande} och \enquote{härstamma} löst.
Den tillför inget om certifieringen av macOS.
Vi hittade ingen studie som prövar påståendet som helhet, ingen källa som
motsäger någon del av det, och de kvalificerande källorna gäller
ordval och Androids användarmiljö.

Gränserna är fakta om våra data och verktyg.
Gallringen gjordes på titelnivå: vi avgjorde de 109 poster som modellen
sorterade som relevanta för påståendet mest efter deras titlar, och
läste en sammanfattning där en fanns, så en relevant studie med en
intetsägande titel kan ha uteslutits.
Sökningarna gick inte genom universitetets nätverk, och Scopus ger
sammanfattningar bara där, så Scopus poster kom utan sammanfattningar;
för de posterna hade både modellen och vi bara titeln.
Semantic Scholar och DBLP kunde inte genomsökas, och IEEE Xplore
genomsöktes bara för de fyra första frågorna, så konferensbidrag som
bara IEEE Xplore indexerar kan saknas för frågorna om macOS, Android och
för motfrågorna.
OpenAlex söktes bara på titel och sammanfattning, och många äldre poster
och böcker saknar sammanfattning; en sökning efter betydelse hittade
verk som nyckelorden hade missat, och de reparerande frågorna fick
tillbaka dem som rör macOS och Android.
Av de citerade källorna läste vi dessa i sin helhet: Ritchie och
Thompson, McKusick, Torvalds, kärnans dokumentation, The Open Groups
sidor, Apples och Androids sidor, Spinellis tidskriftsartiklar, Takahashi
och Takamatsu samt studentrapporten om Android.
Följande lästes bara i sammanfattning eller förlagets beskrivning:
studierna av Khomh med flera, av German med flera, av Mason, av Jung, av
Yamamoto med flera och av Tenev och Duszynski, kapitlen av Tarkoma med
flera och av La Lau, de två böckerna om det inre i macOS och i Android,
vars förlagsbeskrivningar vi såg, boken om öppen källkod, där vi bara
såg en innehållsförteckning, och studentrapporten om DebianRunner.
Inga sidnummer fanns för webbdokumenten, och citatet från Ritchie och
Thompson har ingen sida eftersom kopian som vi läste saknar originalets
sidnummer.
Apples sida är arkiverad, och The Open Groups sidor har inget
revisionsdatum.
The Open Groups registerlistor visade inte macOS den dag vi läste dem,
som beskrivs ovan; en senare läsning kan visa något annat.
En sak som sökningen inte avgjorde: sökningen efter betydelse gav verk
om en rättstvist mellan SCO och andra (ett har titeln \emph{Closed
Source Fights Back: SCO vs.\ The World}), vars ämne och utgång vi inte
kontrollerade.
Om den gällde huruvida Linux innehåller kod från Unix, bör den läsas
innan påståendet om Linux upprepas.
\end{onlyin}

\section{\ensv{Conclusion}{Slutsats}} % sync english=89db2a swedish=c3cc6f

\begin{onlyin}{english} % sync english=9fad17 swedish=ef5ae3
Of today's systems, the BSD systems, and macOS through them, descend from
Unix; Linux does not, but imitates it; and Android is built on Linux's
kernel.
In detail, the BSD systems have a code history that leads from Bell
Labs' Unix through Berkeley and 386BSD to FreeBSD, NetBSD and OpenBSD,
although since 1994 they rest on a release that its former owner agreed
not to sue over.
macOS has a BSD layer derived from FreeBSD and its product pages are in
The Open Group's register of UNIX certified products, the latest for
macOS 26.0 Tahoe, registered on 29 August 2025; the register's overview
lists did not show Apple on 18 November 2026.
Linux was written from scratch as a clone of Unix, and the Linux kernel's
own documentation and Torvalds's first announcement say so; it is
Unix-like in the sense of imitation, even though a Linux-based product
has been registered as a UNIX.
Android uses a modified Linux kernel (a study from 2012 found that
only 0.7\% of the kernel's files were modified) and a user space of its
own, so it is built on Linux's kernel but is not a Linux desktop system.
The deck's main text may therefore say that Linux, macOS and the BSD
systems are alike because the BSD systems and macOS descend from Unix and
Linux imitates it, and that Android builds on the Linux kernel.
It should not say that Linux descends from Unix, that the BSD systems
are certified UNIX systems, or that you can expect the same command line
on Android: no source that we found says so.
\end{onlyin}
\begin{onlyin}{swedish}
Av dagens system härstammar BSD-systemen, och macOS genom dem, från
Unix; Linux gör det inte, men efterliknar det; och Android är byggt på
Linux kärna.
I detalj har BSD-systemen en kodhistoria som leder från Unix vid Bell
Labs genom Berkeley och 386BSD till FreeBSD, NetBSD och OpenBSD, fastän
de sedan 1994 vilar på en utgåva som dess tidigare ägare gick med på att
inte stämma för.
macOS har ett BSD-lager som härletts från FreeBSD, och dess
produktsidor finns i The Open Groups register över UNIX-certifierade
produkter, den senaste för macOS 26.0 Tahoe, registrerad den 29 augusti
2025; registrets översiktslistor visade inte Apple den 18 november 2026.
Linux skrevs från grunden som en klon av Unix, och Linuxkärnans egen
dokumentation och Torvalds första tillkännagivande säger det; det är
Unixliknande i betydelsen efterlikning, även om en Linuxbaserad produkt
har registrerats som ett UNIX.
Android använder en modifierad Linuxkärna (en studie från 2012 fann att
bara 0,7 procent av kärnans filer var ändrade) och en egen
användarmiljö, så det är byggt på Linux kärna men är inte ett
Linuxskrivbordssystem.
Föreläsningens huvudtext får därför säga att Linux, macOS och
BSD-systemen liknar varandra eftersom BSD-systemen och macOS härstammar
från Unix och Linux efterliknar det, och att Android bygger på
Linuxkärnan.
Den bör inte säga att Linux härstammar från Unix, att BSD-systemen är
certifierade UNIX-system, eller att du kan räkna med samma kommandorad
på Android: ingen källa som vi hittade säger det.
\end{onlyin}

\section{\ensv{Further work}{Fortsatt arbete}} % sync english=a44494 swedish=969e69

\begin{onlyin}{english} % sync english=38ecc6 swedish=816ada
This search left some things undone.
The search should be redone through the university's network, so that
Scopus returns abstracts, and the records should then be screened on
their abstracts, not their titles.
IEEE Xplore, Semantic Scholar and DBLP should be searched when their
limits reset or their keys are replaced; the queries S5, S5b, S4b and the
counter queries C1 to C4 failed on IEEE Xplore.
The studies by Khomh and others and by German and others, and the
chapters by Tarkoma and others and by La Lau, are read in abstract only
and should be read in full; the first two need the publisher's
subscription.
The Open Group's certificate of conformance for macOS 26.0 Tahoe was
not opened, and neither was the explanation of why the register's
overview lists do not show Apple; both should be asked of The Open Group.
The litigation that disputed whether Linux contains Unix code should be
read from a court document or a legal history, not from a news summary.

What the field lacks is as follows.
No study measures how much of today's BSD code is Bell Labs code, or
how much Linux and the BSD systems share; Spinellis's repository makes it
possible, and German and others began on the Linux and BSD part.
There is no study of how novices understand the claim that two systems
\enquote{are Unix} or \enquote{are Unix-like}, which would show whether
the wording in \cref{sec:where-cli} helps or misleads.
If such a study showed that little Bell Labs code remains in today's
BSD systems, the text would have to say that they descend from Unix in
design and history, not in code; if it confirmed the code lineage, the
text would stand as it is.
\end{onlyin}
\begin{onlyin}{swedish}
Sökningen lämnade en del ogjort.
Sökningen bör göras om genom universitetets nätverk, så att Scopus ger
sammanfattningar, och posterna bör då gallras på sina sammanfattningar,
inte sina titlar.
IEEE Xplore, Semantic Scholar och DBLP bör genomsökas när deras gränser
återställs eller deras nycklar byts; frågorna S5, S5b, S4b och
motfrågorna C1 till C4 misslyckades i IEEE Xplore.
Studierna av Khomh med flera och av German med flera, och kapitlen av
Tarkoma med flera och av La Lau, är bara lästa i sammanfattning och bör
läsas i sin helhet; de två första kräver förlagets prenumeration.
The Open Groups intyg om överensstämmelse för macOS 26.0 Tahoe öppnades
inte, och inte heller förklaringen till varför registrets
översiktslistor inte visar Apple; båda bör efterfrågas hos The Open
Group.
Rättstvisten om huruvida Linux innehåller kod från Unix bör läsas i ett
domstolsdokument eller en rättshistorisk framställning, inte i en
nyhetssammanfattning.

Det här saknar forskningsfältet.
Ingen studie mäter hur stor del av dagens BSD-kod som är kod från Bell
Labs, eller hur mycket Linux och BSD-systemen delar; Spinellis arkiv gör
det möjligt, och German med flera började på delen om Linux och BSD.
Det finns ingen studie av hur nybörjare förstår påståendet att två
system \enquote{är Unix} eller \enquote{är Unixliknande}, vilket skulle
visa om formuleringen i \cref{sec:where-cli} hjälper eller vilseleder.
Visade en sådan studie att lite kod från Bell Labs finns kvar i dagens
BSD-system, skulle texten behöva säga att de härstammar från Unix i
utformning och historia, inte i kod; bekräftade den kodlinjen, skulle
texten stå som den är.
\end{onlyin}

\section{\ensv{Hits bearing on the claim}{Träffar som rör påståendet}} % sync english=b678a6 swedish=eaf510

% Author's note (verification and reproduction):
%   scholar session : unix-lineage
%   exports         : litteratursokning/unix-lineage.csv (the session,
%                     every record and its decision; the abstract column
%                     is dropped for size, and every record keeps its DOI
%                     and URL, so its abstract can be looked up) and the
%                     -full hit-list tables
%   searched        : 2026-11-18; providers openalex (title and abstract,
%                     via scripts/openalex_title_abstract.py; default
%                     full-text search superseded), wos, scopus, ieee
%                     (daily quota spent after S4: S4b, S5, S5b, C1-C4
%                     failed); s2: key rejected, HTTP 403; dblp: failed,
%                     anti-bot page; not through the KTH tunnel, so Scopus
%                     returned no abstracts
%   classification  : scholar llm context + llm classify --no-examples,
%                     model anthropic/claude-sonnet-5-5; 109 supports or
%                     qualifies rows read at title (and abstract where one
%                     existed) level, 31 rows (26 works) kept, 78 rows
%                     discarded as no-evidence or app-level-only
%   hit-list tables : scholar sessions export unix-lineage -f table
%                     --bearing-only --lang en|sv --label sok-unix-lineage
%                     --track "Claim D"|"Påstående D" --theme ... -o
%                     litteratursokning/unix-lineage-full(-sv)
%   provenance      : ltnotes.bib (CLAIM/FOUND-VIA/QUOTE/VERIFIED)
\begin{onlyin}{english} % sync english=2a400c swedish=2d84d4
% --- scholar LaTeX fragment --------------------------------------------
% Generated by: scholar sessions export -f table
% Session: unix-lineage
% \input this file from the body of a document whose
% preamble loads:
%   \usepackage{longtable}
%   \usepackage{booktabs}
%   \usepackage{array}
% Text that pdfLaTeX cannot set in T1 (Cyrillic, CJK, ...) is
% wrapped in \scholarnonlatin{...}. The \providecommand lines
% below print the text under XeLaTeX and LuaLaTeX (load a font
% that has the script) and print [non-Latin text] under pdfLaTeX.
% To print it under pdfLaTeX too, define the command before the
% \input, e.g. for Cyrillic with \usepackage[T2A,T1]{fontenc}:
%   \newcommand{\scholarnonlatin}[1]{{\fontencoding{T2A}\selectfont #1}}
% ----------------------------------------------------------------------

\ifdefined\XeTeXversion\providecommand{\scholarnonlatin}[1]{#1}\fi
\ifdefined\directlua\providecommand{\scholarnonlatin}[1]{#1}\fi
\providecommand{\scholarnonlatin}[1]{[non-Latin text]}
\begingroup\footnotesize
\begin{longtable}{@{}%
  >{\raggedright\arraybackslash}p{0.44\textwidth}c%
  >{\raggedright\arraybackslash}p{0.13\textwidth}%
  >{\raggedright\arraybackslash}p{0.26\textwidth}@{}}
\caption{Records bearing on the claim, Claim D (19 cited, 42 supporting and 2 qualifying or contradicting it, out of 1267 unique hits; the 693 adjacent and 505 false hits, 6 excluded although classified as bearing on the claim are not listed). The last column gives the reason for each inclusion; for machine-classified rows the model's confidence is given. A row is cited when the text cites the work, whoever decided on its records. Databases: OpenAlex (OA; not open access), IEEE Xplore (IEEE), Web of Science (WoS), Scopus.}\label{tab:sok-unix-lineage}\\
\toprule
Title & Year & Provider & Reason \\
\midrule
\endfirsthead
\multicolumn{4}{@{}l}{\emph{\tablename~\thetable{} (cont.)}}\\
\toprule
Title & Year & Provider & Reason \\
\midrule
\endhead
\midrule \multicolumn{4}{r@{}}{\emph{continued on next page}}\\
\endfoot
\bottomrule
\endlastfoot
A repository of Unix history and evolution & 2016 & OA, WoS, Scopus & \emph{cited}: BSD lineage from Unix (0.75) \\
A Repository with 44 Years of Unix Evolution & 2015 & OA, WoS, Scopus, IEEE & \emph{cited}: BSD lineage from Unix (0.80) \\
Adapting Linux for Mobile Platforms: An Empirical Study of Android & 2012 & WoS, OA, Scopus & \emph{cited}: Android and the Linux kernel (0.80) \\
Applying bioinformatics in the analysis of software variants & 2012 & IEEE & \emph{cited}: BSD lineage from Unix (0.70) \\
Code siblings: Technical and legal implications of copying code between applications & 2009 & IEEE & \emph{cited}: Code flows between Linux and BSD (0.90) \\
Criteria for UNIX-like systems & 1982 & OA & \emph{cited}: Unix-like defined (0.65) \\
DebianRunner: Running Desktop Applications on Android Smartphones & 2010 & OA & \emph{cited}: Android differs from Linux userland (0.70) \\
Embedded Android: Porting, Extending, and Customizing & 2013 & OA & \emph{cited}: Android differs from Linux userland (0.85) \\
Evolution of the Unix System Architecture: An Exploratory Case Study & 2019 & OA, WoS, Scopus, IEEE & \emph{cited}: BSD lineage from Unix (0.85) \\
Mac OS X and iOS Internals: To the Apple's Core & 2012 & OA & \emph{cited}: macOS and its BSD base (0.70) \\
Measuring similarity of large software systems based on source code correspondence & 2005 & WoS, OA, Scopus & \emph{cited}: BSD lineage from Unix (0.70) \\
Mobile operating systems & 2014 & OA & \emph{cited}: Android and the Linux kernel (0.65) \\
POSIX.1 conformance for Android Applications & 2013 & WoS, OA & \emph{cited}: Android differs from Linux userland (0.50) (app-level-only) \\
Starting with Linux & 2015 & OA & \emph{cited}: Loose talk of descent (0.55) \\
The Comparison: Linux vs. Mac OS X & 2009 & OA & \emph{cited}: Loose talk of descent (0.60) \\
The Open Source Alternative: Understanding Risks and Leveraging Opportunities & 2008 & OA & \emph{cited}: Loose talk of descent (0.45) \\
Trend analysis and Classification of Linux distributions & 2017 & OA & \emph{cited}: Loose talk of descent (0.50) \\
Unix and POSIX in a Few Words & 2021 & OA & \emph{cited}: Unix-like defined (0.60) \\
UNIX License Makes Linux the Last Major Piece of the Puzzle & 2013 & OA & \emph{cited}: Linux origin and status (0.85) \\
35 Years of Linux: The Open Source Engine of Global Infrastructure and Innovation & 2026 & OA & supports the claim (0.50) \\
\scholarnonlatin{図解で明解}Mac OS X\scholarnonlatin{のしくみ} & 2003 & OA & supports the claim (0.50) \\
Advanced synchronization in Mac OS X: Extending Unix to SMP and real-time & 2002 & WoS, OA, Scopus & supports the claim (0.80) \\
An Investigation of Patch Porting Practices of the Linux Kernel Ecosystem & 2024 & OA, WoS, Scopus & supports the claim (0.55) \\
Analysis of Android Kernel & 2009 & OA & supports the claim (0.60) \\
Android & 2010 & OA & supports the claim (0.55) \\
Android at risk: current threats stemming from unprotected local and external resources & 2014 & OA & supports the claim (0.50) \\
Comparison Analysis Android and GNU / Linux & 2016 & OA & supports the claim (0.55) \\
Design of Mini Multi-Process Micro-Kernel Embedded OS on ARM & 2013 & OA, Scopus & supports the claim (0.45) (kernel-internals) \\
Die UNIX-Welt & 2018 & OA & supports the claim (0.50) \\
Eliminating the Security Weakness of Linux and UNIX Operating Systems & 2013 & OA, Scopus & supports the claim (0.60) (no-evidence) \\
Embedding Linux & 2001 & OA, IEEE & supports the claim (0.50) \\
Evolusi Sistem Operasi Android: Sebuah Tinjauan Literatur & 2023 & OA & supports the claim (0.50) \\
Fundamentals of Linux & 2018 & OA & supports the claim (0.50) \\
Hello, Android & 2018 & OA & supports the claim (0.50) \\
Hybridizing Your System & 2009 & OA & supports the claim (0.80) \\
Introducing Darwin and the Mac OS X Command Line & 2011 & OA & supports the claim (0.65) \\
Introducing Darwin and the Shell & 2009 & OA & supports the claim (0.55) \\
Introduction to the Linux Environment & 2020 & OA & supports the claim (0.70) \\
Learning Unix for Mac OS X Panther & 2003 & OA & supports the claim (0.55) \\
Linux and the Web & 2021 & OA & supports the claim (0.80) \\
Linux and Unix Security & 2009 & OA, Scopus & supports the claim (0.75) \\
Mac OS X and iOS & 2011 & OA & supports the claim (0.60) \\
Mac OS X Panther in a Nutshell & 2004 & OA & supports the claim (0.60) \\
MAC OS X Version 10.1 Black Book with Cdrom & 2002 & OA & supports the claim (0.55) \\
Mac OS X: The Missing Manual & 2001 & OA & supports the claim (0.55) \\
Mastering Red Hat Linux 9 & 2004 & OA & supports the claim (0.55) \\
Nature Inspired System Analysis. & 2017 & OA & supports the claim (0.65) \\
Open systems storage platforms & 1992 & OA & supports the claim (0.50) \\
Operating system BSD & 2012 & OA & supports the claim (0.75) \\
Operating Systems & 2026 & OA & supports the claim (0.50) \\
Operating Systems Overview & 2016 & OA & supports the claim (0.55) \\
OS X for Hackers at Heart & 2005 & OA, Scopus & supports the claim (0.55) \\
Programming Mac OS X: A Guide for Unix Developers & 2003 & OA & supports the claim (0.55) \\
Similarity of software system and its measurement tool SMMT & 2007 & OA, Scopus & supports the claim (0.75) \\
SystemStarter and the Mac OS X startup process & 2002 & WoS, OA, Scopus & supports the claim (0.55) (no-evidence) \\
T reble: Fast Software Updates by Creating an Equilibrium in an Active Software Ecosystem of Globally Distributed Stakeholders & 2019 & OA, WoS, Scopus & supports the claim (0.50) \\
The design and implementation of the NetBSD rc.d system & 2001 & WoS, OA & supports the claim (0.50) \\
THE LINUX OPERATING SYSTEM: AN INTRODUCTION & 1995 & OA & supports the claim (0.85) \\
The LiteOS Operating System: Towards Unix-Like Abstractions for Wireless Sensor Networks & 2008 & IEEE & supports the claim (0.50) \\
Unix for Mac OS X & 2002 & OA & supports the claim (0.60) \\
UNIX Forensic Analysis & 2010 & OA, Scopus & supports the claim (0.60) \\
13 SOURCE CODE ISN'T & 2021 & IEEE & qualifies or contradicts the claim (0.55) \\
Absolute freebsd, 2nd edition & 2007 & OA & qualifies or contradicts the claim (0.50) \\
\end{longtable}
\endgroup

\end{onlyin}
\begin{onlyin}{swedish}
% --- scholar LaTeX fragment --------------------------------------------
% Generated by: scholar sessions export -f table
% Session: unix-lineage
% \input this file from the body of a document whose
% preamble loads:
%   \usepackage{longtable}
%   \usepackage{booktabs}
%   \usepackage{array}
% ----------------------------------------------------------------------

\begingroup\footnotesize
\begin{longtable}{@{}%
  >{\raggedright\arraybackslash}p{0.44\textwidth}c%
  >{\raggedright\arraybackslash}p{0.13\textwidth}%
  >{\raggedright\arraybackslash}p{0.26\textwidth}@{}}
\caption{Träffar som rör påståendet, Påstående D (22 citerade, 4 som stöder påståendet och 0 som kvalificerar eller motsäger det, av 1267 unika träffar; de 548 angränsande och 622 felträffarna, 71 uteslutna trots att de klassats som relevanta för påståendet listas inte). Skälet till varje inkludering står i sista kolumnen; för maskinklassade rader anges modellens konfidens. Databaser: OpenAlex (OA; inte open access), IEEE Xplore (IEEE), Web of Science (WoS), Scopus.}\label{tab:sok-unix-lineage}\\
\toprule
Titel & År & Leverantör & Skäl \\
\midrule
\endfirsthead
\multicolumn{4}{@{}l}{\emph{\tablename~\thetable{} (forts.)}}\\
\toprule
Titel & År & Leverantör & Skäl \\
\midrule
\endhead
\midrule \multicolumn{4}{r@{}}{\emph{forts.\ på nästa sida}}\\
\endfoot
\bottomrule
\endlastfoot
A repository of Unix history and evolution & 2016 & OA & \emph{citerad}: BSD härstammar från Unix \\
A Repository with 44 Years of Unix Evolution & 2015 & OA & \emph{citerad}: BSD härstammar från Unix \\
Adapting Linux for Mobile Platforms: An Empirical Study of Android & 2012 & WoS, OA & \emph{citerad}: Android och Linuxkärnan \\
An Evolution of Android Operating System and Its Version & 2015 & OA & \emph{citerad}: Android och Linuxkärnan \\
Analysis of Android Kernel & 2009 & OA & \emph{citerad}: Android och Linuxkärnan \\
Android OS CASE STUDY & 2021 & OA & \emph{citerad}: Android och Linuxkärnan \\
Applying bioinformatics in the analysis of software variants & 2012 & IEEE & \emph{citerad}: BSD härstammar från Unix \\
Code siblings: Technical and legal implications of copying code between applications & 2009 & IEEE & \emph{citerad}: Kod flyttas mellan Linux och BSD \\
Comparison Analysis Android and GNU / Linux & 2016 & OA & \emph{citerad}: Android och Linuxkärnan \\
Criteria for UNIX-like systems & 1982 & OA & \emph{citerad}: Unixliknande definierat \\
DebianRunner: Running Desktop Applications on Android Smartphones & 2010 & OA & \emph{citerad}: Androids användarmiljö skiljer sig från Linux \\
Embedded Android: Porting, Extending, and Customizing & 2013 & OA & \emph{citerad}: Androids användarmiljö skiljer sig från Linux \\
Evolution of the Unix System Architecture: An Exploratory Case Study & 2019 & OA & \emph{citerad}: BSD härstammar från Unix \\
Mac OS X and iOS Internals: To the Apple's Core & 2012 & OA & \emph{citerad}: macOS och dess BSD-grund \\
Measuring similarity of large software systems based on source code correspondence & 2005 & WoS, OA & \emph{citerad}: BSD härstammar från Unix \\
Mobile operating systems & 2014 & OA & \emph{citerad}: Android och Linuxkärnan \\
POSIX.1 conformance for Android Applications & 2013 & WoS, OA & \emph{citerad}: Androids användarmiljö skiljer sig från Linux \\
The Comparison: Linux vs. Mac OS X & 2009 & OA & \emph{citerad}: Lös tal om härstamning \\
The Open Source Alternative: Understanding Risks and Leveraging Opportunities & 2008 & OA & \emph{citerad}: Lös tal om härstamning \\
Trend analysis and Classification of Linux distributions & 2017 & OA & \emph{citerad}: Linux ursprung och ställning \\
Unix and POSIX in a Few Words & 2021 & OA & \emph{citerad}: Unixliknande definierat \\
UNIX License Makes Linux the Last Major Piece of the Puzzle & 2013 & OA & \emph{citerad}: Linux ursprung och ställning \\
Half-Century of Unix: History, Preservation, and Lessons Learned & 2017 & OA & stöder påståendet (0.75) \\
The History of Unix in the History of Software & 2017 & OA & stöder påståendet (0.45) \\
Unix History Repository & 2018 & OA & stöder påståendet (0.80) \\
UNIX USG V. 2 AND BSD 4. S: COMPARATIVE STUDY. & 1985 & Scopus & stöder påståendet (0.50) \\
\end{longtable}
\endgroup

\end{onlyin}
